\documentclass[12pt,a4paper]{article}

% --- Pacotes de Suporte ---
\usepackage{cmap} % corrige o mapeamento Unicode de ligaduras (fi/ff/fl)
\usepackage[utf8]{inputenc}
\usepackage[T1]{fontenc}
\usepackage[expansion=false]{microtype}
\usepackage[brazilian]{babel}
\usepackage{amsmath}
\usepackage{amsfonts}
\usepackage{amssymb}
\usepackage{geometry}
\usepackage{bm}
\usepackage{graphicx}
\usepackage{float}
\usepackage{booktabs}
\usepackage{hyperref}
\usepackage{fancyhdr}
\usepackage{titlesec}
\usepackage{tabularx}
\usepackage{longtable}

\hyphenation{
  off-shore Off-shore
  con-tí-nuo con-tí-nu-a
  te-r-mo-di-nâ-mi-ca te-r-mo-di-nâ-mi-cas
  me-câ-ni-ca me-câ-ni-cas
  ro-to-di-nâ-mi-ca ro-to-di-nâ-mi-cas
  ins-tru-men-ta-ção ins-tru-men-ta-ções
}

\numberwithin{equation}{section}
\numberwithin{table}{section}
\numberwithin{figure}{section}

\geometry{margin=2.5cm, headheight=15pt}
\addtolength{\topmargin}{-0.5pt}
\setlength{\emergencystretch}{1em}

\hypersetup{
	colorlinks=true,
	linkcolor=black,
	citecolor=black,
	urlcolor=black,
	pdftitle={Fundamentos de Mecânica do Contínuo, Termofluidos e Sistemas de Medição},
	pdfauthor={Marcelo Moreira Mejias},
	pdfsubject={Bases Teóricas para Máquinas Rotodinâmicas},
	pdfkeywords={mecânica do contínuo, termodinâmica, mecânica dos fluidos, transferência de calor, máquinas rotodinâmicas}
}

\pagestyle{fancy}
\fancyhf{}
\renewcommand{\sectionmark}[1]{\markboth{#1}{}}
\fancyhead[L]{\nouppercase{\leftmark}}
\fancyhead[R]{\thepage}
\renewcommand{\headrulewidth}{0.4pt}

% --- Identificação de volume/revisão ---
\newcommand{\colecaotitulo}{Enciclopédia de Máquinas Rotodinâmicas}
\newcommand{\volumeatual}{Volume I}
\newcommand{\revdata}{Rev. 23/08/2026}

% --- Identidade visual dos capítulos ---
\newcommand{\chapterword}{Capítulo}
\titleformat{\section}
	{\normalfont\Large\bfseries\raggedright}
	{\chapterword~\thesection}{1em}{}
	[\vspace{-0.5em}\rule{\textwidth}{0.4pt}]


\begin{document}


	% --- CAPA ---
	\begin{titlepage}
		\centering
		\vspace*{1.5cm}
		{\large \textsc{\colecaotitulo\ --- \volumeatual}} \\[1cm]
		\rule{\textwidth}{0.4pt} \\[0.8cm]
		{\Huge \textbf{Fundamentos de Mecânica do Contínuo,}} \\[0.3cm]
		{\Huge \textbf{Termofluidos e Sistemas de Medição}} \\[0.5cm]
		{\Large \textbf{Bases Teóricas para Máquinas Rotodinâmicas}} \\[0.8cm]
		\rule{\textwidth}{0.4pt} \\[1.5cm]

		% Placeholder para ilustração conceitual do Volume I
		\vspace{5cm}

		\vfill
		{\large \textbf{Autor: Engenheiro Marcelo Moreira Mejias, M.Sc.}} \\[0.5cm]

		{\large \today} \\[0.3cm]
		{\normalsize \revdata}
	\end{titlepage}

	% --- DEDICATÓRIA ---
	\newpage
	\thispagestyle{empty}
	\vspace*{\fill}
	\begin{flushright}
		\textit{Aos que constroem máquinas a partir dos princípios,}
		\textit{antes de confiar nas conveniências.}
	\end{flushright}
	\vspace*{\fill}


% --- PREFÁCIO ---
\newpage
\section*{Prefácio}
\addcontentsline{toc}{section}{Prefácio}

Este volume reúne os fundamentos teóricos que sustentam o projeto de máquinas
rotodinâmicas em geral e o estudo detalhado de compressores de parafuso
apresentado no Volume II desta coleção. A intenção não é substituir livros
clássicos de mecânica do contínuo, termodinâmica, mecânica dos fluidos,
transferência de calor, eletricidade, mecânica e instrumentação, mas oferecer
uma versão condensada, coerente e orientada ao engenheiro que precisa aplicar
esses conceitos em equipamentos reais.

A escolha dos tópicos reflete as demandas práticas da modelagem de máquinas
rotativas. Mecânica do contínuo fornece a linguagem de tensores e equações de
conservação que regem sólidos e fluidos. A termodinâmica define limites de
eiciência e irreversibilidade. A mecânica dos fluidos e a transferência de
calor traduzem geometria em perda de carga, troca térmica e ruído. A eletricidade
 e a mecânica conectam o motor ao eixo e o eixo aos mancais. A instrumentação
 permite fechar o laço entre modelo e realidade.

Cada capítulo foi organizado para ser lido de forma independente, mas também
como parte de uma sequência que culmina na modelagem integrada de sistemas
rotodinâmicos. As equações são numeradas e referenciadas de modo consistente,
 facilitando a ligação com o Volume II.

\vspace{1em}
\hfill \textit{O Autor}

	% --- SUMÁRIO ---
	\newpage
	\tableofcontents
	\newpage

	% --- LISTAS DE FIGURAS E TABELAS ---
	\phantomsection
	\addcontentsline{toc}{section}{Lista de Figuras}
	\listoffigures
	\newpage
	\phantomsection
	\addcontentsline{toc}{section}{Lista de Tabelas}
	\listoftables
	\newpage

% --- NOMENCLATURA ---
\newpage
\section*{Nomenclatura}
\addcontentsline{toc}{section}{Nomenclatura}

A tabela abaixo reúne os símbolos de uso transversal neste volume. Símbolos
específicos de cada capítulo são definidos localmente. A notação vetorial e
tensorial segue a convenção usual da mecânica do contínuo: vetores em
itálico com seta ($\vec{v}$) ou negrito ($\mathbf{v}$), tensores de segunda
ordem em negrito ($\boldsymbol{\sigma}$), e componentes indicadas por
índices ($v_i$, $\sigma_{ij}$).

{\small
\setlength{\tabcolsep}{3pt}
\renewcommand{\arraystretch}{1.15}

\subsection*{Grandezas Mecânicas e do Contínuo}
\begin{longtable}{>{\raggedright\arraybackslash}p{2.8cm} >{\raggedright\arraybackslash}p{6.5cm} >{\raggedright\arraybackslash}p{2.5cm} >{\raggedright\arraybackslash}p{2.8cm}}
	\toprule
	\textbf{Símbolo} & \textbf{Significado} & \textbf{Unidade} & \textbf{Capítulos} \\ \midrule
	\endhead
	$\mathcal{B}$ & Grandeza extensiva genérica em balanço & conforme a grandeza & 1, 2 \\
	$\beta$ & Grandeza intensiva associada a $\mathcal{B}$ (por unidade de massa) & conforme a grandeza & 1, 2 \\
	$\mathbf{j}_\beta$ & Fluxo não convectivo de $\mathcal{B}$ & conforme a grandeza & 2 \\
	$\dot{\mathcal{G}}_\beta$ & Taxa de geração volumétrica de $\mathcal{B}$ & conforme a grandeza & 2 \\
	$\mathrm{Kn}$ & Número de Knudsen & adimensional & 2 \\
	$\mathbf{F}$ & Gradiente de deformação (não confundir com força $\mathbf{F}$, Caps. 6 e 8) & adimensional & 2 \\
	$J$ & Jacobiano da deformação, $\det\mathbf{F}$ & adimensional & 2 \\
	$\mathbf{C}$ & Tensor de Cauchy-Green à direita & adimensional & 2 \\
	$\mathbf{E}_L$ & Tensor de deformação de Green-Lagrange (não confundir com $E$, módulo de Young) & adimensional & 2 \\
	$\mathbf{S}$ & Tensor taxa de deformação & s$^{-1}$ & 2, 4 \\
	$\mathbf{W}$ & Tensor taxa de rotação, ou spin (não confundir com $W$, razão de umidade) & s$^{-1}$ & 2 \\
	$\mathbf{t}_{(\mathbf{n})}$ & Vetor tensão sobre superfície de normal $\mathbf{n}$ & Pa & 2 \\
	$\lambda_L$ & Primeira constante de Lamé & Pa & 2 \\
	$\lambda_v$ & Segunda viscosidade (viscosidade volumétrica) & Pa\,s & 2 \\
	$\sigma_{VM}$ & Tensão equivalente de von Mises & Pa & 2, 8 \\
	$\mathbf{u}$ & Vetor deslocamento & m & 2, 8 \\
	$\boldsymbol{\varepsilon}$ & Tensor deformação infinitesimal & adimensional & 2 \\
	$\boldsymbol{\sigma}$ & Tensor tensão de Cauchy & Pa & 2 \\
	$\rho$ & Massa específica & kg\,m$^{-3}$ & 2--6 \\
	$\mathbf{v}$ & Vetor velocidade do material & m\,s$^{-1}$ & 2--4 \\
	$\mathbf{a}$ & Vetor aceleração & m\,s$^{-2}$ & 2, 8 \\
	$\boldsymbol{\omega}$ & Vetor velocidade angular & rad\,s$^{-1}$ & 2, 8 \\
	$E$ & Módulo de elasticidade (Young) & Pa & 2, 8 \\
	$\nu$ & Coeficiente de Poisson & adimensional & 2 \\
	$G$ & Módulo de cisalhamento (= segunda constante de Lamé) & Pa & 2 \\
	$\mu$ & Viscosidade dinâmica & Pa\,s & 2, 4 \\
	$\nu_k$ & Viscosidade cinemática & m$^2$\,s$^{-1}$ & 4 \\
	$\mathbf{g}$ & Aceleração da gravidade & m\,s$^{-2}$ & 2, 4, 8 \\
	$J_p$ & Momento polar de inércia (não confundir com $J$, jacobiano, Cap. 2, nem $\mathbf{J}$, densidade de corrente, Cap. 7) & kg\,m$^2$ & 8 \\
	$\tau$ & Torque (não confundir com $T$, temperatura, Caps. 3--5) & N\,m & 7, 8 \\ \bottomrule
\end{longtable}

\subsection*{Grandezas Termodinâmicas e de Transferência}
\begin{longtable}{>{\raggedright\arraybackslash}p{2.8cm} >{\raggedright\arraybackslash}p{6.5cm} >{\raggedright\arraybackslash}p{2.5cm} >{\raggedright\arraybackslash}p{2.8cm}}
	\toprule
	\textbf{Símbolo} & \textbf{Significado} & \textbf{Unidade} & \textbf{Capítulos} \\ \midrule
	\endhead
	$T$ & Temperatura termodinâmica & K & 3--5 \\
	$P$ & Pressão absoluta & Pa, bar & 2--6 \\
	$v$ (minúscula) & Volume específico & m$^3$\,kg$^{-1}$ & 3 \\
	$u$ (minúscula) & Energia interna específica & J\,kg$^{-1}$ & 3 \\
	$h$ (minúscula) & Entalpia específica & J\,kg$^{-1}$ & 3--5 \\
	$s$ (minúscula) & Entropia específica & J\,kg$^{-1}$\,K$^{-1}$ & 3 \\
	$q$ & Calor por unidade de massa & J\,kg$^{-1}$ & 3 \\
	$w$ & Trabalho por unidade de massa & J\,kg$^{-1}$ & 3 \\
	$c_p$ & Calor específico a pressão constante & J\,kg$^{-1}$\,K$^{-1}$ & 3--5 \\
	$c_v$ & Calor específico a volume constante & J\,kg$^{-1}$\,K$^{-1}$ & 3 \\
	$\gamma$ & Razão de calores específicos ($c_p/c_v$) & adimensional & 3, 4 \\
	$\Pi$ & Razão de pressão, $p_2/p_1$ & adimensional & 3 \\
	$n$ & Índice politrópico (não confundir com a normal $\mathbf{n}$, Cap. 2) & adimensional & 3 \\
	$\eta_{is}$ & Rendimento isentrópico & adimensional & 3 \\
	$\eta_p$ & Rendimento politrópico & adimensional & 3 \\
	$\eta_{ex}$ & Rendimento exergético & adimensional & 3 \\
	$Z_c$ & Fator de compressibilidade (não confundir com $Z$, impedância, Cap. 7) & adimensional & 3 \\
	$\psi$ & Exergia específica de fluxo & J\,kg$^{-1}$ & 3 \\
	$\dot{S}_{gen}$ & Taxa de geração de entropia & W\,K$^{-1}$ & 3 \\
	$\dot{E}_d$ & Taxa de destruição de exergia & W & 3 \\
	$T_0$, $p_0$ & Temperatura e pressão do ambiente de referência & K, Pa & 3 \\
	$R$ & Constante específica do gás & J\,kg$^{-1}$\,K$^{-1}$ & 3, 4 \\
	$\dot{Q}$ & Taxa de transferência de calor & W & 3--6 \\
	$\dot{W}$ & Potência (trabalho por unidade de tempo) & W & 2, 3, 8 \\
	$k$ & Condutividade térmica & W\,m$^{-1}$\,K$^{-1}$ & 5 \\
	$h$ (latino) & Coeficiente de convecção de calor & W\,m$^{-2}$\,K$^{-1}$ & 5 \\
	$\varepsilon$ (griego) & Efetividade de trocador de calor & adimensional & 5 \\ \bottomrule
\end{longtable}

\subsection*{Grandezas Elétricas e de Instrumentação}
\begin{longtable}{>{\raggedright\arraybackslash}p{2.8cm} >{\raggedright\arraybackslash}p{6.5cm} >{\raggedright\arraybackslash}p{2.5cm} >{\raggedright\arraybackslash}p{2.8cm}}
	\toprule
	\textbf{Símbolo} & \textbf{Significado} & \textbf{Unidade} & \textbf{Capítulos} \\ \midrule
	\endhead
	$V$ & Tensão elétrica & V & 7 \\
	$I$ & Corrente elétrica (não confundir com o momento de inércia, $J_p$, Cap. 8) & A & 7 \\
	$R$ (latino) & Resistência elétrica & $\Omega$ & 7 \\
	$Z$ (latino) & Impedância elétrica & $\Omega$ & 7 \\
	$P_{el}$ & Potência elétrica ativa & W & 7 \\
	$Q_{el}$ & Potência elétrica reativa & var & 7 \\
	$S_{el}$ & Potência elétrica aparente & VA & 7 \\
	$\cos\varphi$ & Fator de potência (de deslocamento) & adimensional & 7 \\
	$\mathrm{THD}_i$ & Distorção harmônica total de corrente & adimensional & 7 \\
	$f$ & Frequência elétrica/rotativa & Hz & 7, 8 \\
	$\omega = 2\pi f$ & Frequência angular & rad\,s$^{-1}$ & 7, 8 \\
	$\rho_q$ & Densidade volumétrica de carga & C\,m$^{-3}$ & 7 \\
	$\rho_e$ & Resistividade elétrica (não confundir com $\rho_q$, densidade de carga, nem com $\rho$, massa específica) & $\Omega$\,m & 7 \\
	$\mathbf{J}$ & Densidade de corrente (não confundir com $J$, jacobiano, Cap. 2, nem $J_p$, Cap. 8) & A\,m$^{-2}$ & 7 \\
	$\mathbf{E}$ & Campo elétrico (não confundir com $E$, módulo de Young, nem $\mathbf{E}_L$, Cap. 2) & V\,m$^{-1}$ & 7 \\
	$\mathbf{D}$ & Deslocamento elétrico & C\,m$^{-2}$ & 7 \\
	$\mathbf{B}$, $\hat{B}$ & Indução magnética e sua amplitude & T & 7 \\
	$\mathbf{H}$ & Campo magnético & A\,m$^{-1}$ & 7 \\
	$\sigma_e$ & Condutividade elétrica (não confundir com $\boldsymbol{\sigma}$, tensão, Cap. 2) & S\,m$^{-1}$ & 7 \\
	$\varepsilon_e$ & Permissividade elétrica (não confundir com a deformação, Cap. 2, nem com a efetividade, Cap. 5) & F\,m$^{-1}$ & 7 \\
	$\mu_m$, $\mu_0$ & Permeabilidade magnética do meio e do vácuo (não confundir com $\mu$, viscosidade) & H\,m$^{-1}$ & 7 \\
	$\lambda$, $\hat{\Phi}$ & Fluxo concatenado e fluxo por polo (não confundir com $\lambda_L$, $\lambda_v$, Cap. 2) & Wb & 7 \\
	$\lambda_{em}$ & Comprimento de onda eletromagnético & m & 7 \\
	$\delta_p$ & Profundidade pelicular & m & 7 \\
	$p$ & Número de polos do enrolamento (não confundir com $p$, pressão, Cap. 3) & adimensional & 7 \\
	$n_s$, $s$ & Velocidade síncrona e escorregamento (este $s$ não é a entropia específica, Cap. 3) & rpm, adimensional & 7 \\
	$\tau_{em}$ & Torque eletromagnético & N\,m & 7, 8 \\
	$P_{ag}$ & Potência de entreferro & W & 7 \\
	$P_{fe}$ & Perda no núcleo (ferro) & W & 7 \\
	$\eta_{mot}$ & Rendimento do motor & adimensional & 7 \\
	$v_{cm}$ & Tensão de modo comum & V & 7 \\
	$U$ & Incerteza expandida da medição & unidade da grandeza & 9 \\
	$k$ (cobertura) & Fator de cobertura da incerteza & adimensional & 9 \\
	$\sigma_x$ & Desvio padrão experimental & unidade da grandeza & 9 \\ \bottomrule
\end{longtable}

\subsection*{Grandezas de Psicrometria e Transferência de Massa}
\begin{longtable}{>{\raggedright\arraybackslash}p{2.8cm} >{\raggedright\arraybackslash}p{6.5cm} >{\raggedright\arraybackslash}p{2.5cm} >{\raggedright\arraybackslash}p{2.8cm}}
	\toprule
	\textbf{Símbolo} & \textbf{Significado} & \textbf{Unidade} & \textbf{Capítulos} \\ \midrule
	\endhead
	$W$ & Razão de umidade (massa de vapor/massa de ar seco) & kg\,kg$^{-1}$ & 11 \\
	$\phi$ & Umidade relativa & adimensional & 11 \\
	$P_{\text{sat}}$ & Pressão de saturação do vapor d'água & Pa & 11 \\
	$T_{dp}$ & Ponto de orvalho & K, $^{\circ}$C & 11 \\
	$h_{fg}$ & Entalpia de vaporização do água & J\,kg$^{-1}$ & 11, 17 \\
	$c_{p,\text{úmido}}$ & Calor específico do ar úmido & J\,kg$^{-1}$\,K$^{-1}$ & 11 \\
	$D_{AB}$ & Coeficiente de difusão mássica & m$^2$\,s$^{-1}$ & 11 \\
	$q_{eq}$ & Capacidade de equilíbrio de adsorção & kg\,kg$^{-1}$ & 11 \\ \bottomrule
\end{longtable}

\subsection*{Grandezas de Escoamento Compressível}
\begin{longtable}{>{\raggedright\arraybackslash}p{2.8cm} >{\raggedright\arraybackslash}p{6.5cm} >{\raggedright\arraybackslash}p{2.5cm} >{\raggedright\arraybackslash}p{2.8cm}}
	\toprule
	\textbf{Símbolo} & \textbf{Significado} & \textbf{Unidade} & \textbf{Capítulos} \\ \midrule
	\endhead
	$Ma$ & Número de Mach & adimensional & 12 \\
	$a$ & Velocidade do som no meio & m\,s$^{-1}$ & 12, 15 \\
	$\Phi$ & Função de escoamento compressível & adimensional & 12 \\
	$r$ & Razão de pressão crítica/normalizada & adimensional & 12 \\
	$K$ & Coeficiente de perda localizada & adimensional & 12 \\
	$k_p$ & Permeabilidade do meio poroso & m$^2$ & 12 \\ \bottomrule
\end{longtable}

\subsection*{Grandezas de Lubrificação e Tribologia}
\begin{longtable}{>{\raggedright\arraybackslash}p{2.8cm} >{\raggedright\arraybackslash}p{6.5cm} >{\raggedright\arraybackslash}p{2.5cm} >{\raggedright\arraybackslash}p{2.8cm}}
	\toprule
	\textbf{Símbolo} & \textbf{Significado} & \textbf{Unidade} & \textbf{Capítulos} \\ \midrule
	\endhead
	$\mu$ & Viscosidade dinâmica (também em Mecânica dos Fluidos) & Pa\,s & 4, 13 \\
	$\Lambda$ & Número de Sommerfeld & adimensional & 13 \\
	$h_{\text{filme}}$ & Espessura de filme lubrificante & m & 13 \\
	$p_{\text{Hertz}}$ & Pressão de contato de Hertz & Pa & 13 \\
	$\tau$ & Tensão cisalhante no fluido & Pa & 13 \\ \bottomrule
\end{longtable}

\subsection*{Grandezas de Engrenagens e Rolamentos}
\begin{longtable}{>{\raggedright\arraybackslash}p{2.8cm} >{\raggedright\arraybackslash}p{6.5cm} >{\raggedright\arraybackslash}p{2.5cm} >{\raggedright\arraybackslash}p{2.8cm}}
	\toprule
	\textbf{Símbolo} & \textbf{Significado} & \textbf{Unidade} & \textbf{Capítulos} \\ \midrule
	\endhead
	$m$ & Módulo da engrenagem & m & 14 \\
	$z$ & Número de dentes & adimensional & 14 \\
	$Y$ & Fator de forma de Lewis & adimensional & 14 \\
	$P_r$ & Carga equivalente do rolamento & N & 14 \\
	$C$ & Carga dinâmica nominal do rolamento & N & 14 \\
	$L_{10h}$ & Vida nominal do rolamento em horas & h & 14 \\
	$p_{\text{rol}}$ & Expoente de vida do rolamento & adimensional & 14 \\ \bottomrule
\end{longtable}

\subsection*{Grandezas de Rotordinâmica e Acústica}
\begin{longtable}{>{\raggedright\arraybackslash}p{2.8cm} >{\raggedright\arraybackslash}p{6.5cm} >{\raggedright\arraybackslash}p{2.5cm} >{\raggedright\arraybackslash}p{2.8cm}}
	\toprule
	\textbf{Símbolo} & \textbf{Significado} & \textbf{Unidade} & \textbf{Capítulos} \\ \midrule
	\endhead
	$k_{\text{mancal}}$ & Rigidez do mancal & N\,m$^{-1}$ & 15 \\
	$\omega_c$ & Velocidade crítica angular & rad\,s$^{-1}$ & 15 \\
	$GMF$ & Frequência de engrenamento & Hz & 15 \\
	$f_{bp}$ & Frequência de passagem de lóbulos & Hz & 15 \\
	$f_H$ & Frequência de ressonância de Helmholtz & Hz & 15 \\ \bottomrule
\end{longtable}

} % fim small

\newpage

%=============================================================================
% CAPÍTULO 1 — INTRODUÇÃO
%=============================================================================
\section{Introdução ao Volume I}

\subsection{Por que um volume de fundamentos}

A engenharia de máquinas rotodinâmicas avançou a ponto de exigir modelos
quantitativos de todos os fenômenos que governam um equipamento: escoamento
interno e externo, troca térmica, deformação mecânica, conversão
eletromecânica de energia, vibração e medição. Cada um desses fenômenos tem
sua própria teoria clássica, com notação, convenções e hipóteses próprias.
Para quem projeta ou mantém uma máquina real, o desafio não é dominar cada
teoria isoladamente, mas saber quando e como combiná-las.

A dificuldade prática, portanto, raramente está no interior de uma disciplina:
está nas suas fronteiras. Vale um exemplo concreto, do tipo que ocupa um
engenheiro de campo. A temperatura de descarga do primeiro estágio de um
compressor sobe dois graus ao longo de algumas semanas, sem que nenhuma outra
leitura saia da faixa normal. Cada disciplina deste volume oferece uma
explicação legítima para o mesmo sintoma:

\begin{itemize}
	\item \textbf{Termodinâmica:} o índice politrópico efetivo aumentou, e o
	processo se afastou da compressão de referência
	(Capítulo~\ref{sec:termodinamica}).
	\item \textbf{Mecânica dos fluidos:} a fuga interna pelas folgas cresceu, e
	parte do gás já comprimido retorna para ser recomprimido
	(Capítulo~\ref{sec:fluidos}).
	\item \textbf{Transferência de calor:} o trocador perdeu efetividade e
	elevou a temperatura de admissão a jusante, que realimenta o estágio pelo
	circuito de arrefecimento (Capítulo~\ref{sec:calor}).
	\item \textbf{Mecânica do contínuo:} a dilatação diferencial entre rotor e
	carcaça alterou a folga efetiva de operação a quente
	(Capítulo~\ref{sec:continuo}).
	\item \textbf{Instrumentação:} o próprio sensor derivou, e a máquina não
	mudou coisa alguma (Capítulo~\ref{sec:instrumentacao}).
\end{itemize}

As cinco hipóteses são compatíveis com o dado observado. Nenhuma teoria
isolada as distingue, porque cada uma explica o sintoma dentro do seu próprio
domínio e não tem como falar sobre os demais. Separá-las exige escrever as
cinco no mesmo conjunto de variáveis e confrontá-las simultaneamente contra
todas as leituras disponíveis --- que é, precisamente, o que o Volume II desta
coleção chama de Gêmeo Digital. Este volume prepara essa escrita comum.

Os capítulos foram organizados de modo que um mesmo conceito --- por exemplo,
conservação de massa ou balanço de energia --- apareça primeiro em sua forma
geral (Mecânica do Contínuo, Capítulo~\ref{sec:continuo}) e depois em suas
formas aplicadas (Termodinâmica, Capítulo~\ref{sec:termodinamica}; Mecânica
dos Fluidos, Capítulo~\ref{sec:fluidos}; Transferência de Calor,
Capítulo~\ref{sec:calor}). A repetição é deliberada: é ela que torna visível
quais hipóteses foram acrescentadas em cada especialização e, portanto, onde
cada modelo deixa de valer.

\subsection{O esqueleto comum: balanços sobre uma fronteira}

Há uma única sentença por trás de quase todo o conteúdo deste volume: aquilo
que se acumula dentro de uma região do espaço é o que entrou, menos o que
saiu, mais o que foi gerado ali dentro. A Equação~\eqref{eq:v1-c1-1} escreve
essa sentença na forma mais crua possível, para uma grandeza extensiva
genérica $\mathcal{B}$ --- massa, quantidade de movimento, energia, entropia
ou carga elétrica:

\begin{equation}\label{eq:v1-c1-1}
\frac{d\mathcal{B}}{dt} = \dot{\mathcal{B}}_{\text{ent}} -
\dot{\mathcal{B}}_{\text{sai}} + \dot{\mathcal{B}}_{\text{ger}},
\end{equation}
onde $d\mathcal{B}/dt$ é a taxa de acúmulo dentro da região delimitada,
$\dot{\mathcal{B}}_{\text{ent}}$ e $\dot{\mathcal{B}}_{\text{sai}}$ são os
fluxos que atravessam a fronteira, e $\dot{\mathcal{B}}_{\text{ger}}$ é a taxa
de geração interna.

O que distingue um capítulo do outro é o que se substitui em cada termo:

\begin{itemize}
	\item \textbf{Massa:} não há geração ($\dot{\mathcal{B}}_{\text{ger}}=0$) e
	a equação se torna a continuidade.
	\item \textbf{Quantidade de movimento:} a ``geração'' é a resultante das
	forças externas aplicadas, e a equação se torna a segunda lei de Newton
	escrita para uma região do espaço em vez de para um corpo.
	\item \textbf{Energia:} não há geração --- apenas conversão entre formas
	--- e os fluxos de fronteira incluem calor e trabalho; a equação se torna a
	primeira lei da termodinâmica.
	\item \textbf{Entropia:} a geração é obrigatoriamente não negativa, e essa
	única restrição de sinal é a segunda lei.
\end{itemize}

Vale registrar o que a Equação~\eqref{eq:v1-c1-1} \textbf{não} é. Ela não
constitui uma terceira formulação, independente das que virão; é um esquema
mnemônico. A forma rigorosa aparece adiante em dois níveis distintos ---
local e diferencial no Capítulo~\ref{sec:continuo}, integral sobre volumes de
controle finitos no Capítulo~\ref{sec:sistemas} --- e são essas as formas a
usar em qualquer cálculo. A utilidade do esquema é outra: sempre que uma
dedução deste volume parecer desconectada das demais, o exercício produtivo é
perguntar qual $\mathcal{B}$ ela está balanceando e sobre qual fronteira.

\subsection{Escopo e limites deste volume}

Este volume é uma condensação orientada à aplicação, não um substituto dos
tratados clássicos de cada disciplina. A diferença de propósito determina o
que entra e o que fica de fora.

Entram os resultados que um engenheiro precisa ter à mão para montar o modelo
de uma máquina: as equações de conservação e as relações constitutivas
(Capítulo~\ref{sec:continuo}), os limites impostos pela segunda lei e a
contabilidade de irreversibilidades (Capítulo~\ref{sec:termodinamica}), as
correlações que traduzem geometria em perda de carga e em troca térmica
(Capítulos~\ref{sec:fluidos} e~\ref{sec:calor}), o balanço sobre volumes de
controle que amarra subsistemas em série (Capítulo~\ref{sec:sistemas}), a
conversão eletromecânica e o acionamento (Capítulo~\ref{sec:eletrica}), a
dinâmica de rotores e a fadiga (Capítulo~\ref{sec:mecanica}) e a cadeia de
medição com sua incerteza (Capítulo~\ref{sec:instrumentacao}).

Ficam deliberadamente de fora quatro classes de conteúdo:

\begin{itemize}
	\item \textbf{Métodos numéricos e sua implementação} --- discretização,
	malhas, esquemas de solução de sistemas não lineares. Este volume apresenta
	as equações; resolvê-las numericamente é assunto do Volume II e do
	simulador que o acompanha.
	\item \textbf{Formalismo matemático avançado.} Demonstrações de
	existência e unicidade, análise funcional e tratamento rigoroso de soluções
	fracas ficam de fora. As deduções \emph{físicas} são desenvolvidas --- toda
	equação deste volume vem acompanhada da sua base, das suas premissas e do
	caminho que leva até ela ---, mas o rigor buscado é o do engenheiro, não o
	do analista.
	\item \textbf{Projeto normativo.} Dimensionamento por código --- vasos de
	pressão, engrenagens, dispositivos de alívio --- é matéria de norma e de
	projeto, não de modelagem da operação.
	\item \textbf{Tabelas de propriedades.} Valores numéricos de materiais e
	fluidos mudam por fornecedor, lote e condição; reproduzi-los aqui daria a
	impressão de uma precisão que o volume não tem como garantir.
\end{itemize}

Uma consequência prática do último item merece destaque, porque atravessa toda
a coleção. Várias grandezas que aparecem como constantes nas equações --- o
coeficiente de perda de um filtro, o rendimento mecânico de uma transmissão, o
coeficiente que quantifica o vazamento por uma folga --- não são constantes
físicas, e sim valores de catálogo levantados em bancada. Elas variam com o
lote de fabricação, com a montagem e com a instalação. O Volume II trata a
determinação dessas grandezas a partir de dados de operação como um problema
inverso; este volume se limita a assinalar, em cada caso, que o número não é
universal.

\subsection{Convenções de notação, unidades e numeração}

As convenções abaixo valem para todo o volume e são as mesmas do Volume II,
para que o leitor transite entre os dois sem retraduzir símbolos.

\textbf{Unidades.} Todas as grandezas estão no Sistema Internacional. Pressões
são \textbf{absolutas}, salvo indicação explícita em contrário; temperaturas
aparecem em kelvin nas equações e em graus Celsius apenas quando o texto trata
de leitura de instrumento. A distinção não é cosmética: uma razão entre
temperaturas só tem significado físico na escala absoluta, e avaliá-la em
Celsius pode inverter o sinal do resultado sem produzir nenhum sintoma
aparente de erro.

\textbf{Notação vetorial e tensorial.} Vetores em negrito ($\mathbf{v}$),
tensores de segunda ordem em negrito grego ($\boldsymbol{\sigma}$),
componentes indicadas por índices ($v_i$, $\sigma_{ij}$), com soma implícita
sobre índices repetidos.

\textbf{Numeração.} Equações, tabelas e figuras são numeradas por capítulo. O
rótulo interno de cada equação segue o padrão
\texttt{eq:v1-c\textit{C}-\textit{N}}, em que \textit{C} é o capítulo e
\textit{N} a posição dentro dele; o prefixo \texttt{v1} distingue este volume
do Volume II, cujos rótulos usam \texttt{eq:s\textit{C}-\textit{N}}. Toda
equação, tabela e figura é citada pelo número no texto que a antecede, nunca
como ``a equação abaixo'' --- de modo que o leitor saiba o que vai encontrar
antes de encontrá-lo.

\textbf{Símbolos.} O alfabeto grego está quase saturado nesta coleção. Quando
dois significados disputam a mesma letra, o de uso mais restrito recebe
subscrito, e a Nomenclatura registra a coexistência. Onde a distinção depende
do capítulo --- o mesmo $\mu$ nomeia a viscosidade dinâmica no
Capítulo~\ref{sec:fluidos} e uma das constantes de Lamé no
Capítulo~\ref{sec:continuo} --- o significado ativo é declarado no ponto de
uso.

\textbf{Apresentação das equações.} Nenhuma equação deste volume é enunciada
sem contexto. Cada uma vem acompanhada de quatro elementos: a \emph{base} de
que decorre, nomeada explicitamente; as \emph{premissas} que precisam valer;
o \emph{caminho} da dedução, ainda que a álgebra intermediária seja
comprimida; e o \emph{domínio de validade}, isto é, onde ela deixa de servir.
Convém, ao ler, distinguir quatro estatutos diferentes, porque eles não se
questionam da mesma forma:

\begin{itemize}
	\item \textbf{Postulados} --- as leis de balanço e a segunda lei. Não se
	demonstram dentro desta teoria; aceitam-se, e toda a estrutura se apoia
	neles.
	\item \textbf{Teoremas matemáticos} --- transporte de Reynolds, tetraedro
	de Cauchy, divergência, localização. São verdadeiros dadas as suas hipóteses
	de regularidade, e é sempre nelas que se deve olhar quando um resultado
	falha.
	\item \textbf{Relações constitutivas} --- Hooke, Newton, Fourier, gás
	ideal. São modelos de material, ajustados a experimento, válidos em faixa
	limitada.
	\item \textbf{Correlações empíricas} --- fator de atrito, número de
	Nusselt, perda em leito empacotado. São ajustes a dados, com faixa de
	validade declarada e incerteza própria, e não devem ser extrapoladas.
\end{itemize}

O Capítulo~\ref{sec:continuo} é o modelo dessa apresentação e vale como
referência de leitura para os demais: nele os quatro estatutos aparecem em
sequência, do postulado de balanço às relações constitutivas.

\subsection{Relação com o Volume II}

O Volume II desta coleção aplica os fundamentos aqui apresentados a um caso
específico: o compressor de parafuso rotativo isento de óleo em plantas de ar
comprimido offshore. A relação entre os dois volumes é de especialização, não
de resumo e detalhamento: o Volume II não repete a teoria, ele escolhe as
hipóteses que a particularizam para aquela máquina e assume as consequências
dessa escolha.

Dois exemplos tornam a relação concreta. A primeira lei aplicada a um volume
de controle (Capítulo~\ref{sec:termodinamica}) é geral; ao aplicá-la a um
estágio de compressão, o Volume II fixa o processo como politrópico, adota um
índice efetivo calibrado para a máquina e passa a tratar o desvio desse índice
como sintoma de degradação. A relação entre efetividade e número de unidades
de transferência (Capítulo~\ref{sec:calor}) é geral; aplicada a um intercooler,
ela se torna um indicador de incrustação --- mas ambíguo, porque a efetividade
cai igualmente por sujeira interna e por água de resfriamento mais quente, e
desfazer a ambiguidade exige retornar ao produto $UA$, que é a grandeza
invariante às temperaturas de entrada.

A Tabela~\ref{tab:v1-c1-1} resume onde cada fundamento é retomado no Volume II.
Ela funciona como índice de navegação entre os dois volumes: um leitor que
encontre dificuldade num capítulo do Volume II localiza por ela o capítulo
deste volume que estabelece a base correspondente.

\begin{table}[H]
\centering
\caption[Mapeamento entre fundamentos e Volume II]{Fundamentos deste volume
e sua aplicação no Volume II.}
\label{tab:v1-c1-1}
\begin{tabular}{@{}lll@{}}
\toprule
\textbf{Fundamento} & \textbf{Capítulo} & \textbf{Aplicação no Volume II} \\
\midrule
Mecânica do Contínuo & \ref{sec:continuo} & Modelo de estado, balanços \\
Termodinâmica & \ref{sec:termodinamica} & Ciclos de compressão \\
Mecânica dos Fluidos & \ref{sec:fluidos} & Perda de carga, admissão \\
Transferência de Calor & \ref{sec:calor} & Intercooler, aftercooler \\
Sistemas de Transferência & \ref{sec:sistemas} & Vaso pulmão, rede, secagem \\
Eletricidade & \ref{sec:eletrica} & Motor, VSD \\
Mecânica & \ref{sec:mecanica} & Vibração, mancais \\
Instrumentação & \ref{sec:instrumentacao} & Sensores, EKF \\
\bottomrule
\end{tabular}
\end{table}

\subsection{Como usar este volume}

Cada capítulo pode ser consultado isoladamente, mas a ordem proposta segue uma
progressão lógica: dos princípios gerais de conservação às teorias específicas
de fluidos, calor, eletricidade, mecânica e medição. Três percursos de leitura
cobrem os usos previstos:

\begin{itemize}
	\item \textbf{Leitura sequencial}, do Capítulo~\ref{sec:continuo} ao
	Capítulo~\ref{sec:sintese}, para quem quer reconstruir a base inteira. É o
	percurso mais longo e o único em que a repetição deliberada dos balanços
	entre capítulos cumpre seu papel didático.
	\item \textbf{Consulta pontual}, entrando pelo Sumário ou pela
	Nomenclatura. Cada capítulo declara as hipóteses de que depende, de modo
	que uma equação consultada fora de contexto venha acompanhada da sua
	condição de validade.
	\item \textbf{Ponte a partir do Volume II}, entrando pela
	Tabela~\ref{tab:v1-c1-1}. É o percurso de quem já está resolvendo um
	problema concreto e precisa apenas do fundamento que o sustenta.
\end{itemize}

O Capítulo~\ref{sec:sintese} fecha o volume mostrando como esses fragmentos se
unem na modelagem de uma máquina rotodinâmica: é ele que reconstitui, a partir
das peças, a cadeia completa que o Volume II percorre em detalhe.

Uma advertência final sobre o uso das equações reunidas aqui. Toda equação
carrega hipóteses, e nenhuma vale fora delas. Escoamento incompressível,
propriedades constantes, regime permanente, pequenas deformações, linearidade
--- cada uma dessas simplificações é razoável em alguma faixa e desastrosa
fora dela. Quando um resultado calculado contradiz a observação, a primeira
suspeita não deve recair sobre a medição, e sim sobre a hipótese que deixou de
ser válida.


%=============================================================================
% CAPÍTULO 2 — MECÂNICA DO CONTÍNUO
%=============================================================================
\section{Mecânica do Contínuo}\label{sec:continuo}

Este capítulo estabelece a linguagem comum de todo o volume. Os capítulos
seguintes não introduzem física nova: eles especializam o que está aqui,
acrescentando hipóteses. A termodinâmica (Capítulo~\ref{sec:termodinamica}) é
o balanço de energia e de entropia sob hipóteses de equilíbrio local; a
mecânica dos fluidos (Capítulo~\ref{sec:fluidos}) é o balanço de quantidade de
movimento com uma relação constitutiva viscosa; a transferência de calor
(Capítulo~\ref{sec:calor}) é o balanço de energia com a lei de Fourier. Vale,
portanto, ler este capítulo com atenção às \emph{premissas}, porque são elas
--- e não as equações --- que determinam onde cada modelo posterior deixa de
valer.

A exposição segue a ordem lógica da disciplina, que não é a ordem histórica:
primeiro a base de mecânica clássica de que tudo decorre, com seus limites de
validade (Seção~\ref{sub:v1-c2-bases}); depois a hipótese que autoriza tratar
matéria discreta como campo contínuo (Seção~\ref{sub:v1-c2-hipotese}); depois
a descrição do movimento e da
deformação, que é pura geometria e não depende de qual material se trata
(Seções~\ref{sub:v1-c2-cinematica} e~\ref{sub:v1-c2-deformacao}); depois o
conceito de tensão, com o teorema que garante a existência do tensor
(Seção~\ref{sub:v1-c2-tensao}); depois as leis de balanço, que valem para
qualquer material (Seção~\ref{sub:v1-c2-balanco}); e só no fim as relações
constitutivas, que são onde o material específico entra
(Seção~\ref{sub:v1-c2-constitutivas}).

\subsection{Bases físicas: de onde vem a mecânica do contínuo}\label{sub:v1-c2-bases}

A mecânica do contínuo não é uma teoria autônoma. Ela é a mecânica clássica de
sistemas de partículas, reescrita em linguagem de campo depois de uma operação
de média. Começar por aí não é preciosismo: é o que explica \emph{por que} as
leis de balanço são exatamente quatro, por que a tensão é um tensor simétrico e
por que tudo isso deixa de valer em certas condições. Quem parte diretamente das
equações de campo tem de aceitar essas três coisas como dogma.

\textbf{O ponto de partida.} A mecânica clássica postula --- não demonstra ---
que existe uma classe de referenciais, ditos \textbf{inerciais}, nos quais a
taxa de variação da quantidade de movimento de um corpo iguala a força
resultante sobre ele. Para um sistema de $N$ partículas de massas $m_i$, a
soma sobre todas elas produz a Equação~\eqref{eq:v1-c2-1}:
\begin{equation}\label{eq:v1-c2-1}
\sum_k \mathbf{F}^{\text{ext}}_k =
\frac{d}{dt}\sum_{i=1}^{N} m_i \mathbf{v}_i .
\end{equation}

O detalhe decisivo da Equação~\eqref{eq:v1-c2-1} é o que \emph{não} aparece
nela: as forças internas. Pela terceira lei de Newton, elas ocorrem aos pares
iguais e opostos e se cancelam na soma; e se além disso forem centrais ---
dirigidas ao longo da reta que une as partículas ---, cancelam-se também os
seus momentos. Daí decorrem dois resultados que serão reencontrados adiante já
na forma de campo: o movimento do conjunto depende apenas das forças externas
(ancestral discreto da equação de Cauchy, Equação~\eqref{eq:v1-c2-24}) e o
balanço de momento angular não gera equação independente, apenas uma restrição
de simetria (ancestral da Equação~\eqref{eq:v1-c2-15}).

\textbf{Por que essas leis de conservação, e não outras.} As grandezas
conservadas na mecânica clássica não foram escolhidas por conveniência: cada
uma corresponde a uma simetria do espaço-tempo, correspondência estabelecida
pelo \textbf{teorema de Noether}. A Tabela~\ref{tab:v1-c2-1} resume o
argumento. Se as leis da física não mudam quando se desloca a origem do tempo,
existe uma grandeza conservada, e ela é a energia; se não mudam quando se
desloca a origem do espaço, a grandeza conservada é a quantidade de movimento;
se não mudam quando se gira o sistema de coordenadas, é o momento angular.

\begin{table}[H]
\centering
\caption[Simetrias e grandezas conservadas]{Simetrias do espaço-tempo e as
grandezas conservadas correspondentes, e onde cada balanço reaparece na forma
de campo.}
\label{tab:v1-c2-1}
\small
\begin{tabular}{@{}lll@{}}
\toprule
\textbf{Simetria} & \textbf{Grandeza conservada} & \textbf{Forma de campo} \\
\midrule
Homogeneidade do tempo & Energia & Eq.~\eqref{eq:v1-c2-25} \\
Homogeneidade do espaço & Quantidade de movimento & Eq.~\eqref{eq:v1-c2-24} \\
Isotropia do espaço & Momento angular & Eq.~\eqref{eq:v1-c2-15} \\
Postulado independente & Massa & Eq.~\eqref{eq:v1-c2-22} \\
\bottomrule
\end{tabular}
\end{table}

A última linha da Tabela~\ref{tab:v1-c2-1} merece nota, porque é uma
honestidade que costuma ser omitida: a conservação de massa \textbf{não} é
consequência de uma simetria na mecânica clássica --- é um postulado
independente, e um postulado que a relatividade posteriormente absorveu no
balanço de energia. Na faixa de velocidades e energias das máquinas tratadas
nesta coleção, tratá-la como lei separada é exato para todos os efeitos
práticos.

\textbf{A questão do referencial, que em máquinas rotativas não é acadêmica.}
A Equação~\eqref{eq:v1-c2-1} vale em referencial inercial. Ao escrevê-la num
referencial que gira com velocidade angular $\boldsymbol{\omega}$ --- o que é
natural quando se analisa um rotor ---, a aceleração absoluta se decompõe como
na Equação~\eqref{eq:v1-c2-2}:
\begin{equation}\label{eq:v1-c2-2}
\mathbf{a}_{\text{abs}} = \mathbf{a}_{\text{rel}} +
\underbrace{2\,\boldsymbol{\omega}\times\mathbf{v}_{\text{rel}}}_{\text{Coriolis}} +
\underbrace{\boldsymbol{\omega}\times(\boldsymbol{\omega}\times\mathbf{r})}_{\text{centrípeta}} +
\underbrace{\dot{\boldsymbol{\omega}}\times\mathbf{r}}_{\text{de Euler}} .
\end{equation}
Os três termos adicionais não são forças reais: são consequência de o
referencial não ser inercial. Ao passá-los para o outro lado da equação, eles
aparecem como as forças aparentes centrífuga, de Coriolis e de Euler, que
podem ser incorporadas ao termo de força de corpo $\rho\mathbf{b}$ das seções
seguintes. É essa incorporação que autoriza aplicar todo o formalismo deste
capítulo a um rotor em rotação. Ignorá-la é a origem de erros grosseiros: a
tensão centrífuga que dilata um rotor de compressor e consome parte da folga de
projeto vem exatamente do segundo termo, e não existe no referencial fixo.

\textbf{Os limites de validade da base.} Tudo o que segue neste volume é
mecânica clássica não relativística, e convém saber onde ela para:

\begin{itemize}
	\item \textbf{Velocidades muito abaixo da luz.} As velocidades periféricas
	de rotores estão na casa de $10^2\,$m\,s$^{-1}$, oito ordens de grandeza
	abaixo de $c$. Correções relativísticas são irrelevantes.
	\item \textbf{Efeitos quânticos desprezíveis.} O tratamento clássico exige
	que o comprimento de onda térmico de de Broglie seja muito menor que o
	espaçamento médio entre moléculas. Para ar em condições ambientes, o
	primeiro é da ordem de $10^{-11}\,$m e o segundo da ordem de
	$10^{-9}\,$m --- a separação é confortável.
	\item \textbf{Ausência de campos eletromagnéticos acoplados ao movimento.}
	As forças eletromagnéticas do motor são tratadas à parte, no
	Capítulo~\ref{sec:eletrica}, e entram na mecânica apenas como torque
	aplicado no eixo.
\end{itemize}

\textbf{Da partícula ao campo.} Resta o passo que dá nome ao capítulo. Um
volume de gás de tamanho de engenharia contém da ordem de $10^{19}$ moléculas
por centímetro cúbico, e nenhum modelo pode acompanhar cada trajetória. A saída
é abandonar a descrição individual e trabalhar com médias sobre volumes
pequenos em escala macroscópica e enormes em escala molecular. As leis de
balanço sobrevivem intactas a essa média --- somar sobre partículas e depois
mediar, ou mediar e depois somar, dá o mesmo resultado ---, e é essa
sobrevivência que permite reescrevê-las como equações de campo. O que a média
custa é informação: perde-se o acesso ao movimento individual, e as grandezas
que dele dependiam (viscosidade, condutividade térmica, calor específico)
deixam de ser dedutíveis e passam a ser \emph{medidas}. É por isso que a
Seção~\ref{sub:v1-c2-constitutivas} existe e não pode ser deduzida das leis de
balanço.

A validade dessa média é o assunto da próxima seção.

\subsection{A hipótese do contínuo e seus limites}\label{sub:v1-c2-hipotese}

A matéria é discreta, mas as equações que a descrevem em engenharia são
equações de campo: supõem que densidade, velocidade e temperatura tenham valor
definido em cada ponto do espaço. Conciliar as duas coisas é o papel da
\textbf{hipótese do contínuo}.

A construção é a seguinte. Associa-se a cada ponto geométrico um volume de
promediação $\Delta V$ e define-se a densidade como a razão entre a massa
contida e esse volume. Se $\Delta V$ for grande demais, a razão passa a
refletir variações macroscópicas do próprio escoamento; se for pequeno demais,
a razão flutua violentamente conforme moléculas entram e saem. Entre os dois
extremos existe --- quando existe --- um \emph{patamar}, uma faixa de $\Delta
V$ em que a razão é praticamente constante. É esse valor de patamar que se
chama densidade no ponto. Note que isso \textbf{não} é o limite matemático
$\Delta V \to 0$: esse limite não existe fisicamente, e escrevê-lo é um abuso
de notação consagrado.

O critério quantitativo que decide se o patamar existe é o \textbf{número de
Knudsen}, definido na Equação~\eqref{eq:v1-c2-3} como a razão entre o livre
caminho médio molecular $\lambda_m$ e o comprimento característico $L$ do
problema:
\begin{equation}\label{eq:v1-c2-3}
\mathrm{Kn} = \frac{\lambda_m}{L}.
\end{equation}
A hipótese do contínuo é aceita para $\mathrm{Kn} \lesssim 10^{-2}$. Entre
$10^{-2}$ e $10^{-1}$ o meio ainda é contínuo no interior, mas a condição de
não deslizamento na parede falha (regime de deslizamento); acima disso, as
equações deste capítulo perdem validade e é preciso recorrer à teoria cinética
dos gases.

Esse limite não é acadêmico nas máquinas tratadas nesta coleção. O livre
caminho médio do ar a $1\,$bar e $20\,^{\circ}$C é da ordem de $70\,$nm, e o
vazamento pelas folgas internas de um compressor de parafuso ocorre em fendas
de dezenas de micrometros --- $\mathrm{Kn}$ da ordem de $10^{-3}$, ainda
confortavelmente contínuo. Como $\lambda_m$ é inversamente proporcional à
pressão, porém, a mesma folga sob vácuo sairia da faixa. A hipótese vale aqui
por causa da pressão de operação, não por causa da geometria, e é assim que ela
deve ser verificada: pelo número, não pela familiaridade.

Uma segunda premissa acompanha a primeira e é frequentemente esquecida: a
existência do patamar exige \textbf{separação de escalas}, isto é, que o
comprimento macroscópico seja muito maior que a escala da microestrutura. Em um
meio poroso, num leito de dessecante ou numa camada de incrustação, a
microestrutura tem dimensão milimétrica, e as equações de campo só fazem
sentido depois de uma segunda promediação, sobre um volume elementar
representativo. É por isso que o escoamento em leito empacotado não é descrito
por Navier-Stokes diretamente, e sim por relações promediadas, tratadas no
Capítulo~\ref{sec:fluidos}.

\subsection{Cinemática: descrições material e espacial}\label{sub:v1-c2-cinematica}

Aceita a hipótese do contínuo, o corpo passa a ser um conjunto contínuo de
\emph{pontos materiais}. Escolhe-se uma \textbf{configuração de referência} ---
em geral o corpo indeformado, num instante $t_0$ --- e identifica-se cada ponto
material pela posição $\mathbf{X}$ que ele ocupa nela. O movimento do corpo é
então a família de aplicações da Equação~\eqref{eq:v1-c2-4}, que leva cada
ponto material à sua posição atual $\mathbf{x}$:
\begin{equation}\label{eq:v1-c2-4}
\mathbf{x} = \boldsymbol{\chi}(\mathbf{X},t).
\end{equation}

A Equação~\eqref{eq:v1-c2-4} carrega três premissas que raramente são
enunciadas e que valem a pena explicitar, porque são elas que excluem
comportamentos fisicamente impossíveis:

\begin{enumerate}
	\item $\boldsymbol{\chi}$ é \textbf{bijetiva} em cada instante --- dois
	pontos materiais distintos não podem ocupar a mesma posição, e um ponto
	material não pode estar em dois lugares. É a impenetrabilidade da matéria.
	\item $\boldsymbol{\chi}$ é \textbf{continuamente diferenciável} em
	$\mathbf{X}$ e em $t$ --- o corpo não se rompe nem se dobra sobre si mesmo.
	Trincas e superfícies de descontinuidade violam essa premissa e exigem
	tratamento à parte.
	\item A orientação se conserva, condição que a
	Seção~\ref{sub:v1-c2-deformacao} escreverá como $J>0$ --- o material não
	pode ser ``virado pelo avesso''.
\end{enumerate}

Daí decorrem duas maneiras de descrever qualquer grandeza, e a escolha entre
elas organiza toda a modelagem posterior:

\textbf{Descrição material (ou lagrangeana):} a grandeza é função de
$(\mathbf{X},t)$ --- acompanha-se um ponto material específico. É a descrição
natural para sólidos, em que a configuração de referência é identificável e o
interesse é o deslocamento a partir dela. O campo de deslocamento é
$\mathbf{u}(\mathbf{X},t) = \boldsymbol{\chi}(\mathbf{X},t) - \mathbf{X}$.

\textbf{Descrição espacial (ou euleriana):} a grandeza é função de
$(\mathbf{x},t)$ --- observa-se um ponto fixo do espaço e registra-se o que
passa por ele. É a descrição natural para fluidos, em que não há configuração
de referência útil e ninguém deseja rastrear uma partícula de gás
individualmente.

Um trem de compressão exige as duas ao mesmo tempo: o rotor que dilata e flete
é tratado materialmente, o gás que o atravessa é tratado espacialmente, e a
folga entre ambos --- a grandeza que o Volume II monitora --- é justamente onde
as duas descrições se encontram.

A ponte entre elas é a \textbf{derivada material}, que responde a uma pergunta
euleriana com significado lagrangeano: qual a taxa de variação de uma grandeza
\emph{para o ponto material} que neste instante passa por aqui? Seja
$\varphi(\mathbf{x},t)$ um campo espacial qualquer. Ao longo da trajetória de um
ponto material, $\mathbf{x} = \boldsymbol{\chi}(\mathbf{X},t)$, e a regra da
cadeia dá a Equação~\eqref{eq:v1-c2-5}:
\begin{equation}\label{eq:v1-c2-5}
\frac{D\varphi}{Dt} \equiv
\left.\frac{\partial \varphi}{\partial t}\right|_{\mathbf{X}} =
\underbrace{\frac{\partial \varphi}{\partial t}}_{\text{local}} +
\underbrace{\mathbf{v}\cdot\nabla\varphi}_{\text{convectivo}},
\qquad
\mathbf{v} = \left.\frac{\partial \boldsymbol{\chi}}{\partial t}\right|_{\mathbf{X}}.
\end{equation}

A leitura dos dois termos é diretamente instrumental. Um termopar fixo na
tubulação mede $\partial T/\partial t$: a variação no seu ponto. A partícula de
gás que passa por ele experimenta $DT/Dt$, que inclui o aquecimento sofrido
apenas por ter se deslocado para uma região mais quente. Em regime permanente o
primeiro termo é nulo e o segundo não --- um escoamento estacionário acelera e
aquece partículas o tempo todo. Confundir os dois é a origem mais comum de erro
de sinal em balanços de máquinas em operação contínua.

\subsection{Deformação: gradiente, medidas finitas e linearização}\label{sub:v1-c2-deformacao}

Deformação é uma noção geométrica: mede quanto as distâncias entre pontos
materiais vizinhos mudaram. Nada nesta seção depende do material.

O objeto fundamental é o \textbf{gradiente de deformação} $\mathbf{F}$,
definido na Equação~\eqref{eq:v1-c2-6}, que transporta um segmento
infinitesimal da configuração de referência para a atual:
\begin{equation}\label{eq:v1-c2-6}
\mathbf{F} = \frac{\partial \mathbf{x}}{\partial \mathbf{X}},
\qquad F_{iJ} = \frac{\partial x_i}{\partial X_J},
\qquad d\mathbf{x} = \mathbf{F}\,d\mathbf{X}.
\end{equation}

Seu determinante, o \textbf{jacobiano} da Equação~\eqref{eq:v1-c2-7}, é a razão
entre o volume atual e o de referência de um mesmo elemento material:
\begin{equation}\label{eq:v1-c2-7}
J = \det \mathbf{F}, \qquad dV = J\,dV_0.
\end{equation}
A premissa de conservação de orientação da seção anterior é agora quantitativa:
$J>0$ sempre, e $J=1$ caracteriza movimento isocórico (que preserva volume).
A idealização de \emph{incompressibilidade}, usada à exaustão nos capítulos de
fluidos, é exatamente a imposição $J \equiv 1$.

$\mathbf{F}$ sozinho, porém, não serve como medida de deformação, e a razão é
essencial: ele não distingue deformação de rotação de corpo rígido. Um rotor
que gira sem se deformar tem $\mathbf{F}$ diferente da identidade, embora
nenhuma distância material tenha mudado. A separação formal é dada pelo
\textbf{teorema da decomposição polar}, Equação~\eqref{eq:v1-c2-8}, válido para
qualquer $\mathbf{F}$ com $J>0$:
\begin{equation}\label{eq:v1-c2-8}
\mathbf{F} = \mathbf{R}\,\mathbf{U} = \mathbf{V}\,\mathbf{R},
\end{equation}
onde $\mathbf{R}$ é ortogonal próprio (rotação pura) e $\mathbf{U}$,
$\mathbf{V}$ são simétricos positivos definidos (estiramento puro).

Como não se deseja carregar $\mathbf{R}$, constroem-se medidas que o eliminem.
O produto $\mathbf{C} = \mathbf{F}^{\mathsf{T}}\mathbf{F} = \mathbf{U}^2$
--- tensor de Cauchy-Green à direita --- é imune à rotação, e dele se define o
tensor de deformação de \textbf{Green-Lagrange}, Equação~\eqref{eq:v1-c2-9}:
\begin{equation}\label{eq:v1-c2-9}
\mathbf{E}_L = \tfrac{1}{2}\left(\mathbf{C} - \mathbf{I}\right)
= \tfrac{1}{2}\left(\mathbf{F}^{\mathsf{T}}\mathbf{F} - \mathbf{I}\right).
\end{equation}
Por construção, $\mathbf{E}_L = \mathbf{0}$ para qualquer movimento de corpo
rígido --- que é precisamente a propriedade exigida de uma medida de
deformação.

Escrevendo $\mathbf{F} = \mathbf{I} + \nabla_{\!X}\mathbf{u}$ e substituindo na
Equação~\eqref{eq:v1-c2-9}, obtém-se a forma exata da
Equação~\eqref{eq:v1-c2-10}:
\begin{equation}\label{eq:v1-c2-10}
\mathbf{E}_L = \tfrac{1}{2}\left( \nabla_{\!X}\mathbf{u} +
\nabla_{\!X}\mathbf{u}^{\mathsf{T}} +
\nabla_{\!X}\mathbf{u}^{\mathsf{T}}\,\nabla_{\!X}\mathbf{u} \right).
\end{equation}

Aqui entra a única aproximação desta seção, e ela merece nome e critério.
Sob a \textbf{hipótese de pequenas deformações}, $\|\nabla\mathbf{u}\| \ll 1$,
o termo quadrático da Equação~\eqref{eq:v1-c2-10} é de segunda ordem e pode ser
descartado. Pela mesma hipótese, $\partial/\partial X_J \approx \partial/\partial
x_j$ --- as configurações de referência e atual quase coincidem, e deixa de ser
necessário distinguir em qual delas se deriva. O resultado é o \textbf{tensor
de deformação infinitesimal} da Equação~\eqref{eq:v1-c2-11}:
\begin{equation}\label{eq:v1-c2-11}
\varepsilon_{ij} = \frac{1}{2}\left( \frac{\partial u_i}{\partial x_j} +
\frac{\partial u_j}{\partial x_i} \right).
\end{equation}

Duas consequências práticas dessa linearização, ambas relevantes para máquinas
rotativas. A primeira: como o termo descartado é quadrático, o erro relativo de
$\varepsilon$ é da ordem do próprio $\|\nabla\mathbf{u}\|$; para deformações
elásticas de aço, da ordem de $10^{-3}$, o erro é desprezível, e a Equação
\eqref{eq:v1-c2-11} é plenamente adequada à dilatação térmica de rotores e
carcaças. A segunda, e mais traiçoeira: a hipótese restringe o \emph{gradiente
de deslocamento}, não o deslocamento. Uma rotação de corpo rígido finita tem
gradiente de deslocamento grande, e a Equação~\eqref{eq:v1-c2-11} aplicada
diretamente a um rotor girando prevê deformação onde não há nenhuma. Em
máquinas rotativas, portanto, a deformação infinitesimal só é legítima quando
escrita num referencial que acompanha a rotação.

É útil, por fim, separar $\varepsilon$ em duas partes com significados físicos
distintos, como faz a Equação~\eqref{eq:v1-c2-12}: uma parcela esférica, que
representa mudança de volume sem mudança de forma, e uma parcela desviadora
$\mathbf{e}$, que representa mudança de forma a volume constante:
\begin{equation}\label{eq:v1-c2-12}
\boldsymbol{\varepsilon} = \underbrace{\tfrac{1}{3}\,\varepsilon_{kk}\,\mathbf{I}}_{\text{volumétrica}}
+ \underbrace{\mathbf{e}}_{\text{desviadora}},
\qquad \mathrm{tr}\,\mathbf{e} = 0 .
\end{equation}
A separação não é cosmética: materiais dúcteis escoam por cisalhamento, isto é,
respondem à parte desviadora, e é essa observação que sustenta o critério de
von Mises adiante e o dimensionamento à fadiga do Capítulo~\ref{sec:mecanica}.

Para fluidos, a grandeza de interesse não é a deformação acumulada --- que não
tem sentido, pois um fluido não guarda memória da configuração de referência
--- e sim a sua \emph{taxa}. Decompondo o gradiente de velocidade em partes
simétrica e antissimétrica, a Equação~\eqref{eq:v1-c2-13} define o tensor
\textbf{taxa de deformação} $\mathbf{S}$ e o tensor \textbf{taxa de rotação}
(ou spin) $\mathbf{W}$:
\begin{equation}\label{eq:v1-c2-13}
\nabla\mathbf{v} = \mathbf{S} + \mathbf{W},
\qquad
\mathbf{S} = \tfrac{1}{2}\!\left(\nabla\mathbf{v} + \nabla\mathbf{v}^{\mathsf{T}}\right),
\qquad
\mathbf{W} = \tfrac{1}{2}\!\left(\nabla\mathbf{v} - \nabla\mathbf{v}^{\mathsf{T}}\right).
\end{equation}
$\mathbf{W}$ descreve rotação local instantânea de corpo rígido e, pela mesma
razão que $\mathbf{R}$ foi eliminado acima, não pode gerar tensão viscosa ---
premissa que a Seção~\ref{sub:v1-c2-constitutivas} usará ao escrever a lei de
Newton da viscosidade em função de $\mathbf{S}$ apenas.

\subsection{Tensão: o teorema de Cauchy}\label{sub:v1-c2-tensao}

As forças que atuam sobre uma porção de contínuo dividem-se em dois tipos:
\textbf{forças de corpo}, que agem à distância sobre a massa (gravidade, forças
eletromagnéticas, e as forças de inércia aparentes num referencial girante), e
\textbf{forças de superfície}, que a vizinhança exerce por contato através da
fronteira. As primeiras se escrevem como $\rho\mathbf{b}$ por unidade de
volume. As segundas exigem mais cuidado, e é delas que nasce o conceito de
tensão.

Define-se o \textbf{vetor tensão} $\mathbf{t}_{(\mathbf{n})}$ como a força de
contato por unidade de área que o material do lado para o qual aponta o vetor
normal unitário $\mathbf{n}$ exerce sobre o material do outro lado. O ponto
crucial --- e a razão de todo o edifício funcionar --- é que
$\mathbf{t}_{(\mathbf{n})}$ depende do ponto e da orientação do corte: um mesmo
ponto material sujeito a um mesmo estado de solicitação apresenta trações
diferentes em planos diferentes. Em princípio, seriam necessárias infinitas
informações, uma por orientação possível.

O \textbf{teorema do tetraedro de Cauchy} mostra que três bastam. O argumento é
o seguinte. Isola-se um tetraedro infinitesimal com três faces sobre os planos
coordenados, de áreas $A\,n_1$, $A\,n_2$, $A\,n_3$, e a quarta face de área $A$
e normal $\mathbf{n}$. Aplica-se o balanço de quantidade de movimento a esse
tetraedro. As forças de superfície são proporcionais às áreas, isto é, a $h^2$,
onde $h$ é a altura do tetraedro; já a força de corpo e o termo de inércia são
proporcionais ao volume, isto é, a $h^3$. Fazendo $h \to 0$, os termos de
volume se anulam mais rapidamente que os de superfície, e o que sobra é uma
relação puramente algébrica entre as quatro trações --- a
Equação~\eqref{eq:v1-c2-14}:
\begin{equation}\label{eq:v1-c2-14}
\mathbf{t}_{(\mathbf{n})} = \boldsymbol{\sigma}\,\mathbf{n},
\qquad t_i = \sigma_{ij}\,n_j .
\end{equation}

O conteúdo do teorema não é a fórmula, e sim a existência do objeto
$\boldsymbol{\sigma}$: as trações em todas as infinitas orientações possíveis
ficam determinadas por nove números, e esses nove números se transformam como
um tensor de segunda ordem. É esse resultado que autoriza falar em ``estado de
tensão em um ponto''. As premissas empregadas foram apenas duas, e ambas devem
ser lembradas: a tração depende da superfície apenas através de $\mathbf{n}$
(postulado de Cauchy --- não depende, por exemplo, da curvatura do corte), e a
aceleração e a força de corpo são limitadas, o que permite descartar os termos
de volume.

Um segundo resultado reduz nove para seis. Aplicando o balanço de
\textbf{momento angular} a um elemento material e supondo que não existam
momentos distribuídos de corpo nem tensões de momento --- premissa que define o
chamado contínuo não polar, ou de Boltzmann ---, obtém-se a simetria enunciada
na Equação~\eqref{eq:v1-c2-15}:
\begin{equation}\label{eq:v1-c2-15}
\boldsymbol{\sigma} = \boldsymbol{\sigma}^{\mathsf{T}},
\qquad \sigma_{ij} = \sigma_{ji}.
\end{equation}
A simetria é uma consequência de uma lei de balanço, não uma definição, e ela
falha em meios com microestrutura polarizável. Nos materiais desta coleção,
vale.

Sendo simétrico e real, $\boldsymbol{\sigma}$ admite três autovalores reais
$\sigma_1 \ge \sigma_2 \ge \sigma_3$ --- as \textbf{tensões principais} --- e
três direções principais mutuamente ortogonais, nas quais a tração é puramente
normal, sem cisalhamento. Os coeficientes do polinômio característico são os
invariantes de $\boldsymbol{\sigma}$, e são invariantes justamente por não
dependerem do sistema de coordenadas escolhido: qualquer critério de falha
fisicamente admissível tem de ser escrito em função deles, e não das
componentes individuais.

Pela mesma razão que motivou a Equação~\eqref{eq:v1-c2-12}, separa-se a tensão
em pressão hidrostática e parte desviadora, como na
Equação~\eqref{eq:v1-c2-16}:
\begin{equation}\label{eq:v1-c2-16}
\boldsymbol{\sigma} = -p\,\mathbf{I} + \mathbf{s},
\qquad
p = -\tfrac{1}{3}\,\mathrm{tr}\,\boldsymbol{\sigma},
\qquad
\mathrm{tr}\,\mathbf{s} = 0 .
\end{equation}
A evidência experimental de que metais dúcteis não escoam sob pressão
hidrostática elevada, por maior que ela seja, leva a construir o critério de
escoamento apenas sobre $\mathbf{s}$. A medida escalar usual é a \textbf{tensão
equivalente de von Mises} da Equação~\eqref{eq:v1-c2-17}:
\begin{equation}\label{eq:v1-c2-17}
\sigma_{VM} = \sqrt{\tfrac{3}{2}\,\mathbf{s}\!:\!\mathbf{s}}
= \sqrt{\tfrac{1}{2}\left[(\sigma_1-\sigma_2)^2 + (\sigma_2-\sigma_3)^2 +
(\sigma_3-\sigma_1)^2\right]} .
\end{equation}
É $\sigma_{VM}$ --- e não uma componente isolada --- que se compara à tensão
admissível do material no dimensionamento de eixos e engrenagens e na
contagem de dano por fadiga do Capítulo~\ref{sec:mecanica}.

\subsection{Leis de balanço}\label{sub:v1-c2-balanco}

As leis de balanço são \textbf{postulados}, não teoremas: não se demonstram a
partir de algo mais fundamental dentro desta teoria. O que se demonstra é como
convertê-las de uma porção de matéria que se move para uma região fixa do
espaço --- e é essa conversão, e não a lei em si, que produz as equações
efetivamente usadas em engenharia.

Seja $\mathcal{B}$ uma grandeza extensiva qualquer e $\beta$ a grandeza
correspondente por unidade de massa. O postulado de balanço, escrito para um
\textbf{volume material} $V_m(t)$ --- uma porção fixa de matéria, que se move e
se deforma com o corpo ---, afirma que a taxa de variação do conteúdo iguala o
que entra por fluxo não convectivo através da fronteira mais o que é gerado
internamente, conforme a Equação~\eqref{eq:v1-c2-18}:
\begin{equation}\label{eq:v1-c2-18}
\frac{D}{Dt}\int_{V_m(t)} \rho\,\beta \; dV =
-\oint_{S_m(t)} \mathbf{j}_\beta \cdot \mathbf{n} \; dA +
\int_{V_m(t)} \dot{\mathcal{G}}_\beta \; dV .
\end{equation}
Aqui $\mathbf{j}_\beta$ é o fluxo que atravessa a fronteira por mecanismo
difusivo ou de contato --- o fluxo de calor, a tração superficial --- e
$\dot{\mathcal{G}}_\beta$ é a taxa de geração por unidade de volume.

O obstáculo prático é que $V_m(t)$ se move: os limites da integral dependem do
tempo, e a derivada não passa para dentro. Quem remove o obstáculo é o
\textbf{teorema de transporte de Reynolds}, Equação~\eqref{eq:v1-c2-19}, que
relaciona a derivada material da integral sobre o volume material à derivada
parcial da integral sobre um \textbf{volume de controle} $V_c$ fixo no espaço,
escolhido de modo a coincidir com $V_m$ no instante considerado:
\begin{equation}\label{eq:v1-c2-19}
\frac{D}{Dt}\int_{V_m(t)} \rho\,\beta \; dV =
\underbrace{\frac{\partial}{\partial t}\int_{V_c} \rho\,\beta \; dV}_{\text{acúmulo}} +
\underbrace{\oint_{S_c} \rho\,\beta \,(\mathbf{v}\cdot\mathbf{n}) \; dA}_{\text{fluxo convectivo líquido}} .
\end{equation}

A dedução é curta e vale registrá-la, porque é dela que sai a estrutura
``acúmulo mais o que sai menos o que entra'' de todos os balanços posteriores.
Muda-se a variável de integração para a configuração de referência, onde os
limites são fixos, usando $dV = J\,dV_0$ da Equação~\eqref{eq:v1-c2-7}; a
derivada passa então para dentro da integral. Aparece $DJ/Dt$, dado pela
\textbf{fórmula de expansão de Euler}, Equação~\eqref{eq:v1-c2-20}:
\begin{equation}\label{eq:v1-c2-20}
\frac{DJ}{Dt} = J \, (\nabla\cdot\mathbf{v}).
\end{equation}
Retornando à configuração atual e reagrupando com a
Equação~\eqref{eq:v1-c2-5}, o integrando resultante é
$\partial(\rho\beta)/\partial t + \nabla\cdot(\rho\beta\,\mathbf{v})$; aplicar
o \textbf{teorema da divergência} ao segundo termo produz a integral de
superfície da Equação~\eqref{eq:v1-c2-19}. A interpretação de
$\nabla\cdot\mathbf{v}$ como taxa de dilatação volumétrica, aliás, é leitura
direta da Equação~\eqref{eq:v1-c2-20}.

Combinando as Equações~\eqref{eq:v1-c2-18} e~\eqref{eq:v1-c2-19}, aplicando o
teorema da divergência também ao termo de fluxo e observando que o resultado
vale para \emph{qualquer} volume de controle, chega-se ao passo final. Se uma
integral se anula sobre todo subvolume e o integrando é contínuo, então o
integrando é identicamente nulo --- esse é o \textbf{teorema da localização},
cuja premissa de continuidade é exatamente o que falha em choques e interfaces.
O resultado é a forma local, ou diferencial, da Equação~\eqref{eq:v1-c2-21}:
\begin{equation}\label{eq:v1-c2-21}
\frac{\partial (\rho\beta)}{\partial t} + \nabla\cdot(\rho\beta\,\mathbf{v})
= -\nabla\cdot\mathbf{j}_\beta + \dot{\mathcal{G}}_\beta .
\end{equation}

A Equação~\eqref{eq:v1-c2-21} é o resultado central deste capítulo. Todas as
equações de conservação usadas no volume saem dela por substituição, conforme
resume a Tabela~\ref{tab:v1-c2-2}. Vale notar o que \emph{não} se supôs para
chegar até aqui: nada foi dito sobre o material. A equação vale igualmente para
aço, ar, óleo ou dessecante.

\begin{table}[H]
\centering
\caption[Especializações da forma local de balanço]{Especializações da forma
local de balanço, Equação~\eqref{eq:v1-c2-21}, para cada grandeza conservada.}
\label{tab:v1-c2-2}
\small
\begin{tabular}{@{}lllll@{}}
\toprule
$\mathcal{B}$ & $\beta$ & $\mathbf{j}_\beta$ & $\dot{\mathcal{G}}_\beta$ & \textbf{Lei resultante} \\
\midrule
Massa & $1$ & $\mathbf{0}$ & $0$ & Continuidade, Eq.~\eqref{eq:v1-c2-22} \\
Qtd. de movimento & $\mathbf{v}$ & $-\boldsymbol{\sigma}$ & $\rho\mathbf{b}$ & Cauchy, Eq.~\eqref{eq:v1-c2-24} \\
Momento angular & $\mathbf{x}\times\mathbf{v}$ & --- & --- & Simetria de $\boldsymbol{\sigma}$, Eq.~\eqref{eq:v1-c2-15} \\
Energia total & $e + v^2/2$ & $\mathbf{q} - \boldsymbol{\sigma}\mathbf{v}$ & $\rho\,\mathbf{b}\cdot\mathbf{v}$ & Primeira lei, Eq.~\eqref{eq:v1-c2-25} \\
Entropia & $s$ & $\mathbf{q}/T$ & $\dot{\mathcal{G}}_s \ge 0$ & Segunda lei, Eq.~\eqref{eq:v1-c2-26} \\
\bottomrule
\end{tabular}
\end{table}

\textbf{Conservação de massa.} Com $\beta = 1$, sem fluxo difusivo e sem
geração, a Equação~\eqref{eq:v1-c2-21} fornece a equação da continuidade,
Equação~\eqref{eq:v1-c2-22}:
\begin{equation}\label{eq:v1-c2-22}
\frac{\partial \rho}{\partial t} + \nabla \cdot (\rho \mathbf{v}) = 0.
\end{equation}
Para massa específica constante, ela se reduz à condição de incompressibilidade
da Equação~\eqref{eq:v1-c2-23}, que é a mesma imposição $J\equiv1$ escrita em
termos de taxa:
\begin{equation}\label{eq:v1-c2-23}
\nabla \cdot \mathbf{v} = 0.
\end{equation}

\textbf{Conservação da quantidade de movimento.} Com $\beta = \mathbf{v}$,
fluxo $-\boldsymbol{\sigma}$ (o sinal decorre de $\mathbf{t}=\boldsymbol\sigma
\mathbf{n}$ ser entrada de quantidade de movimento) e geração $\rho\mathbf{b}$,
e usando a Equação~\eqref{eq:v1-c2-22} para simplificar o lado esquerdo,
obtém-se a \textbf{equação do movimento de Cauchy},
Equação~\eqref{eq:v1-c2-24}:
\begin{equation}\label{eq:v1-c2-24}
\rho \frac{D\mathbf{v}}{Dt} = \nabla \cdot \boldsymbol{\sigma} + \rho \mathbf{b}.
\end{equation}
É a segunda lei de Newton escrita para o contínuo. No caso particular de
aceleração nula, $D\mathbf{v}/Dt = \mathbf{0}$, ela se reduz às equações de
equilíbrio estático, $\partial\sigma_{ij}/\partial x_j + \rho b_i = 0$, usadas
no dimensionamento estrutural de carcaças e vasos.

\textbf{Conservação de energia.} Com $\beta = e + v^2/2$, fluxo
$\mathbf{q} - \boldsymbol{\sigma}\mathbf{v}$ (condução mais trabalho das forças
de superfície) e geração $\rho\,\mathbf{b}\cdot\mathbf{v}$, resulta a
Equação~\eqref{eq:v1-c2-25}:
\begin{equation}\label{eq:v1-c2-25}
\rho \frac{D}{Dt}\left( e + \frac{v^2}{2} \right) =
\nabla \cdot (\boldsymbol{\sigma}\mathbf{v}) - \nabla \cdot \mathbf{q} +
\rho\, \mathbf{b}\cdot\mathbf{v},
\end{equation}
onde $e$ é a energia interna específica e $\mathbf{q}$ o vetor fluxo de calor.
É desta equação --- e não de um postulado independente --- que o
Capítulo~\ref{sec:termodinamica} extrai a primeira lei na forma usual para
volumes de controle.

\textbf{Balanço de entropia.} Com $\beta = s$ e fluxo $\mathbf{q}/T$, a
Equação~\eqref{eq:v1-c2-21} produz a \textbf{desigualdade de Clausius-Duhem},
Equação~\eqref{eq:v1-c2-26}:
\begin{equation}\label{eq:v1-c2-26}
\rho \frac{Ds}{Dt} + \nabla\cdot\!\left(\frac{\mathbf{q}}{T}\right)
= \dot{\mathcal{G}}_s \ge 0 .
\end{equation}
A segunda lei da termodinâmica aparece assim como o único balanço deste
capítulo cujo termo de geração tem sinal imposto. Toda a análise de exergia do
Capítulo~\ref{sec:termodinamica} é consequência dessa única restrição de sinal,
e é ela que impede que uma relação constitutiva proposta seja fisicamente
inadmissível.

\subsection{Relações constitutivas}\label{sub:v1-c2-constitutivas}

As leis de balanço não bastam para resolver problema algum, e a razão é
aritmética. Um simples inventário deixa isso evidente: as incógnitas são
$\rho$, as três componentes de $\mathbf{v}$, as seis componentes independentes
de $\boldsymbol{\sigma}$, a energia interna $e$, as três componentes de
$\mathbf{q}$ e a temperatura $T$ --- quinze ao todo. As equações disponíveis
são a continuidade (uma), o movimento de Cauchy (três) e a energia (uma):
cinco. Faltam dez.

As dez equações que faltam são as \textbf{relações constitutivas}, e é
exatamente nelas que o material entra. Elas não são deduzidas das leis de
balanço; são modelos, ajustados a experimento, sujeitos a três princípios que
restringem as formas admissíveis:

\begin{itemize}
	\item \textbf{Determinismo:} a tensão no instante atual depende da história
	do movimento, não do futuro.
	\item \textbf{Ação local:} a tensão em um ponto depende do movimento numa
	vizinhança infinitesimal dele, não do corpo inteiro.
	\item \textbf{Objetividade material} (indiferença ao referencial): a
	resposta do material não pode depender do observador. É esse princípio que
	proíbe $\boldsymbol{\sigma}$ de depender de $\mathbf{F}$ diretamente ou de
	$\mathbf{W}$, e obriga a usar medidas imunes à rotação como $\mathbf{C}$ e
	$\mathbf{S}$ --- o que fecha, agora com justificativa formal, as escolhas
	feitas na Seção~\ref{sub:v1-c2-deformacao}.
\end{itemize}

\textbf{Sólido elástico linear e isotrópico.} Supondo resposta instantânea,
linear em $\boldsymbol{\varepsilon}$ e sem direção preferencial, a forma mais
geral admissível é a \textbf{lei de Hooke generalizada},
Equação~\eqref{eq:v1-c2-27}, com apenas duas constantes independentes:
\begin{equation}\label{eq:v1-c2-27}
\sigma_{ij} = \lambda_L \, \varepsilon_{kk} \, \delta_{ij} +
2G \, \varepsilon_{ij},
\end{equation}
onde $\lambda_L$ é a primeira constante de Lamé, $G$ o módulo de cisalhamento
(igual à segunda constante de Lamé) e $\delta_{ij}$ o delta de Kronecker. Que
sejam exatamente duas é consequência da isotropia, não uma simplificação
adicional. Na prática de engenharia mede-se o par módulo de Young e coeficiente
de Poisson, relacionado ao primeiro pela Equação~\eqref{eq:v1-c2-28}:
\begin{equation}\label{eq:v1-c2-28}
\lambda_L = \frac{E\,\nu}{(1+\nu)(1-2\nu)},
\qquad
G = \frac{E}{2(1+\nu)} .
\end{equation}
O denominador $(1-2\nu)$ explicita por que $\nu \to 1/2$ corresponde a material
incompressível: $\lambda_L$ diverge, e qualquer tentativa de variar o volume
exige tensão infinita.

\textbf{Fluido newtoniano.} Um fluido em repouso não sustenta cisalhamento, de
modo que a tensão se reduz à pressão termodinâmica; em movimento, acrescenta-se
uma parcela viscosa que, pelo princípio da objetividade, só pode depender de
$\mathbf{S}$, conforme as Equações~\eqref{eq:v1-c2-29} e~\eqref{eq:v1-c2-30}:
\begin{equation}\label{eq:v1-c2-29}
\sigma_{ij} = -p\,\delta_{ij} + \tau_{ij},
\end{equation}
\begin{equation}\label{eq:v1-c2-30}
\tau_{ij} = 2\mu\,S_{ij} + \lambda_v \,(\nabla\cdot\mathbf{v})\,\delta_{ij},
\end{equation}
com $\mu$ a viscosidade dinâmica e $\lambda_v$ a segunda viscosidade. Esta
última é de medição difícil, e adota-se quase universalmente a \textbf{hipótese
de Stokes}, Equação~\eqref{eq:v1-c2-31}:
\begin{equation}\label{eq:v1-c2-31}
\lambda_v = -\tfrac{2}{3}\,\mu
\qquad\Longrightarrow\qquad
\tau_{ij} = 2\mu\,S_{ij} - \tfrac{2}{3}\mu\,(\nabla\cdot\mathbf{v})\,\delta_{ij}.
\end{equation}
A hipótese de Stokes equivale a supor viscosidade volumétrica nula, ou seja,
que a compressão não dissipa por si mesma. É exata para gases monoatômicos e
uma aproximação para os demais; seu efeito é desprezível em escoamentos com
$\nabla\cdot\mathbf{v}$ pequeno, e por isso raramente compromete a modelagem de
um trem de compressão, onde a dissipação relevante é de cisalhamento. Convém,
ainda assim, saber que é hipótese, e não lei.

Substituindo as Equações~\eqref{eq:v1-c2-29} e~\eqref{eq:v1-c2-31} na
Equação~\eqref{eq:v1-c2-24} obtêm-se as equações de \textbf{Navier-Stokes} para
fluido compressível, retomadas no Capítulo~\ref{sec:fluidos}.

\textbf{Condução de calor.} Fecham-se as três incógnitas de $\mathbf{q}$ com a
\textbf{lei de Fourier}, Equação~\eqref{eq:v1-c2-32}, cujo sinal negativo é
imposto pela Equação~\eqref{eq:v1-c2-26}: calor fluindo do frio para o quente
geraria entropia negativa.
\begin{equation}\label{eq:v1-c2-32}
\mathbf{q} = -k\,\nabla T .
\end{equation}
Com ela, e com uma equação de estado calórica que ligue $e$ a $T$, o sistema
fica fechado. A Tabela~\ref{tab:v1-c2-3} reúne as relações constitutivas usadas
neste volume.

\begin{table}[H]
\centering
\caption[Relações constitutivas usadas no volume]{Relações constitutivas para
sólidos elásticos, fluidos newtonianos e condução de calor.}
\label{tab:v1-c2-3}
\small
\resizebox{\textwidth}{!}{%
\begin{tabular}{@{}llll@{}}
\toprule
\textbf{Material/mecanismo} & \textbf{Relação constitutiva} & \textbf{Parâmetros} & \textbf{Premissa} \\
\midrule
Sólido elástico linear & $\sigma_{ij} = \lambda_L\varepsilon_{kk}\delta_{ij} + 2G\varepsilon_{ij}$ & $E$, $\nu$ & Isotropia, $\|\nabla\mathbf{u}\|\ll1$ \\
Fluido newtoniano & $\tau_{ij} = 2\mu S_{ij} + \lambda_v(\nabla\cdot\mathbf{v})\delta_{ij}$ & $\mu$, $\lambda_v$ & Linearidade em $\mathbf{S}$ \\
Fluido newtoniano (Stokes) & $\tau_{ij} = 2\mu S_{ij} - \frac{2}{3}\mu(\nabla\cdot\mathbf{v})\delta_{ij}$ & $\mu$ & Viscosidade volumétrica nula \\
Fluido inviscido & $\sigma_{ij} = -p\,\delta_{ij}$ & nenhum & $\mu \to 0$ \\
Condução de calor & $\mathbf{q} = -k\nabla T$ & $k$ & Isotropia, meio sem memória \\
\bottomrule
\end{tabular}}
\end{table}


%=============================================================================
% CAPÍTULO 3 — TERMODINÂMICA
%=============================================================================
\section{Termodinâmica}\label{sec:termodinamica}

O Capítulo~\ref{sec:continuo} chegou a duas equações de campo --- o balanço de
energia, Equação~\eqref{eq:v1-c2-25}, e o de entropia,
Equação~\eqref{eq:v1-c2-26} --- que valem ponto a ponto e não mencionam
temperatura, entalpia nem calor específico. A termodinâmica de engenharia
trabalha com essas grandezas e com \emph{estados}, não com campos. Este
capítulo faz a ponte: mostra sob qual hipótese o formalismo de estados é
legítimo, deriva dele a primeira e a segunda leis na forma usada em máquinas, e
termina no trabalho de compressão --- a equação que o Volume II usa para cada
estágio.

A ordem é deliberada. Quase todo curso introdutório apresenta a primeira lei
como postulado e a entalpia como definição conveniente. Aqui as duas
\emph{aparecem}: a primeira lei como especialização de um balanço já
estabelecido, e a entalpia como o agrupamento que absorve o trabalho de
fronteira nas portas de um volume de controle. Quem já conhece o resultado
ganha em saber de onde ele vem; quem não conhece evita decorá-lo.

\subsection{A hipótese do equilíbrio local}\label{sub:v1-c3-equilibrio}

Há uma tensão evidente entre os dois capítulos. Propriedades como temperatura,
entropia e entalpia são definidas, na termodinâmica clássica, para sistemas em
\textbf{equilíbrio}: um sistema em equilíbrio não tem gradientes, não escoa e
não troca nada. Um compressor em operação é o oposto disso --- há gradiente de
pressão, de temperatura e de velocidade em toda parte, e nada está em
equilíbrio.

A reconciliação é a \textbf{hipótese do equilíbrio local}: embora o sistema
como um todo esteja fora de equilíbrio, cada partícula material pode ser
tratada, a cada instante, como um sistema em equilíbrio interno, com
temperatura, pressão e entropia bem definidas. Formalmente, supõe-se que a
mesma relação fundamental $s = s(u,v)$ da termodinâmica de equilíbrio vale
localmente, conforme a Equação~\eqref{eq:v1-c3-1}:
\begin{equation}\label{eq:v1-c3-1}
s(\mathbf{x},t) = s\big(u(\mathbf{x},t),\, v(\mathbf{x},t)\big).
\end{equation}

É essa hipótese --- e apenas ela --- que autoriza escrever $Ds/Dt$ na
Equação~\eqref{eq:v1-c2-26} e atribuir a $T$ o significado de temperatura
termodinâmica. Sua validade exige separação de escalas de tempo: o tempo de
relaxação molecular, isto é, o intervalo em que colisões restabelecem o
equilíbrio local, deve ser muito menor que o tempo característico de variação
macroscópica. Para gases em condições ordinárias o tempo de relaxação é da
ordem de $10^{-10}\,$s, contra tempos de processo de milissegundos ou mais
--- a margem é de várias ordens de grandeza. A hipótese falha em ondas de
choque muito intensas e em gases rarefeitos, exatamente as mesmas situações
que já haviam violado a hipótese do contínuo da
Seção~\ref{sub:v1-c2-hipotese}. Não é coincidência: as duas dependem do mesmo
número de Knudsen.

Aceito o equilíbrio local, resta saber quantas propriedades independentes
descrevem cada ponto. A resposta é o \textbf{postulado de estado}: o estado
intensivo de uma substância compressível simples --- aquela sem efeitos
significativos de superfície, campo elétrico ou magnético --- fica determinado
por \emph{duas} propriedades intensivas independentes. Toda tabela de vapor,
todo diagrama e toda equação de estado deste volume são consequência dessa
contagem. Vale reforçar o que ``independentes'' significa: pressão e
temperatura não servem como par na região de saturação, onde são dependentes
entre si.

Antes de tudo isso, porém, é preciso justificar que temperatura exista. Essa é
a função da \textbf{lei zero}: se dois corpos estão em equilíbrio térmico com
um terceiro, estão em equilíbrio térmico entre si. A relação ``estar em
equilíbrio térmico com'' é, portanto, uma relação de equivalência, e as classes
que ela define podem ser rotuladas por um número --- a temperatura empírica. A
lei zero garante a existência do rótulo; a segunda lei, mais adiante, dará a
ele uma escala absoluta e independente da substância usada para medi-lo.

\subsection{Primeira lei: a energia, do campo para o volume de controle}\label{sub:v1-c3-primeira}

Considere a Equação~\eqref{eq:v1-c2-25} integrada sobre um volume de controle
finito, com o teorema da divergência aplicado aos termos de fluxo. Os termos se
organizam naturalmente em quatro grupos: acúmulo no interior, transporte
convectivo pelas portas, condução de calor pela fronteira e trabalho das forças
de superfície.

O terceiro grupo merece atenção, porque é dele que nasce a entalpia. O trabalho
das forças de superfície nas portas de entrada e saída não é livre: para
empurrar massa para dentro contra a pressão $p$, é preciso realizar trabalho. A
potência associada a uma porta de área $A$ com velocidade normal $v_n$ é
$p\,A\,v_n$, e como $\dot{m} = \rho A v_n$, isso equivale a $\dot{m}\,p\,v$ ---
o produto da vazão mássica pelo volume específico e pela pressão. Esse termo
acompanha sempre e obrigatoriamente a energia interna transportada, de modo que
a combinação que atravessa uma porta nunca é $u$ sozinha, e sim $u + pv$.
Define-se, então, a \textbf{entalpia específica} da
Equação~\eqref{eq:v1-c3-2}:
\begin{equation}\label{eq:v1-c3-2}
h \equiv u + p\,v .
\end{equation}
A entalpia não é uma forma de energia --- é uma contabilidade. Ela existe para
que o trabalho de escoamento não precise ser escrito separadamente toda vez.

Separando ainda o trabalho restante em trabalho de eixo $\dot{W}$ --- o que a
máquina efetivamente entrega ou recebe --- e desprezando o trabalho viscoso na
fronteira, resulta a forma usual da primeira lei para volume de controle,
Equação~\eqref{eq:v1-c3-3}:
\begin{equation}\label{eq:v1-c3-3}
\frac{dE_{VC}}{dt} = \dot{Q} - \dot{W} +
\sum_{e}\dot{m}_e\!\left(h_e + \frac{v_e^2}{2} + gz_e\right) -
\sum_{s}\dot{m}_s\!\left(h_s + \frac{v_s^2}{2} + gz_s\right),
\end{equation}
onde os índices $e$ e $s$ referem-se a entradas e saídas. As premissas
embutidas, que raramente são listadas, são três: propriedades uniformes em cada
porta, ausência de trabalho viscoso na fronteira e --- do capítulo anterior ---
equilíbrio local, sem o qual $h_e$ e $h_s$ não teriam sentido.

O caso de um sistema fechado, sem portas, é o caso particular da
Equação~\eqref{eq:v1-c3-4}, integrada no tempo entre dois estados de
equilíbrio:
\begin{equation}\label{eq:v1-c3-4}
\Delta U = Q - W .
\end{equation}
Vale a convenção de sinais adotada em todo o volume: $Q$ é positivo quando
entra no sistema e $W$ é positivo quando o sistema realiza trabalho. Numa
máquina de compressão, portanto, $W$ é negativo --- ou, como é praxe em
engenharia de compressores e como fará o Volume II, troca-se o sinal e fala-se
em trabalho \emph{consumido}, positivo. A troca é legítima desde que
declarada; não declará-la é a origem mais frequente de erro de sinal em
balanços de planta.

\subsection{Segunda lei, temperatura absoluta e entropia}\label{sub:v1-c3-segunda}

A segunda lei nasceu de duas constatações experimentais equivalentes.
O enunciado de \textbf{Kelvin-Planck} afirma que nenhum dispositivo operando em
ciclo pode converter integralmente calor de uma única fonte em trabalho. O de
\textbf{Clausius} afirma que calor não passa espontaneamente de um corpo frio
para um quente. Demonstra-se que negar um implica negar o outro, de modo que os
dois enunciam a mesma lei.

Da segunda lei decorre o resultado de \textbf{Carnot}: todas as máquinas
térmicas reversíveis operando entre os mesmos dois reservatórios têm a mesma
eficiência, e nenhuma máquina real a supera. O ponto notável é que essa
eficiência não pode depender do fluido de trabalho --- se dependesse, acoplar
duas máquinas de fluidos diferentes produziria trabalho líquido sem fonte
quente. Como só resta a temperatura de que depender, define-se a \textbf{escala
termodinâmica absoluta} pela Equação~\eqref{eq:v1-c3-5}:
\begin{equation}\label{eq:v1-c3-5}
\left(\frac{Q_L}{Q_H}\right)_{\text{rev}} = \frac{T_L}{T_H}
\qquad\Longrightarrow\qquad
\eta_{\text{Carnot}} = 1 - \frac{T_L}{T_H} .
\end{equation}
É esta a razão de fundo pela qual as equações deste volume exigem temperatura
em kelvin: a escala absoluta não é uma convenção de unidade, é a única em que a
Equação~\eqref{eq:v1-c3-5} é verdadeira.

Aplicando o resultado de Carnot a um ciclo arbitrário decomposto em ciclos
elementares, obtém-se a \textbf{desigualdade de Clausius},
Equação~\eqref{eq:v1-c3-6}:
\begin{equation}\label{eq:v1-c3-6}
\oint \frac{\delta Q}{T} \le 0,
\end{equation}
com igualdade para processos reversíveis. Como a integral cíclica se anula no
caso reversível, o integrando é diferencial exata, e existe uma função de
estado cuja variação ele mede --- a \textbf{entropia},
Equação~\eqref{eq:v1-c3-7}:
\begin{equation}\label{eq:v1-c3-7}
dS \equiv \left(\frac{\delta Q}{T}\right)_{\text{rev}},
\qquad
\Delta S \ge \int \frac{\delta Q}{T} .
\end{equation}
A entropia é definida por um processo reversível, mas é função de estado: sua
variação entre dois estados independe do caminho, inclusive de caminhos
irreversíveis, que são apenas o modo mais comum de ir de um ao outro.

Para um volume de controle, integrar a Equação~\eqref{eq:v1-c2-26} do capítulo
anterior --- ou, equivalentemente, aplicar a Equação~\eqref{eq:v1-c3-7} com
contabilidade de portas --- fornece o balanço de entropia da
Equação~\eqref{eq:v1-c3-8}:
\begin{equation}\label{eq:v1-c3-8}
\frac{dS_{VC}}{dt} = \sum_{e}\dot{m}_e s_e - \sum_{s}\dot{m}_s s_s +
\sum_j \frac{\dot{Q}_j}{T_{b,j}} + \dot{S}_{gen},
\qquad \dot{S}_{gen} \ge 0 ,
\end{equation}
onde $T_{b,j}$ é a temperatura da fronteira onde ocorre cada troca $\dot Q_j$.
O termo $\dot{S}_{gen}$ é exatamente o $\dot{\mathcal{G}}_s$ da
Tabela~\ref{tab:v1-c2-2}, integrado sobre o volume: a mesma grandeza, escrita
em duas linguagens. Sua não negatividade é a única restrição de sinal de toda
a teoria, e é dela que sai tudo o que se afirma sobre eficiência máxima.

\subsection{Relações de propriedade}\label{sub:v1-c3-propriedades}

As leis de balanço, aqui como no Capítulo~\ref{sec:continuo}, não bastam: é
preciso ligar $u$, $h$ e $s$ às variáveis mensuráveis $p$, $v$ e $T$.

Combinando a primeira lei para um processo reversível de sistema fechado
($\delta q = T\,ds$ e $\delta w = p\,dv$) obtêm-se as \textbf{relações de
Gibbs}, Equações~\eqref{eq:v1-c3-9} e~\eqref{eq:v1-c3-10}, a segunda das quais
segue da primeira diferenciando a Equação~\eqref{eq:v1-c3-2}:
\begin{equation}\label{eq:v1-c3-9}
T\,ds = du + p\,dv,
\end{equation}
\begin{equation}\label{eq:v1-c3-10}
T\,ds = dh - v\,dp .
\end{equation}
Foram deduzidas supondo reversibilidade, mas relacionam apenas \emph{funções de
estado} --- e por isso valem para qualquer processo, reversível ou não. Este é
um ponto que costuma escapar: a hipótese entrou na dedução, não no resultado.

Definem-se os \textbf{calores específicos} pela Equação~\eqref{eq:v1-c3-11},
como as sensibilidades da energia interna e da entalpia à temperatura, cada uma
mantendo constante a variável natural da sua função:
\begin{equation}\label{eq:v1-c3-11}
c_v \equiv \left(\frac{\partial u}{\partial T}\right)_v,
\qquad
c_p \equiv \left(\frac{\partial h}{\partial T}\right)_p .
\end{equation}

\textbf{Gás ideal.} O modelo de gás ideal supõe moléculas de volume próprio
desprezível e sem forças de interação a distância, o que conduz à equação de
estado da Equação~\eqref{eq:v1-c3-12}:
\begin{equation}\label{eq:v1-c3-12}
p\,v = R\,T .
\end{equation}
Dessas mesmas premissas decorre que a energia interna depende apenas da
temperatura --- resultado de Joule ---, e daí, com a
Equação~\eqref{eq:v1-c3-2}, que a entalpia também. Diferenciando $h = u + RT$
resulta a relação de Mayer, Equação~\eqref{eq:v1-c3-13}:
\begin{equation}\label{eq:v1-c3-13}
c_p - c_v = R,
\qquad
\gamma \equiv \frac{c_p}{c_v} .
\end{equation}
Substituindo a Equação~\eqref{eq:v1-c3-12} na
Equação~\eqref{eq:v1-c3-10} e integrando com $c_p$ constante, obtém-se a
variação de entropia da Equação~\eqref{eq:v1-c3-14}, que é a forma efetivamente
usada em máquinas de compressão:
\begin{equation}\label{eq:v1-c3-14}
s_2 - s_1 = c_p \ln\!\frac{T_2}{T_1} - R \ln\!\frac{p_2}{p_1} .
\end{equation}

\textbf{Quando o gás ideal não serve.} O desvio se mede pelo \textbf{fator de
compressibilidade} da Equação~\eqref{eq:v1-c3-15}:
\begin{equation}\label{eq:v1-c3-15}
Z_c = \frac{p\,v}{R\,T} .
\end{equation}
Para ar em temperatura ambiente, $Z_c$ difere de 1 por menos de $1\,\%$ até
algumas dezenas de bar, e o modelo ideal é adequado a um trem de ar comprimido
industrial. O mesmo não vale perto da saturação nem em gases de processo
pesados a alta pressão, casos em que uma equação de estado real se torna
obrigatória. A verificação correta é calcular $Z_c$, não presumir pela
familiaridade com o fluido.

\subsection{Processos ideais e o índice politrópico}\label{sub:v1-c3-processos}

Um processo \textbf{isentrópico} é adiabático e reversível. Impondo
$s_2 = s_1$ na Equação~\eqref{eq:v1-c3-14} e usando a
Equação~\eqref{eq:v1-c3-13}, obtém-se diretamente a relação da
Equação~\eqref{eq:v1-c3-16}:
\begin{equation}\label{eq:v1-c3-16}
\frac{T_2}{T_1} = \left(\frac{p_2}{p_1}\right)^{\frac{\gamma-1}{\gamma}},
\qquad
p\,v^{\gamma} = \text{const.}
\end{equation}
Note que a segunda forma não foi postulada: ela é consequência da primeira, da
equação de estado e da hipótese de calores específicos constantes.

Compressão real não é isentrópica --- há troca de calor com a carcaça e há
irreversibilidade. A generalização usual substitui $\gamma$ por um expoente
ajustado, o \textbf{índice politrópico} $n$ da
Equação~\eqref{eq:v1-c3-17}:
\begin{equation}\label{eq:v1-c3-17}
p\,v^{n} = \text{const.},
\qquad
\frac{T_2}{T_1} = \left(\frac{p_2}{p_1}\right)^{\frac{n-1}{n}} .
\end{equation}
Diferente de $\gamma$, que é propriedade do gás, $n$ é propriedade do
\emph{processo}: descreve o quanto o caminho real se afasta do isentrópico.
Refrigeração intensa durante a compressão puxa $n$ abaixo de $\gamma$;
irreversibilidade e aquecimento por atrito o empurram acima. Justamente por
isso $n$ é observável: invertendo a Equação~\eqref{eq:v1-c3-17} para as duas
medições que qualquer estágio instrumentado fornece --- razão de pressão e
razão de temperatura ---, obtém-se a Equação~\eqref{eq:v1-c3-18}:
\begin{equation}\label{eq:v1-c3-18}
n = \frac{1}{1 - \dfrac{\ln(T_2/T_1)}{\ln(p_2/p_1)}} .
\end{equation}
É esta equação que transforma $n$ de parâmetro de ajuste em indicador de
condição: um deslocamento lento de $n$ ao longo de semanas, a razão de pressão
constante, é assinatura de degradação --- e é assim que o Volume II o utiliza.

A Tabela~\ref{tab:v1-c3-1} reúne os processos elementares. Ela traz duas
colunas de trabalho, e a distinção entre elas é a armadilha mais cara deste
capítulo, tratada a seguir.

\begin{table}[H]
\centering
\caption[Processos ideais para gás ideal]{Processos termodinâmicos ideais para
gás ideal, com o trabalho de sistema fechado ($\int p\,dv$) e o de escoamento
em regime permanente ($\int v\,dp$).}
\label{tab:v1-c3-1}
\small
\resizebox{\textwidth}{!}{%
\begin{tabular}{@{}lllll@{}}
\toprule
\textbf{Processo} & \textbf{Relação} & $\Delta u$ & $w_{\text{fechado}} = \int p\,dv$ & $w_{\text{escoamento}} = \int v\,dp$ \\
\midrule
Isocórico ($v$ const.) & $p/T$ const. & $c_v\Delta T$ & $0$ & $v\,\Delta p$ \\
Isobárico ($p$ const.) & $v/T$ const. & $c_v\Delta T$ & $p\,\Delta v$ & $0$ \\
Isotérmico ($T$ const.) & $pv$ const. & $0$ & $RT\ln(v_2/v_1)$ & $RT\ln(p_1/p_2)$ \\
Isentrópico & $pv^{\gamma}$ const. & $c_v\Delta T$ & $\dfrac{p_1v_1 - p_2v_2}{\gamma-1}$ & $\dfrac{\gamma}{\gamma-1}\,(p_1v_1 - p_2v_2)$ \\
Politrópico & $pv^{n}$ const. & $c_v\Delta T$ & $\dfrac{p_1v_1 - p_2v_2}{n-1}$ & $\dfrac{n}{n-1}\,(p_1v_1 - p_2v_2)$ \\
\bottomrule
\end{tabular}}
\end{table}

\subsection{Trabalho de compressão}\label{sub:v1-c3-trabalho}

Esta seção produz a equação que o Volume II aplica a cada estágio, e convém
segui-la passo a passo, porque o resultado difere daquele que a intuição de
sistema fechado sugere.

Para escoamento em regime permanente, desprezando variações de energia cinética
e potencial, a primeira lei por unidade de massa dá $\delta q + \delta w = dh$.
Num processo reversível, $\delta q = T\,ds$. Substituindo ambos na
Equação~\eqref{eq:v1-c3-10} resulta $\delta w = v\,dp$, isto é, o trabalho
específico da Equação~\eqref{eq:v1-c3-19}:
\begin{equation}\label{eq:v1-c3-19}
w = \int_{p_1}^{p_2} v \; dp .
\end{equation}

Aqui está a armadilha anunciada. O trabalho de um sistema \emph{fechado} é
$\int p\,dv$; o de uma máquina de escoamento é $\int v\,dp$. Não são iguais, e
para um processo politrópico diferem por um fator $n$, conforme a
Equação~\eqref{eq:v1-c3-20}:
\begin{equation}\label{eq:v1-c3-20}
\frac{\int v\,dp}{\int p\,dv} = n .
\end{equation}
Usar a fórmula de sistema fechado para dimensionar um compressor subestima a
potência em cerca de $35\,\%$ para $n = 1{,}35$. O compressor de parafuso, por
mais que comprima gás dentro de uma câmara aparentemente fechada, é uma máquina
de escoamento: admite, comprime e descarrega continuamente, e as etapas de
admissão e descarga --- em que a fronteira se move sob pressão constante ---
são justamente o que a Equação~\eqref{eq:v1-c3-19} contabiliza e a de sistema
fechado ignora.

Substituindo a Equação~\eqref{eq:v1-c3-17} na Equação~\eqref{eq:v1-c3-19} e
integrando, com $p_1v_1 = RT_1$, obtém-se o \textbf{trabalho politrópico
específico} da Equação~\eqref{eq:v1-c3-21}:
\begin{equation}\label{eq:v1-c3-21}
w_{\text{pol}} = \frac{n}{n-1}\,R\,T_1
\left[ \left(\frac{p_2}{p_1}\right)^{\frac{n-1}{n}} - 1 \right] .
\end{equation}
Fazendo $n = \gamma$ recupera-se o trabalho isentrópico, que é a referência
mínima para compressão adiabática.

A comparação entre o real e o ideal se faz por dois rendimentos distintos, e
confundi-los é erro comum. O \textbf{rendimento isentrópico} da
Equação~\eqref{eq:v1-c3-22} compara o trabalho ideal adiabático ao real; o
\textbf{rendimento politrópico} da Equação~\eqref{eq:v1-c3-23} compara os
mesmos trabalhos de forma incremental, ao longo do caminho:
\begin{equation}\label{eq:v1-c3-22}
\eta_{is} = \frac{w_{is}}{w_{\text{real}}}
= \frac{\Pi^{\frac{\gamma-1}{\gamma}} - 1}{\dfrac{T_2}{T_1} - 1},
\qquad \Pi \equiv \frac{p_2}{p_1},
\end{equation}
\begin{equation}\label{eq:v1-c3-23}
\eta_p = \frac{\gamma-1}{\gamma}\cdot\frac{n}{n-1} .
\end{equation}
Eliminando $n$ entre as duas, chega-se à
Equação~\eqref{eq:v1-c3-24}, que revela o ponto prático:
\begin{equation}\label{eq:v1-c3-24}
\eta_{is} = \frac{\Pi^{\frac{\gamma-1}{\gamma}} - 1}
{\Pi^{\frac{\gamma-1}{\gamma\,\eta_p}} - 1} .
\end{equation}
$\eta_p$ é aproximadamente constante ao longo do mapa da máquina, ao passo que
$\eta_{is}$ \emph{cai com a razão de pressão mesmo sem nenhuma degradação}.
Comparar $\eta_{is}$ medido em dois pontos de operação diferentes, ou contra
uma garantia de fábrica cotada em outra razão de pressão, produz um diagnóstico
falso de perda de desempenho. É $\eta_p$ que se compara.

Por fim, a Equação~\eqref{eq:v1-c3-19} explica por que se comprime em estágios.
Somando o trabalho de dois estágios com resfriamento intermediário perfeito
(isto é, com o gás retornando a $T_1$ antes do segundo estágio) e derivando em
relação à pressão intermediária $p_i$, a condição de mínimo fornece a
Equação~\eqref{eq:v1-c3-25}:
\begin{equation}\label{eq:v1-c3-25}
\frac{\partial w_{\text{total}}}{\partial p_i} = 0
\qquad\Longrightarrow\qquad
p_i^{*} = \sqrt{p_1\,p_2},
\end{equation}
ou seja, razões de compressão iguais nos dois estágios. O ganho vem de reduzir
$v$ antes da segunda compressão: como o trabalho é $\int v\,dp$, resfriar entre
estágios comprime um gás mais denso e barateia cada incremento de pressão. Um
intercooler degradado, portanto, não apenas eleva a temperatura de descarga ---
ele desloca o ponto ótimo de repartição, e é essa dupla consequência que o
Volume II monitora.

\subsection{Exergia e destruição de disponibilidade}\label{sub:v1-c3-exergia}

A primeira lei contabiliza quantidade de energia; a segunda diz que nem toda
ela é aproveitável. A grandeza que mede a parte aproveitável é a
\textbf{exergia}: o máximo trabalho útil extraível de uma corrente até que ela
atinja o equilíbrio com o ambiente de referência, definido por $T_0$ e $p_0$.

A construção é direta e vale acompanhá-la, porque explica a origem do termo
$T_0 s$. Multiplica-se o balanço de entropia,
Equação~\eqref{eq:v1-c3-8}, por $T_0$ e subtrai-se do balanço de energia,
Equação~\eqref{eq:v1-c3-3}. O que sobra em cada porta é a combinação
$h - T_0 s$, e descontando o valor dessa mesma combinação no estado de
referência obtém-se a \textbf{exergia específica de fluxo} da
Equação~\eqref{eq:v1-c3-26}:
\begin{equation}\label{eq:v1-c3-26}
\psi = (h - h_0) - T_0 (s - s_0) + \frac{v^2}{2} + g z .
\end{equation}
O termo de calor, pelo mesmo procedimento, aparece multiplicado pelo fator de
Carnot, e o balanço de exergia em regime permanente fica como na
Equação~\eqref{eq:v1-c3-27}:
\begin{equation}\label{eq:v1-c3-27}
\sum_e \dot{m}_e \psi_e +
\sum_j \dot{Q}_j\!\left(1 - \frac{T_0}{T_j}\right) =
\sum_s \dot{m}_s \psi_s + \dot{W} + \dot{E}_d .
\end{equation}
Diferente da energia, a exergia \textbf{não se conserva}: o termo
$\dot{E}_d \ge 0$ é uma perda real, não um resíduo de contabilidade.

Da própria construção sai o \textbf{teorema de Gouy-Stodola},
Equação~\eqref{eq:v1-c3-28}, que quantifica essa perda:
\begin{equation}\label{eq:v1-c3-28}
\dot{E}_d = T_0 \, \dot{S}_{gen} .
\end{equation}
Sua utilidade prática é converter a geração de entropia --- grandeza abstrata,
de unidade pouco intuitiva --- em watts perdidos, comparáveis diretamente com a
potência de eixo. E permite ordenar as perdas de uma planta por importância
real: uma troca térmica com grande diferença de temperatura pode gerar mais
$\dot{E}_d$ que um vazamento visível, sem que nada pareça errado nas leituras.

O rendimento correspondente, definido na Equação~\eqref{eq:v1-c3-29}, mede
quanto da exergia fornecida sobrevive no produto:
\begin{equation}\label{eq:v1-c3-29}
\eta_{ex} = \frac{\text{exergia recuperada}}{\text{exergia fornecida}}
= 1 - \frac{\dot{E}_d}{\dot{E}_{x,\text{forn}}} .
\end{equation}
A Tabela~\ref{tab:v1-c3-2} resume as premissas de cada resultado central deste
capítulo, e serve como lista de verificação antes de aplicar qualquer um deles
fora do contexto em que foi deduzido.

\begin{table}[H]
\centering
\caption[Premissas dos resultados do capítulo]{Resultados centrais do capítulo
e as premissas de que dependem.}
\label{tab:v1-c3-2}
\small
\resizebox{\textwidth}{!}{%
\begin{tabular}{@{}lll@{}}
\toprule
\textbf{Resultado} & \textbf{Equação} & \textbf{Premissas} \\
\midrule
Formalismo de estados & \eqref{eq:v1-c3-1} & Equilíbrio local; $\mathrm{Kn}\ll1$ \\
Primeira lei, volume de controle & \eqref{eq:v1-c3-3} & Portas uniformes; sem trabalho viscoso na fronteira \\
Relações de Gibbs & \eqref{eq:v1-c3-9}, \eqref{eq:v1-c3-10} & Substância compressível simples; valem sempre \\
Relação de Mayer, $\Delta s$ & \eqref{eq:v1-c3-13}, \eqref{eq:v1-c3-14} & Gás ideal; $c_p$ constante na faixa \\
Relação isentrópica & \eqref{eq:v1-c3-16} & Gás ideal; adiabático e reversível \\
Trabalho de escoamento & \eqref{eq:v1-c3-19} & Regime permanente; reversível; $\Delta$KE e $\Delta$PE desprezíveis \\
Trabalho politrópico & \eqref{eq:v1-c3-21} & As anteriores, mais $n$ constante no estágio \\
Pressão intermediária ótima & \eqref{eq:v1-c3-25} & Resfriamento intermediário perfeito; mesmo $n$ nos estágios \\
Gouy-Stodola & \eqref{eq:v1-c3-28} & Ambiente de referência fixo ($T_0$, $p_0$) \\
\bottomrule
\end{tabular}}
\end{table}


%=============================================================================
% CAPÍTULO 4 — MECÂNICA DOS FLUIDOS
%=============================================================================
\section{Mecânica dos Fluidos}\label{sec:fluidos}

\subsection{Escoamento de fluidos}

A mecânica dos fluidos estuda o movimento de fluidos e suas interações com
fronteiras. Para um fluido newtoniano, as equações de Navier-Stokes são
obtidas substituindo a lei constitutiva na conservação da quantidade de
movimento:
\begin{equation}\label{eq:v1-c4-1}
\rho\frac{D\mathbf{v}}{Dt} = -\nabla p + \mu\nabla^2\mathbf{v} +
\left(\lambda + \mu\right)\nabla\left(\nabla\cdot\mathbf{v}\right) + \rho\mathbf{g}.
\end{equation}

Para fluidos incompressíveis e com viscosidade constante, essa equação
simplifica-se para
\begin{equation}\label{eq:v1-c4-2}
\rho\frac{D\mathbf{v}}{Dt} = -\nabla p + \mu\nabla^2\mathbf{v} + \rho\mathbf{g}.
\end{equation}

\subsection{Equação de Bernoulli}

Para escoamento permanente, incompressível, inviscido e ao longo de uma linha
de corrente, a equação de Bernoulli relaciona pressão, velocidade e altura:
\begin{equation}\label{eq:v1-c4-3}
\frac{P}{\rho g} + \frac{v^2}{2g} + z = \text{constante}.
\end{equation}

Em forma de pressão,
\begin{equation}\label{eq:v1-c4-4}
P + \frac{1}{2}\rho v^2 + \rho g z = \text{constante}.
\end{equation}

\subsection{Escoamento em dutos}

Para dutos circulares em regime permanente, a perda de carga distribuída é
dada pela equação de Darcy-Weisbach:
\begin{equation}\label{eq:v1-c4-5}
\Delta P = f_D \, \frac{L}{D} \, \frac{\rho v^2}{2},
\end{equation}
onde $f_D$ é o fator de atrito de Darcy, $L$ é o comprimento do duto e $D$ é o
diâmetro interno.

O número de Reynolds define o regime do escoamento:
\begin{equation}\label{eq:v1-c4-6}
Re_D = \frac{\rho v D}{\mu} = \frac{vD}{\nu_k}.
\end{equation}
Para $Re_D < 2300$ o escoamento é laminar; para $Re_D > 4000$, turbulento.

\subsection{Números adimensionais}

A Tabela~\ref{tab:v1-c4-1} lista os principais grupos adimensionais usados na
análise de escoamentos térmicos e rotodinâmicos.

\begin{table}[H]
\centering
\caption{Números adimensionais recorrentes em máquinas rotodinâmicas.}
\label{tab:v1-c4-1}
\begin{tabular}{@{}llll@{}}
\toprule
\textbf{Número} & \textbf{Definição} & \textbf{Razão física} & \textbf{Uso} \\
\midrule
Reynolds & $Re = \rho v L / \mu$ & Inércia/viscosidade & Regime do escoamento \\
Mach & $Ma = v/a$ & Velocidade/velocidade do som & Compressibilidade \\
Prandtl & $Pr = \mu c_p / k$ & Difusão viscosa/térmica & Escoamentos com calor \\
Nusselt & $Nu = hL/k$ & Convecção/condução & Transferência de calor \\
Euler & $Eu = \Delta P / (\rho v^2)$ & Pressão/inércia & Máquinas hidráulicas \\
\bottomrule
\end{tabular}
\end{table}

\subsection{Camada limite}

A camada limite é a região adjacente a uma superfície sólida onde os efeitos
viscosos são importantes. Sua espessura $\delta$ cresce ao longo da placa
de acordo com
\begin{equation}\label{eq:v1-c4-7}
\delta(x) \approx \frac{5x}{\sqrt{Re_x}},
\end{equation}
para uma placa plana laminar. O conhecimento da camada limite é essencial
para estimar arrasto, perdas por atrito e transferência de calor em
superfícies de máquinas rotativas.


%=============================================================================
% CAPÍTULO 5 — TRANSFERÊNCIA DE CALOR
%=============================================================================
\section{Transferência de Calor}\label{sec:calor}

\subsection{Modos de transferência de calor}

A transferência de calor ocorre por três modos: condução, convecção e
radiação. Em máquinas rotodinâmicas, os três estão presentes, embora a
convecção forçada domine em trocadores de calor e resfriadores.

\subsection{Condução}

A condução em meios isotrópicos é descrita pela lei de Fourier:
\begin{equation}\label{eq:v1-c5-1}
\mathbf{q} = -k \, \nabla T,
\end{equation}
onde $k$ é a condutividade térmica do material.

Para uma parede plana unidimensional em regime permanente, o fluxo de calor é
\begin{equation}\label{eq:v1-c5-2}
q'' = \frac{k}{L}(T_1 - T_2).
\end{equation}

Em coordenadas cilíndricas (tubos), a resistência térmica por unidade de
comprimento é
\begin{equation}\label{eq:v1-c5-3}
R'_{cond} = \frac{\ln(r_2/r_1)}{2\pi k}.
\end{equation}

\subsection{Convecção}

A convecção obedece à lei do resfriamento de Newton:
\begin{equation}\label{eq:v1-c5-4}
q'' = h\,(T_s - T_\infty),
\end{equation}
onde $h$ é o coeficiente de transferência de calor por convecção e $T_s$ e
$T_\infty$ são as temperaturas da superfície e do fluido longe dela.

O coeficiente $h$ é geralmente obtido por correlações adimensionais, como
\begin{equation}\label{eq:v1-c5-5}
Nu = C \, Re^m \, Pr^n,
\end{equation}
com constantes $C$, $m$ e $n$ que dependem da geometria e do regime.

\subsection{Radiação}

A radiação térmica entre duas superfícies cinzas grandes e próximas é
\begin{equation}\label{eq:v1-c5-6}
q'' = \varepsilon_r \, \sigma \, (T_1^4 - T_2^4),
\end{equation}
onde $\varepsilon_r$ é a emissividade efetiva, $\sigma$ é a constante de
Stefan-Boltzmann e as temperaturas devem estar em Kelvin.

\subsection{Trocadores de calor}

A taxa de transferência de calor em um trocador pode ser calculada por
\begin{equation}\label{eq:v1-c5-7}
\dot{Q} = U \, A \, \Delta T_{ml},
\end{equation}
onde $U$ é o coeficiente global de transferência de calor, $A$ é a área de
troca e $\Delta T_{ml}$ é a diferença média logarítmica de temperatura (LMTD):
\begin{equation}\label{eq:v1-c5-8}
\Delta T_{ml} = \frac{\Delta T_1 - \Delta T_2}{\ln(\Delta T_1 / \Delta T_2)}.
\end{equation}

O método da efetividade-NTU fornece uma alternativa quando as temperaturas de
saída não são conhecidas:
\begin{equation}\label{eq:v1-c5-9}
\varepsilon = \frac{\dot{Q}}{\dot{Q}_{max}}, \qquad
NTU = \frac{UA}{C_{min}}.
\end{equation}
A relação entre $\varepsilon$, $NTU$ e $C_r = C_{min}/C_{max}$ depende da
configuração do trocador (paralelo, contracorrente, casco-tubo etc.).

\begin{table}[H]
\centering
\caption{Resistências térmicas em série para uma parede composta.}
\label{tab:v1-c5-1}
\begin{tabular}{@{}lll@{}}
\toprule
\textbf{Modo} & \textbf{Resistência} & \textbf{Unidade} \\
\midrule
Condução plana & $R = L/(kA)$ & K\,W$^{-1}$ \\
Condução cilíndrica & $R = \ln(r_2/r_1)/(2\pi k L)$ & K\,W$^{-1}$ \\
Convecção & $R = 1/(hA)$ & K\,W$^{-1}$ \\
Global & $1/(UA) = \sum R_i$ & K\,W$^{-1}$ \\
\bottomrule
\end{tabular}
\end{table}


%=============================================================================
% CAPÍTULO 6 — SISTEMAS DE TRANSFERÊNCIA DE MASSA E ENERGIA
%=============================================================================
\section{Sistemas de Transferência de Massa e Energia}\label{sec:sistemas}

\subsection{Volumes de controle}

A análise de sistemas de transferência de massa e energia baseia-se na
formulção de volumes de controle. Um volume de controle é uma região fixa no
espaço através da qual massa, energia e momento podem fluir. As equações de
balanço para um volume de controle em regime permanente são:

\textbf{Massa:}
\begin{equation}\label{eq:v1-c6-1}
\sum_e \dot{m}_e = \sum_s \dot{m}_s.
\end{equation}

\textbf{Energia:}
\begin{equation}\label{eq:v1-c6-2}
\dot{Q} - \dot{W} = \sum_s \dot{m}_s\left(h_s + \frac{v_s^2}{2} + gz_s\right) - \sum_e \dot{m}_e\left(h_e + \frac{v_e^2}{2} + gz_e\right).
\end{equation}

\textbf{Quantidade de movimento:}
\begin{equation}\label{eq:v1-c6-3}
\sum \mathbf{F} = \sum_s \dot{m}_s \mathbf{v}_s - \sum_e \dot{m}_e \mathbf{v}_e.
\end{equation}

\subsection{Regimes permanente e transiente}

Em regime permanente, as propriedades em qualquer ponto do volume de controle
não variam com o tempo. Em regime transiente, o acúmulo deve ser considerado:
\begin{equation}\label{eq:v1-c6-4}
\frac{dM_{VC}}{dt} = \sum_e \dot{m}_e - \sum_s \dot{m}_s.
\end{equation}

Aplicado à energia:
\begin{equation}\label{eq:v1-c6-5}
\frac{dE_{VC}}{dt} = \dot{Q} - \dot{W} + \sum_e \dot{m}_e\theta_e - \sum_s \dot{m}_s\theta_s,
\end{equation}
onde $\theta = h + v^2/2 + gz$ representa a entalpia de estagnação específica.

\subsection{Sistemas de armazenamento e distribuição}

Vasos de pressão, tanques e redes de tubulações são exemplos de sistemas de
armazenamento e distribuição. Para um vaso de volume $V$ contendo gás ideal,
a equação de estado fornece
\begin{equation}\label{eq:v1-c6-6}
P\,V = m\,R\,T.
\end{equation}

A derivada temporal da pressão em um vaso com entrada e saída é obtida
combinando a conservação de massa com a equação de estado, resultado usado no
cálculo de autonomia de sistemas de ar comprimido.


%=============================================================================
% CAPÍTULO 7 — FUNDAMENTOS DE ELETRICIDADE
%=============================================================================
\section{Fundamentos de Eletricidade}\label{sec:eletrica}

Os capítulos anteriores seguiram sempre a mesma construção: uma lei de balanço
sobre uma fronteira, uma hipótese que autoriza tratar matéria discreta como
campo, e relações constitutivas que precisam ser medidas. Este capítulo faz a
mesma construção para o eletromagnetismo, e a faz por uma razão prática: a
teoria de circuitos --- lei de Ohm, impedância, fasores --- não é uma teoria
primitiva. Ela é um \emph{limite} das equações de campo, exatamente como a
mecânica do contínuo é um limite da mecânica de partículas, e tem condições de
validade tão verificáveis quanto o número de Knudsen da
Seção~\ref{sub:v1-c2-hipotese}. Começar pela lei de Ohm é começar depois do
ponto em que as respostas importantes já foram dadas.

Isso não é preciosismo acadêmico neste contexto específico. Os três mecanismos
de falha elétrica que dominam a mortalidade de um motor acionado por inversor
--- sobretensão nos terminais por reflexão de onda, aquecimento localizado por
desbalanço entre fases e erosão de rolamento por corrente de eixo --- são todos
fenômenos que a teoria de circuitos concentrados \textbf{não} descreve, porque
todos vivem justamente onde a hipótese de parâmetros concentrados falha. Um
volume de fundamentos que parasse na lei de Ohm entregaria ao Volume II um
motor que só sabe adoecer de excesso de corrente.

A ordem do capítulo repete a do Capítulo~\ref{sec:continuo}: primeiro a base
de campo e a lei de conservação de que tudo decorre
(Seção~\ref{sub:v1-c7-bases}); depois a hipótese que autoriza substituir campos
por circuitos, com seu critério quantitativo e seus três modos de falha
(Seção~\ref{sub:v1-c7-concentrados}); depois as relações constitutivas do meio,
que são onde o material entra (Seção~\ref{sub:v1-c7-constitutivas}); depois o
formalismo de fasores e a contabilidade de potência
(Seções~\ref{sub:v1-c7-fasores} e~\ref{sub:v1-c7-potencia}); depois o sistema
trifásico e sua decomposição em sequências
(Seção~\ref{sub:v1-c7-trifasico}); e só então a conversão eletromecânica
(Seção~\ref{sub:v1-c7-conversao}), a máquina de indução
(Seção~\ref{sub:v1-c7-inducao}), o acionamento por inversor
(Seção~\ref{sub:v1-c7-inversor}) e a contabilidade de perdas
(Seção~\ref{sub:v1-c7-perdas}).

\textbf{Nota sobre colisões de símbolos.} A notação consagrada do
eletromagnetismo colide com a dos capítulos anteriores em cinco pontos, e as
colisões são resolvidas aqui de uma vez: a condutividade elétrica é
$\sigma_e$ (o tensor de tensão de Cauchy é $\boldsymbol{\sigma}$, Cap.~2); a
permissividade é $\varepsilon_e$ (a deformação é $\boldsymbol{\varepsilon}$,
Cap.~2, e a efetividade de trocador é $\varepsilon$, Cap.~5); a permeabilidade
magnética é $\mu_m$ (a viscosidade dinâmica é $\mu$, Caps.~2 e~4); a potência
ativa é $P_{el}$ e a reativa $Q_{el}$ (a pressão é $P$, Caps.~2--6, e a taxa de
calor é $\dot{Q}$, Caps.~3--6); e o campo elétrico é $\mathbf{E}$, em negrito
por ser vetor, distinto do módulo de Young $E$ (Caps.~2 e~8) e do tensor de
Green-Lagrange $\mathbf{E}_L$ (Cap.~2). Duas colisões \textbf{não} são
resolvidas, deliberadamente, porque a notação alternativa isolaria este volume
do Volume II e da literatura de máquinas elétricas: o escorregamento é $s$,
como a entropia específica do Capítulo~\ref{sec:termodinamica}, e o número de
polos é $p$, como a pressão em minúscula. Nos dois casos o contexto é
inequívoco --- não há entropia em uma equação de escorregamento nem pressão em
uma de velocidade síncrona --- e a coincidência é sinalizada onde aparece.

\subsection{Bases físicas: carga, campo e força}\label{sub:v1-c7-bases}

A grandeza primitiva do eletromagnetismo é a \textbf{carga elétrica}. Ela tem
duas propriedades que a mecânica do contínuo não precisou postular para a
massa: é \emph{quantizada}, em múltiplos inteiros da carga elementar
$e = 1{,}602\times10^{-19}\,$C, e existe em dois sinais. A primeira propriedade
é irrelevante na escala de engenharia --- um ampère corresponde a cerca de
$6\times10^{18}$ cargas elementares por segundo, e a granularidade se dilui
exatamente como a granularidade molecular se diluiu na
Seção~\ref{sub:v1-c2-hipotese}. A segunda é decisiva: por haver dois sinais, a
matéria comum é neutra em escala macroscópica, e é essa neutralidade que
permite que correntes enormes atravessem um condutor sem que ele acumule carga.

A lei de conservação correspondente é o caso $\mathcal{B} =$ carga da
Equação~\eqref{eq:v1-c1-1}, prometido na introdução deste volume e cobrado
agora. Escrita na forma local, com $\rho_q$ a densidade volumétrica de carga e
$\mathbf{J}$ a densidade de corrente, ela é a Equação~\eqref{eq:v1-c7-1}:
\begin{equation}\label{eq:v1-c7-1}
\frac{\partial \rho_q}{\partial t} + \nabla\cdot\mathbf{J} = 0 .
\end{equation}
A Equação~\eqref{eq:v1-c7-1} é formalmente idêntica à continuidade de massa,
Equação~\eqref{eq:v1-c2-22}, com $\rho_q$ no lugar de $\rho$ e $\mathbf{J}$ no
lugar de $\rho\mathbf{v}$ --- e, como aquela, tem geração nula. Carga não é
criada nem destruída; apenas se desloca. (O símbolo $\mathbf{J}$ para densidade
de corrente não deve ser confundido com $J$, o jacobiano da deformação do
Capítulo~\ref{sec:continuo}, nem com $J_p$, o momento polar de inércia do
Capítulo~\ref{sec:mecanica}.)

\textbf{As equações de campo.} O que a mecânica do contínuo obtém por média
sobre partículas, o eletromagnetismo obtém por postulado experimental
consolidado: as quatro equações de \textbf{Maxwell}, reunidas na
Equação~\eqref{eq:v1-c7-2}, que relacionam o campo elétrico $\mathbf{E}$, a
indução magnética $\mathbf{B}$, o deslocamento elétrico $\mathbf{D}$ e o campo
magnético $\mathbf{H}$ às suas fontes:
\begin{equation}\label{eq:v1-c7-2}
\begin{aligned}
\nabla\cdot\mathbf{D} &= \rho_q, &\qquad
\nabla\times\mathbf{E} &= -\frac{\partial\mathbf{B}}{\partial t}, \\[4pt]
\nabla\cdot\mathbf{B} &= 0, &\qquad
\nabla\times\mathbf{H} &= \mathbf{J} + \frac{\partial\mathbf{D}}{\partial t} .
\end{aligned}
\end{equation}

Cada uma tem uma leitura de engenharia direta, e vale registrá-las porque são
elas que reaparecem disfarçadas em todo o resto do capítulo. A lei de Gauss
($\nabla\cdot\mathbf{D} = \rho_q$) diz que carga é fonte de campo elétrico ---
é o que torna possível o acoplamento capacitivo parasita da
Seção~\ref{sub:v1-c7-concentrados}. A ausência de monopolos
($\nabla\cdot\mathbf{B} = 0$) diz que linhas de campo magnético se fecham sobre
si mesmas --- é o que obriga o fluxo a percorrer um circuito magnético fechado,
estator--entreferro--rotor--entreferro--estator, e o que faz do entreferro o
elemento dominante da relutância. A lei de \textbf{Faraday}
($\nabla\times\mathbf{E} = -\partial\mathbf{B}/\partial t$) diz que campo
magnético variável gera campo elétrico --- é o princípio do transformador, do
gerador e, neste capítulo, da própria indução de corrente no rotor gaiola. A lei
de \textbf{Ampère-Maxwell} diz que corrente e campo elétrico variável geram
campo magnético --- é o princípio do eletroímã e do enrolamento de estator.

Duas observações estruturais. A primeira: a Equação~\eqref{eq:v1-c7-1}
\emph{não} é independente da Equação~\eqref{eq:v1-c7-2}. Tomando o divergente
da lei de Ampère-Maxwell e usando $\nabla\cdot(\nabla\times\mathbf{H})=0$
recupera-se exatamente a conservação de carga. É a mesma situação encontrada no
Capítulo~\ref{sec:continuo}, em que o balanço de momento angular não produziu
equação nova, apenas a simetria da Equação~\eqref{eq:v1-c2-15}: uma lei que
parecia adicional estava contida na estrutura. A segunda: o termo
$\partial\mathbf{D}/\partial t$, a \emph{corrente de deslocamento}, é o único
acréscimo de Maxwell às leis experimentais que o precederam, e é ele que torna
o conjunto consistente com a conservação de carga --- e que, ao mesmo tempo,
faz as equações admitirem solução ondulatória com velocidade
$c = 1/\sqrt{\varepsilon_e\mu_m}$. Guardar essa velocidade é essencial: é dela
que sai o critério de validade da próxima seção.

\textbf{A ponte com a mecânica.} O elo entre este capítulo e todos os
anteriores é a força de \textbf{Lorentz}, Equação~\eqref{eq:v1-c7-3}, escrita
aqui na forma volumétrica que se acopla diretamente ao termo de força de corpo
$\rho\mathbf{b}$ da equação de Cauchy, Equação~\eqref{eq:v1-c2-24}:
\begin{equation}\label{eq:v1-c7-3}
\mathbf{f}_{L} = \rho_q\mathbf{E} + \mathbf{J}\times\mathbf{B} .
\end{equation}
A Seção~\ref{sub:v1-c2-bases} listou, entre os limites de validade da base
mecânica, a ``ausência de campos eletromagnéticos acoplados ao movimento'', e
remeteu a este capítulo. A Equação~\eqref{eq:v1-c7-3} é o cumprimento dessa
promessa: nas máquinas tratadas nesta coleção o acoplamento existe, mas é
unidirecional e localizado. O campo produz força no rotor, essa força integra
em torque, e o torque entra na dinâmica mecânica como condição de contorno; o
movimento do rotor, por sua vez, realimenta o campo apenas através da tensão
induzida --- não através de deformação do meio. É essa unidirecionalidade
prática que autoriza tratar motor e máquina acionada como dois modelos
acoplados por uma única variável escalar, o torque de eixo, em vez de um único
problema magneto-mecânico.

\textbf{O balanço de energia do campo.} Falta a contabilidade energética, que é
o que fecha o círculo com o Capítulo~\ref{sec:termodinamica}. Manipulando a
Equação~\eqref{eq:v1-c7-2} obtém-se o teorema de \textbf{Poynting},
Equação~\eqref{eq:v1-c7-4}:
\begin{equation}\label{eq:v1-c7-4}
\frac{\partial u_{em}}{\partial t} + \nabla\cdot(\mathbf{E}\times\mathbf{H})
= -\,\mathbf{J}\cdot\mathbf{E},
\qquad
u_{em} = \tfrac{1}{2}\left(\mathbf{E}\cdot\mathbf{D} +
\mathbf{H}\cdot\mathbf{B}\right).
\end{equation}
A Equação~\eqref{eq:v1-c7-4} é, mais uma vez, a Equação~\eqref{eq:v1-c1-1} com
$\mathcal{B}$ igual à energia: $u_{em}$ é a densidade de energia acumulada no
campo, o vetor de Poynting $\mathbf{E}\times\mathbf{H}$ é o fluxo que atravessa
a fronteira, e $\mathbf{J}\cdot\mathbf{E}$ é a taxa com que energia
eletromagnética sai do campo e entra na matéria. Em um condutor ôhmico esse
termo vale $\sigma_e|\mathbf{E}|^2 \geq 0$: é a perda Joule, e o sinal não
negativo garantido pela positividade de $\sigma_e$ é exatamente a
manifestação eletromagnética da desigualdade de entropia,
Equação~\eqref{eq:v1-c2-26}. Toda perda Joule deste capítulo --- no cabo, no
estator, no rotor, no IGBT --- é geração de entropia no sentido do
Capítulo~\ref{sec:termodinamica}, e aparece como fonte térmica nos balanços do
Capítulo~\ref{sec:calor}. Não há duas contabilidades: há uma só.

Uma consequência da Equação~\eqref{eq:v1-c7-4} que costuma surpreender e que
importa para o dimensionamento térmico: a energia não viaja \emph{dentro} do
condutor, mas no campo ao redor dele. O condutor guia a energia; o que se
dissipa dentro dele é apenas a parcela $\mathbf{J}\cdot\mathbf{E}$ que
``vaza'' para a matéria. É por isso que a temperatura de um cabo depende do
meio que o cerca, e não apenas da corrente que ele conduz.

\subsection{A hipótese de parâmetros concentrados e seus
limites}\label{sub:v1-c7-concentrados}

A Seção~\ref{sub:v1-c2-hipotese} mostrou que tratar matéria discreta como campo
contínuo exige uma separação de escalas, medida pelo número de Knudsen. A
teoria de circuitos exige uma redução análoga, e em sentido oposto: ela
substitui \emph{campos distribuídos no espaço} por um punhado de elementos
concentrados ($R$, $L$, $C$) ligados por fios ideais. A pergunta é a mesma ---
sob qual condição a substituição é legítima? --- e a resposta tem a mesma
forma: existe um número adimensional que decide.

A condição é que o circuito seja \textbf{eletricamente curto}: que o tempo de
propagação de um sinal através dele seja desprezível em face do tempo em que
os sinais variam. Equivalentemente, que a maior dimensão $\ell$ do circuito
seja muito menor que o comprimento de onda $\lambda_{em}$ associado à
frequência de operação, como na Equação~\eqref{eq:v1-c7-5}:
\begin{equation}\label{eq:v1-c7-5}
\frac{\ell}{\lambda_{em}} \ll 1,
\qquad \lambda_{em} = \frac{v_p}{f},
\end{equation}
onde $v_p$ é a velocidade de propagação no meio --- próxima de $c$ no ar,
tipicamente $1{,}5$ a $2\times10^{8}\,$m\,s$^{-1}$ em cabo isolado. Satisfeita a
Equação~\eqref{eq:v1-c7-5}, os atrasos de propagação somem, o campo em cada
instante é o campo estático correspondente às fontes daquele instante (regime
\textbf{quase-estático}), e as duas leis de Kirchhoff se tornam teoremas em vez
de postulados.

\textbf{De Maxwell a Kirchhoff.} A lei dos nós é a
Equação~\eqref{eq:v1-c7-1} integrada sobre um volume que envolve o nó, com o
termo de acúmulo $\partial\rho_q/\partial t$ desprezado --- o nó não armazena
carga:
\begin{equation}\label{eq:v1-c7-6}
\sum_k I_k = 0 .
\end{equation}
A lei das malhas é a lei de Faraday integrada sobre a superfície apoiada na
malha, com o fluxo magnético concatenado desprezado ou, quando não desprezível,
\emph{atribuído a um indutor} explícito no lugar de permanecer distribuído:
\begin{equation}\label{eq:v1-c7-7}
\sum_k V_k = 0 .
\end{equation}
Esta última observação é a mais importante das duas, porque explica o que um
elemento concentrado realmente é: $L$ não é um componente físico, é o
\emph{nome} dado ao fluxo magnético que a malha concatena; $C$ é o nome dado à
carga que se acumula entre dois condutores; $R$ é o nome dado à parcela
$\mathbf{J}\cdot\mathbf{E}$ dissipada. Todo campo que sobra fora dos elementos
declarados no esquema não deixa de existir --- ele apenas deixa de ser
modelado. Os modos de falha da hipótese, abaixo, são exatamente os casos em que
o campo não modelado é grande demais para ser ignorado.

\textbf{Modo 1: propagação e reflexão de onda.} A frequência que importa na
Equação~\eqref{eq:v1-c7-5} não é a frequência fundamental, e sim a maior
frequência com conteúdo relevante. Para a rede em $60\,$Hz,
$\lambda_{em}\approx 5000\,$km, e mesmo uma instalação offshore inteira é
eletricamente minúscula --- a razão da Equação~\eqref{eq:v1-c7-5} é da ordem de
$10^{-5}$. Para o mesmo cabo alimentado por um inversor PWM, porém, o que
trafega são degraus de tensão com tempo de subida $t_r$ da ordem de
$50$ a $200\,$ns, cujo conteúdo espectral se estende até cerca de
$0{,}35/t_r$, isto é, alguns megahertz. O comprimento de onda cai para a casa
de dezenas de metros, e a razão da Equação~\eqref{eq:v1-c7-5} deixa de ser
pequena.

A consequência quantitativa é o \textbf{comprimento crítico de cabo}: quando o
cabo é longo o bastante para que o degrau leve mais de meio tempo de subida
para percorrê-lo e voltar --- com $v_p = 1{,}5\times10^{8}\,$m\,s$^{-1}$ e
$t_r = 100\,$ns, isso ocorre a partir de cerca de $7{,}5\,$m de cabo ---, a onda refletida na descontinuidade de impedância
--- o motor, cuja impedância de surto é uma ordem de grandeza maior que a do
cabo --- chega ao terminal antes que o degrau termine e se soma a ele. A
tensão nos terminais do motor pode aproximar-se do dobro da tensão do
barramento CC. É este o mecanismo que degrada o isolamento do enrolamento e que
o Volume II modela em suas consequências, pela lei de Arrhenius: a vida do
isolamento cai porque a solicitação elétrica por evento é maior, não porque a
corrente média subiu. Um modelo de circuito concentrado não tem como prevê-lo,
porque nele o cabo é um fio ideal e não existe onda alguma para refletir.

\textbf{Modo 2: efeito pelicular e proximidade.} O parâmetro $R$ só é constante
enquanto a corrente se distribui uniformemente na seção do condutor. Campo
magnético variável induz correntes que se opõem ao próprio fluxo, expulsando a
corrente para a periferia; a profundidade característica dessa penetração é a
\textbf{profundidade pelicular} da Equação~\eqref{eq:v1-c7-8}:
\begin{equation}\label{eq:v1-c7-8}
\delta_p = \sqrt{\frac{2}{\omega\,\mu_m\,\sigma_e}} .
\end{equation}
Para o cobre em $60\,$Hz, $\delta_p \approx 8{,}5\,$mm --- maior que o raio da
maioria dos condutores usados, e por isso o efeito é modesto na frequência
industrial. Em alguns megahertz, $\delta_p$ cai para dezenas de micrometros: a
mesma seção de cobre oferece uma resistência ordens de grandeza maior às
componentes de alta frequência do PWM do que à fundamental. O corolário
prático é que perdas de cabo e de enrolamento sob acionamento por inversor não
podem ser estimadas com a resistência em corrente contínua.

\textbf{Modo 3: capacitâncias parasitas.} O terceiro modo é o mais silencioso e
o de maior consequência mecânica. Entre o enrolamento do estator e o núcleo,
entre estator e rotor através do entreferro, e entre rotor e carcaça através do
filme de óleo do rolamento, existem capacitâncias que nenhum esquema elétrico
declara --- e que a lei de Gauss garante existirem, porque há condutores
separados por dielétricos. Em $60\,$Hz sua impedância $1/(\omega C)$ é tão alta
que elas são, de fato, circuito aberto. Nas frequências de comutação e,
sobretudo, nos $\mathrm{d}v/\mathrm{d}t$ das bordas do PWM, sua impedância cai
o suficiente para conduzir corrente mensurável. Essa corrente encontra o
caminho que houver, e um dos caminhos disponíveis atravessa o rolamento. Os
elementos $C_{sr}$, $C_{rg}$ e $C_{\text{bearing}}$ que o Volume II combina em
um divisor capacitivo para estimar a tensão de eixo são precisamente esses
capacitores não declarados; a Seção~\ref{sub:v1-c7-inversor} fecha a origem da
tensão que os excita.

A Tabela~\ref{tab:v1-c7-1} organiza o critério e os três modos, na mesma
função que a Tabela~\ref{tab:v1-c2-1} exerce no Capítulo~\ref{sec:continuo}:
dizer onde a teoria para.

\begin{table}[H]
\centering
\caption[Validade da hipótese de parâmetros concentrados]{Hipótese de
parâmetros concentrados: critério de validade e modos de falha relevantes em
acionamento por inversor.}
\label{tab:v1-c7-1}
\small
\resizebox{\textwidth}{!}{%
\begin{tabular}{@{}llll@{}}
\toprule
\textbf{Regime} & \textbf{Critério} & \textbf{O que deixa de valer} & \textbf{Consequência física} \\
\midrule
Quase-estático & $\ell/\lambda_{em}\ll1$ & --- (hipótese válida) & Circuito $R$, $L$, $C$ é exato o bastante \\
Propagação & $\ell \gtrsim v_p t_r/2$ & Fio ideal & Reflexão; sobretensão no terminal do motor \\
Pelicular & $r_{\text{cond}} \gtrsim \delta_p$ & $R$ constante & Perda extra nas componentes de alta frequência \\
Capacitivo & $1/(\omega C)$ comparável às demais & Circuito aberto entre partes & Corrente de modo comum; corrente de eixo \\
\bottomrule
\end{tabular}}
\end{table}

\subsection{Relações constitutivas do meio}\label{sub:v1-c7-constitutivas}

As equações de Maxwell relacionam quatro campos, mas são apenas duas equações
vetoriais independentes: faltam relações que digam como o \emph{material}
responde. Elas têm exatamente o mesmo estatuto epistemológico das relações
constitutivas da Seção~\ref{sub:v1-c2-constitutivas} --- não são dedutíveis das
leis de balanço, são medidas --- e as três centrais são a
Equação~\eqref{eq:v1-c7-9}:
\begin{equation}\label{eq:v1-c7-9}
\mathbf{J} = \sigma_e\mathbf{E},
\qquad
\mathbf{D} = \varepsilon_e\mathbf{E},
\qquad
\mathbf{B} = \mu_m\mathbf{H} .
\end{equation}
A primeira é a \textbf{lei de Ohm na forma local}. A forma familiar
$V = R\,I$ da Equação~\eqref{eq:v1-c7-10} é a sua integração ao longo de um
condutor de comprimento $L_c$ e seção $A_c$ uniformes:
\begin{equation}\label{eq:v1-c7-10}
V = R\,I,
\qquad
R = \frac{L_c}{\sigma_e A_c} .
\end{equation}
Vale registrar o que essa integração pressupõe, porque é o mesmo tipo de
premissa oculta que o Capítulo~\ref{sec:continuo} insiste em explicitar:
seção uniforme, corrente distribuída uniformemente nela (o que a
Equação~\eqref{eq:v1-c7-8} já mostrou não ser garantido) e $\sigma_e$
homogêneo e isotrópico.

\textbf{A condutividade depende da temperatura, e isso não é detalhe.} Em
metais, a resistividade cresce quase linearmente com a temperatura na faixa de
operação de máquinas, conforme a Equação~\eqref{eq:v1-c7-11}:
\begin{equation}\label{eq:v1-c7-11}
\rho_{e}(T) = \rho_{e,20}\left[1 + \alpha_T\,(T - 20\,^{\circ}\mathrm{C})\right],
\qquad
\sigma_e = 1/\rho_e ,
\end{equation}
com $\rho_e$ a resistividade elétrica --- que não é a densidade de carga
$\rho_q$ da Equação~\eqref{eq:v1-c7-1} nem a massa específica $\rho$ dos
capítulos anteriores; a tripla coincidência de letra é herdada da notação
consagrada e resolvida aqui pelos subscritos ---, e
com $\alpha_T \approx 3{,}9\times10^{-3}\,\mathrm{K}^{-1}$ para o cobre. O
efeito é grande: um enrolamento que sai de $20\,^{\circ}$C e estabiliza em
$100\,^{\circ}$C tem resistência cerca de $31\,\%$ maior. Três consequências
percorrem o resto desta coleção. Primeira: a perda Joule de um motor quente não
é a de um motor frio para a mesma corrente. Segunda: a resistência do rotor
$R_r$ que aparece na expressão de torque cresce com a temperatura, deslocando a
curva torque--escorregamento --- é exatamente por essa via, e não por um termo
aditivo, que o rendimento do motor no Volume II depende de $T$. Terceira, e
mais útil na prática de manutenção: a própria resistência do enrolamento é um
termômetro. Medir $R$ e inverter a Equação~\eqref{eq:v1-c7-11} dá a temperatura
\emph{média} do cobre, que nenhum RTD pontual mede.

\textbf{O meio magnético não é linear.} A terceira relação da
Equação~\eqref{eq:v1-c7-9} é a mais frágil das três. Em materiais
ferromagnéticos --- o aço do estator e do rotor --- $\mu_m$ não é constante:
cresce, atinge máximo e desaba quando os domínios magnéticos se alinham e o
material \textbf{satura}, tipicamente entre $1{,}5$ e $2{,}0\,$T nos aços
elétricos usuais. Acima da saturação, aumentar a corrente praticamente não
aumenta o fluxo; apenas aquece o cobre. Toda a lógica do controle
tensão/frequência da Seção~\ref{sub:v1-c7-inversor} existe por causa dessa
curva.

Além de não linear, a relação é \emph{com memória}: o percurso de magnetização
descreve um laço de \textbf{histerese}, cuja área é energia perdida por ciclo.
Somada às correntes parasitas induzidas no próprio núcleo, ela compõe a perda
no ferro, descrita empiricamente pela relação de \textbf{Steinmetz},
Equação~\eqref{eq:v1-c7-12}:
\begin{equation}\label{eq:v1-c7-12}
P_{fe} = \underbrace{k_h\,f\,\hat{B}^{\,a}}_{\text{histerese}} +
\underbrace{k_e\,f^{2}\,\hat{B}^{2}}_{\text{correntes parasitas}},
\end{equation}
onde $\hat{B}$ é a amplitude da indução, $f$ a frequência, e $k_h$, $k_e$ e o
expoente $a$ (usualmente entre $1{,}6$ e $2{,}2$) são constantes ajustadas por
ensaio do material. A dependência quadrática do segundo termo com $f$ é a razão
física de o núcleo ser laminado e não maciço: laminar reduz a área dos
percursos disponíveis para as correntes parasitas. E é também a razão de o
conteúdo harmônico do PWM aquecer o ferro mais do que a fundamental sozinha
sugeriria. A Equação~\eqref{eq:v1-c7-12} é o conteúdo físico da parcela
``núcleo'' que o Volume II agrega, junto com atrito e ventilação, na constante
$P_{\text{fixas}}$ de ensaio.

A Tabela~\ref{tab:v1-c7-2} reúne as relações constitutivas deste capítulo no
mesmo formato da Tabela~\ref{tab:v1-c2-3}, para que a comparação entre as duas
seja imediata: em ambos os casos, o que entra na tabela são parâmetros
\emph{medidos}, com faixa de validade declarada.

\begin{table}[H]
\centering
\caption[Relações constitutivas eletromagnéticas]{Relações constitutivas
eletromagnéticas, parâmetros e premissas --- contraparte da
Tabela~\ref{tab:v1-c2-3}.}
\label{tab:v1-c7-2}
\small
\resizebox{\textwidth}{!}{%
\begin{tabular}{@{}llll@{}}
\toprule
\textbf{Meio/mecanismo} & \textbf{Relação constitutiva} & \textbf{Parâmetros} & \textbf{Premissa} \\
\midrule
Condutor ôhmico & $\mathbf{J} = \sigma_e\mathbf{E}$ & $\sigma_e(T)$ & Isotropia; $\sigma_e$ independente de $|\mathbf{E}|$ \\
Dielétrico linear & $\mathbf{D} = \varepsilon_e\mathbf{E}$ & $\varepsilon_e$ & Sem polarização remanente; abaixo da rigidez \\
Meio magnético linear & $\mathbf{B} = \mu_m\mathbf{H}$ & $\mu_m$ & Longe da saturação; sem histerese \\
Aço elétrico (real) & $B = B(H)$, com laço & $\hat{B}_{sat}$, $k_h$, $k_e$, $a$ & Ensaio na faixa de $f$ e $\hat{B}$ de uso \\
Perda no núcleo & Eq.~\eqref{eq:v1-c7-12} & $k_h$, $k_e$, $a$ & Ajuste empírico; laminação preservada \\
\bottomrule
\end{tabular}}
\end{table}

\subsection{Regime permanente senoidal e fasores}\label{sub:v1-c7-fasores}

Aceita a hipótese de parâmetros concentrados, resta resolver as equações
diferenciais do circuito. Quando a excitação é senoidal de frequência única e o
circuito é linear e invariante no tempo, existe um atalho que não custa
nenhuma aproximação adicional.

O argumento é o seguinte, e convém enunciá-lo porque é ele que delimita o uso.
A resposta de um sistema linear invariante no tempo a $e^{j\omega t}$ é
$H(j\omega)\,e^{j\omega t}$: a exponencial complexa é autofunção do sistema, e
só a amplitude e a fase mudam. Representar cada grandeza senoidal pelo número
complexo que carrega essa amplitude e essa fase --- o \textbf{fasor} --- e
substituir $\mathrm{d}/\mathrm{d}t$ por $j\omega$ converte as equações
diferenciais do circuito em equações algébricas complexas. A
Equação~\eqref{eq:v1-c7-13} define a impedância que resulta dessa substituição
para os três elementos e para a associação série:
\begin{equation}\label{eq:v1-c7-13}
Z_R = R,
\quad
Z_L = j\omega L,
\quad
Z_C = \frac{1}{j\omega C},
\quad
Z_{RLC} = R + j\left(\omega L - \frac{1}{\omega C}\right).
\end{equation}
A parte real de $Z$ é onde há dissipação; a parte imaginária, a
\textbf{reatância}, é onde há apenas armazenamento --- no campo magnético do
indutor ou no campo elétrico do capacitor, os dois termos de $u_{em}$ na
Equação~\eqref{eq:v1-c7-4}. (O símbolo $Z$ aqui não deve ser confundido com
$Z_c$, o fator de compressibilidade do Capítulo~\ref{sec:termodinamica}, como
já registrado na Nomenclatura.)

O formalismo tem duas condições de morte, e as duas são realidade cotidiana em
um acionamento: ele exige \textbf{linearidade} --- e a saturação magnética da
Seção~\ref{sub:v1-c7-constitutivas} é não linear --- e exige \textbf{frequência
única} --- e a saída de um inversor PWM é tudo menos uma senoide pura. A saída
não é abandonar os fasores, e sim aplicá-los frequência a frequência
(superposição) e depois somar as contribuições, procedimento que a
Seção~\ref{sub:v1-c7-potencia} mostra ter consequências não triviais na
contabilidade de potência.

Um resultado da Equação~\eqref{eq:v1-c7-13} merece destaque por aparecer onde
não se espera. Em $\omega_0 = 1/\sqrt{LC}$ as duas reatâncias se cancelam e a
impedância cai a $R$: é a \textbf{ressonância}. O conjunto formado pela
indutância do cabo e pela capacitância parasita do conjunto cabo--motor
possui ressonância própria, tipicamente na faixa de centenas de quilohertz.
Quando o espectro do PWM excita essa frequência, a oscilação que aparece
superposta ao degrau de tensão nos terminais do motor não é ruído de medição:
é o circuito ressoando, e ela se soma à sobretensão de reflexão discutida na
Seção~\ref{sub:v1-c7-concentrados}.

\subsection{Potência: ativa, reativa e o efeito dos
harmônicos}\label{sub:v1-c7-potencia}

A potência instantânea é o produto $p_{el}(t) = v(t)\,i(t)$ --- sem qualquer
hipótese. Para tensão e corrente senoidais defasadas de $\varphi$, a média
sobre um ciclo produz a \textbf{potência ativa} da
Equação~\eqref{eq:v1-c7-14}, que é a única parcela que efetivamente transporta
energia da fonte para a carga:
\begin{equation}\label{eq:v1-c7-14}
P_{el} = V_{rms}\,I_{rms}\cos\varphi .
\end{equation}
O produto $V_{rms}I_{rms}$ sem o cosseno é a \textbf{potência aparente}
$S_{el}$, e a parcela em quadratura é a \textbf{potência reativa} $Q_{el}$;
juntas formam o triângulo de potências da Equação~\eqref{eq:v1-c7-15}:
\begin{equation}\label{eq:v1-c7-15}
Q_{el} = V_{rms}\,I_{rms}\sin\varphi,
\qquad
S_{el} = V_{rms}\,I_{rms},
\qquad
S_{el}^{2} = P_{el}^{2} + Q_{el}^{2} .
\end{equation}

Vale dizer o que $Q_{el}$ \emph{é}, porque a definição algébrica esconde a
física. A potência reativa é energia que vai e volta entre a fonte e os campos
de armazenamento do circuito --- o campo magnético do entreferro e das
indutâncias de dispersão, no caso de um motor --- com média nula ao longo do
ciclo. Ela não é consumida. Mas é transportada: o cabo, o transformador e o
gerador precisam conduzir a corrente correspondente, e essa corrente dissipa
$R\,I_{rms}^2$ em cada elemento do caminho. É essa assimetria --- reativo não
consumido, mas pago em perda de transporte --- que justifica economicamente a
correção do fator de potência. O ganho é quadrático: elevar $\cos\varphi$ de
$0{,}70$ para $0{,}95$ reduz a corrente em $26\,\%$ e as perdas de transporte
em $46\,\%$, para a mesma potência ativa entregue.

\textbf{Por que trifásico.} Um sistema monofásico entrega potência
\emph{pulsante}: com $v = \hat{V}\cos\omega t$ e
$i = \hat{I}\cos(\omega t - \varphi)$, o produto contém um termo constante e um
termo oscilando em $2\omega$. Em um motor, esse termo de $2\omega$ é torque
pulsante, com todas as consequências vibratórias que o
Capítulo~\ref{sec:mecanica} trata. Em um sistema trifásico equilibrado, ao
contrário, as três contribuições instantâneas somam-se a uma constante --- as
parcelas oscilantes das três fases estão defasadas de $120^{\circ}$ e se
cancelam --- e a potência total é a da Equação~\eqref{eq:v1-c7-16}:
\begin{equation}\label{eq:v1-c7-16}
p_{el,3\phi}(t) = 3\,V_{f}I_{f}\cos\varphi = \sqrt{3}\,V_L I_L\cos\varphi
= \text{constante},
\end{equation}
onde $V_f$, $I_f$ são grandezas de fase e $V_L$, $I_L$ de linha, com
$V_L = \sqrt{3}\,V_f$ em ligação estrela. A constância da
Equação~\eqref{eq:v1-c7-16} é a verdadeira razão de o acionamento industrial
ser trifásico, e não uma questão de economia de cobre: torque constante, sem
pulsação em $2f$.

\textbf{O que os harmônicos fazem com essa contabilidade.} A
Equação~\eqref{eq:v1-c7-15} pressupõe senoides. O retificador de entrada de um
inversor drena corrente em pulsos, não em senoide, e nesse caso o triângulo de
potências deixa de fechar: existe uma terceira parcela, a potência de
distorção, e $S_{el}^2 = P_{el}^2 + Q_{el}^2 + D^2$. A distorção é quantificada
pela taxa de distorção harmônica total da Equação~\eqref{eq:v1-c7-17}:
\begin{equation}\label{eq:v1-c7-17}
\mathrm{THD}_i = \frac{\sqrt{\sum_{h\geq2} I_h^2}}{I_1},
\end{equation}
onde $I_h$ é o valor eficaz do harmônico de ordem $h$ e $I_1$ o da fundamental.
O fator de potência verdadeiro se decompõe então em dois fatores distintos,
Equação~\eqref{eq:v1-c7-18} --- o de deslocamento, que é o $\cos\varphi$
clássico, e o de distorção, que não tem nada a ver com defasagem:
\begin{equation}\label{eq:v1-c7-18}
\mathrm{FP} = \underbrace{\cos\varphi_1}_{\text{deslocamento}}\cdot
\underbrace{\frac{1}{\sqrt{1+\mathrm{THD}_i^{2}}}}_{\text{distorção}} .
\end{equation}
A Equação~\eqref{eq:v1-c7-18} explica um resultado de campo que confunde:
um acionamento com retificador a diodos apresenta $\cos\varphi_1$ próximo da
unidade --- e, ainda assim, fator de potência medido em torno de $0{,}9$ ou
menos, inteiramente por conta do segundo fator. Corrigi-lo com banco de
capacitores não funciona; pior, o banco pode ressoar com a indutância da rede
na frequência de algum harmônico. A correção correta é reduzir a própria
distorção (reator de linha, retificador de 12 pulsos, filtro ativo). Além
disso, cada harmônico de corrente dissipa perda Joule adicional --- agravada
pelo efeito pelicular da Equação~\eqref{eq:v1-c7-8} --- e perda extra no ferro
pela Equação~\eqref{eq:v1-c7-12}: é essa soma que fundamenta o fator de
desclassificação (\textit{derating}) que o Volume II calcula a partir de
$\mathrm{THD}_i$.

\subsection{Sistemas trifásicos e componentes
simétricas}\label{sub:v1-c7-trifasico}

A Equação~\eqref{eq:v1-c7-16} vale para um sistema \emph{equilibrado}. Redes
reais --- e as redes isoladas de plataformas offshore em particular, onde não
há barramento infinito para absorver assimetrias de carga --- não são
perfeitamente equilibradas. É preciso, portanto, um formalismo capaz de
descrever três fasores quaisquer, e ele existe desde Fortescue: qualquer terno
de fasores $(\mathbf{V}_a,\mathbf{V}_b,\mathbf{V}_c)$ se decompõe de maneira
única na soma de três ternos --- um de \textbf{sequência positiva}
(equilibrado, girando na ordem $abc$), um de \textbf{sequência negativa}
(equilibrado, girando na ordem inversa) e um de \textbf{sequência zero} (as
três fases em fase). Com $a = e^{\,j2\pi/3}$, a decomposição é a
Equação~\eqref{eq:v1-c7-19}:
\begin{equation}\label{eq:v1-c7-19}
\begin{bmatrix}\mathbf{V}_0\\ \mathbf{V}_+\\ \mathbf{V}_-\end{bmatrix}
= \frac{1}{3}
\begin{bmatrix}
1 & 1 & 1\\
1 & a & a^2\\
1 & a^2 & a
\end{bmatrix}
\begin{bmatrix}\mathbf{V}_a\\ \mathbf{V}_b\\ \mathbf{V}_c\end{bmatrix} .
\end{equation}

A utilidade da Equação~\eqref{eq:v1-c7-19} não é algébrica, é física, e está
concentrada em um único fato: \textbf{a máquina de indução apresenta impedâncias
muito diferentes a cada sequência}. À sequência positiva ela oferece a
impedância de operação normal, alta, porque o campo gira junto com o rotor e o
escorregamento visto é o pequeno $s$ de regime. À sequência negativa o campo
gira ao contrário, e o rotor a vê com escorregamento próximo de $2-s\approx2$:
a impedância correspondente é a de rotor bloqueado, tipicamente de um quinto a
um sexto da de sequência positiva.

Daí decorrem, imediatamente, os três sintomas que o Volume II monitora:

\begin{itemize}
	\item \textbf{Amplificação do desbalanço.} Um desbalanço de tensão de
	poucos pontos percentuais produz um desbalanço de corrente cinco a seis
	vezes maior, na razão inversa das impedâncias de sequência. É por isso que
	a corrente é o indicador sensível, e a tensão o indicador tardio.
	\item \textbf{Aquecimento localizado.} A corrente de sequência negativa
	dissipa perda Joule concentrada, e o faz predominantemente no rotor, onde
	é induzida na frequência $(2-s)f \approx 2f$ --- justamente onde o efeito
	pelicular da Equação~\eqref{eq:v1-c7-8} concentra a corrente na superfície
	das barras. O rotor aquece desproporcionalmente, e o rotor é a parte que
	nenhum RTD instrumenta.
	\item \textbf{Torque pulsante em $2f$.} A interação entre o campo de
	sequência positiva e as correntes de sequência negativa produz um torque
	alternado em $2f$, que aparece no espectro de vibração e no espectro de
	corrente --- um sintoma elétrico que se manifesta como sintoma mecânico.
\end{itemize}

O primeiro item é a base física do índice de desbalanço que o Volume II
computa segundo a NEMA MG~1 \cite{nemamg1}, e da razão de essa norma
recomendar um limite tão apertado --- $1\,\%$ nos terminais em carga nominal
--- com curva de desclassificação acima desse valor: o desbalanço de tensão
parece pequeno porque é medido na grandeza errada.

A sequência zero merece uma nota, porque explica um detalhe de projeto: em
ligação triângulo ou estrela isolada, as correntes de sequência zero não têm
caminho de retorno e simplesmente não circulam. Esse é um dos motivos de
motores industriais serem ligados sem neutro --- e também a razão de a tensão
de modo comum gerada pelo inversor, que é uma grandeza de sequência zero, não
encontrar caminho pelos enrolamentos e procurar caminho pelas capacitâncias
parasitas da Seção~\ref{sub:v1-c7-concentrados}.

\textbf{Nota sobre transformações de eixos.} Para regime permanente
senoidal, fasores e componentes simétricas bastam. Para transitórios --- e para
controle vetorial --- é necessário um modelo de coeficientes constantes, o que
se obtém pelas transformações de Clarke ($abc \to \alpha\beta$) e Park
($\alpha\beta \to dq$, em referencial girante). O efeito de Park é eliminar a
dependência explícita do tempo nas indutâncias mútuas entre estator e rotor,
convertendo a máquina em algo formalmente equivalente a uma máquina de corrente
contínua, com uma componente de corrente controlando fluxo e outra controlando
torque. É esse o fundamento do controle vetorial mencionado no
Volume II; o desenvolvimento completo pertence à literatura de acionamentos
\cite{fitzgerald2003}.

\subsection{Conversão eletromecânica de energia}\label{sub:v1-c7-conversao}

Chega-se ao ponto que justifica o capítulo dentro de um volume sobre máquinas
rotodinâmicas: como campo eletromagnético vira torque de eixo. A resposta não
depende do tipo de máquina, e é mais um balanço de energia sobre uma fronteira
--- a Equação~\eqref{eq:v1-c1-1} outra vez, agora com a fronteira envolvendo o
\emph{campo de acoplamento} entre a parte elétrica e a parte mecânica:
\begin{equation}\label{eq:v1-c7-20}
\underbrace{\mathrm{d}W_{el}}_{\text{entra pelos terminais}} =
\underbrace{\mathrm{d}W_{campo}}_{\text{acumula no campo}} +
\underbrace{\mathrm{d}W_{mec}}_{\text{sai pelo eixo}} +
\underbrace{\mathrm{d}W_{perdas}}_{\text{Joule + ferro}} .
\end{equation}
A Equação~\eqref{eq:v1-c7-20} diz que não há conversão direta de energia
elétrica em mecânica: toda ela transita pelo campo de acoplamento. Essa
observação, que parece formal, tem consequência prática imediata --- a
capacidade de torque de uma máquina é limitada pelo quanto de energia ela
consegue armazenar no entreferro, e não pela corrente que os terminais
suportam.

Isolando a parcela mecânica, e trabalhando com a \textbf{coenergia} magnética
$W'_m(i,\theta)$ --- a área complementar sob a curva $\lambda(i)$, que é a forma
conveniente quando a variável independente é a corrente --- o torque
eletromagnético resulta da Equação~\eqref{eq:v1-c7-21}:
\begin{equation}\label{eq:v1-c7-21}
\tau_{em} = \left.\frac{\partial W'_m(i,\theta)}{\partial\theta}\right|_{i=\text{const}},
\qquad
W'_m = \int_0^i \lambda(i',\theta)\,\mathrm{d}i' ,
\end{equation}
onde $\theta$ é a posição angular do rotor e $\lambda$ o fluxo concatenado. A
Equação~\eqref{eq:v1-c7-21} é o resultado central da conversão eletromecânica e
contém, como casos particulares, todos os torques de máquina elétrica: existe
torque sempre que a energia armazenada no campo \emph{depende da posição do
rotor}. Se essa dependência vem da variação de relutância, tem-se o torque de
relutância; se vem da variação da indutância mútua entre enrolamentos de
estator e rotor, tem-se o torque das máquinas de indução e síncronas. Se o
fluxo concatenado não dependesse de $\theta$, não haveria torque algum, por
mais corrente que se injetasse.

\textbf{A ponte com o Capítulo~\ref{sec:continuo}.} Há uma segunda maneira de
calcular a mesma força, e ela é conceitualmente mais satisfatória porque
reconduz o eletromagnetismo à linguagem de tensão do Capítulo~2. A força por
unidade de área na superfície do rotor pode ser escrita como o divergente de um
tensor, o \textbf{tensor de tensões de Maxwell}, cujas componentes normal e
tangencial no entreferro são a Equação~\eqref{eq:v1-c7-22}:
\begin{equation}\label{eq:v1-c7-22}
\sigma_{n} = \frac{B_n^{2} - B_t^{2}}{2\mu_0},
\qquad
\sigma_{t} = \frac{B_n B_t}{\mu_0} .
\end{equation}
onde $B_n$ e $B_t$ são as componentes normal e tangencial da indução no
entreferro e $\mu_0 = 4\pi\times10^{-7}\,$H\,m$^{-1}$ é a permeabilidade do
vácuo, aplicável porque o entreferro é ar.
A componente tangencial $\sigma_t$ é uma \emph{tensão de cisalhamento
aplicada na superfície do rotor}, no sentido exato do vetor tensão
$\mathbf{t}_{(\mathbf{n})}$ do teorema de Cauchy,
Seção~\ref{sub:v1-c2-tensao}. O torque é a sua integral sobre a superfície do
entreferro, multiplicada pelo raio, o que dá a Equação~\eqref{eq:v1-c7-23}:
\begin{equation}\label{eq:v1-c7-23}
\tau_{em} = 2\pi r^{2} L_{fe}\,\bar{\sigma}_{t},
\end{equation}
com $r$ o raio do entreferro, $L_{fe}$ o comprimento axial do pacote e
$\bar{\sigma}_t$ a tensão tangencial média. A leitura de projeto da
Equação~\eqref{eq:v1-c7-23} é direta e é a razão de ela estar aqui: como
$\bar{\sigma}_t$ é limitada pela saturação do aço a algumas dezenas de
quilopascals em máquinas industriais fechadas, \textbf{o torque é proporcional
ao volume do rotor}. Potência pode ser obtida aumentando a rotação; torque, só
aumentando volume. É essa relação que explica por que a carcaça de um motor
acompanha o torque nominal, e não a potência --- e por que a alternativa a um
motor grande de baixa rotação é um motor menor de alta rotação com redutor, que
é precisamente a escolha de arquitetura do trem de potência analisado no
Volume II.

\subsection{O campo girante e a máquina de indução}\label{sub:v1-c7-inducao}

\textbf{Por que existe um campo girante.} Três enrolamentos deslocados de
$120^{\circ}$ no espaço, percorridos por correntes deslocadas de $120^{\circ}$
no tempo, produzem uma força magnetomotriz resultante cuja soma trigonométrica
é uma onda de amplitude constante que \emph{gira} no entreferro --- não uma
onda pulsante. Nenhuma peça se move para produzi-la; a rotação é do padrão
espacial de campo. A velocidade dessa rotação é fixada pela frequência e pelo
número de polos $p$ do enrolamento, e é a \textbf{velocidade síncrona} da
Equação~\eqref{eq:v1-c7-24}:
\begin{equation}\label{eq:v1-c7-24}
n_s = \frac{120\,f}{p} \quad [\mathrm{rpm}],
\qquad
\omega_s = \frac{2\pi n_s}{60} \quad [\mathrm{rad\,s^{-1}}] .
\end{equation}
A velocidade angular síncrona $\omega_s$ da Equação~\eqref{eq:v1-c7-24} é
mecânica, como no Volume II; a angular elétrica correspondente é
$\omega_e = 2\pi f = (p/2)\,\omega_s$, e as duas coincidem apenas na máquina de
dois polos. Aqui $p$ é o \emph{número de polos} (não o número de pares), convenção adotada
para coincidir com a do Volume II; o número de pares é $p/2$. A colisão com $p$
minúsculo de pressão, dos capítulos de termodinâmica, é aparente apenas: não há
contexto em que as duas grandezas compareçam à mesma equação. Um motor de
quatro polos em $60\,$Hz gira, portanto, com campo a $1800\,$rpm.

\textbf{Por que o rotor precisa escorregar.} A máquina de indução não tem
alimentação no rotor: a corrente que o percorre é \emph{induzida}. Pela lei de
Faraday, Equação~\eqref{eq:v1-c7-2}, só há tensão induzida se houver variação
de fluxo concatenado, e só há variação se houver \emph{movimento relativo}
entre o campo girante e o rotor. Se o rotor alcançasse a velocidade síncrona,
não haveria movimento relativo, nem tensão induzida, nem corrente, nem torque.
A máquina de indução é obrigada a operar abaixo do síncrono, e a medida dessa
diferença é o \textbf{escorregamento} da Equação~\eqref{eq:v1-c7-25}:
\begin{equation}\label{eq:v1-c7-25}
s = \frac{n_s - n}{n_s} .
\end{equation}
(Este $s$ é o escorregamento, não a entropia específica do
Capítulo~\ref{sec:termodinamica}; a coincidência de símbolo é consagrada nas
duas disciplinas e foi mantida por compatibilidade com o Volume II.) Em carga
nominal, $s$ fica tipicamente entre $1\,\%$ e $3\,\%$ --- o que faz do próprio
escorregamento uma medida indireta de carga, e é assim que ele é usado no
Volume II.

\textbf{A repartição da potência de entreferro.} O circuito equivalente por
fase representa o rotor pela resistência $R_r/s$ referida ao estator. A
potência que atravessa o entreferro é a Equação~\eqref{eq:v1-c7-26}:
\begin{equation}\label{eq:v1-c7-26}
P_{ag} = 3\,I_r^{2}\,\frac{R_r}{s} .
\end{equation}
Dessa potência, a parcela efetivamente dissipada no cobre do rotor é
$3I_r^2R_r$, e o restante é potência mecânica. Dividindo uma pela outra
obtém-se a repartição exata da Equação~\eqref{eq:v1-c7-27}, que é um dos
resultados mais úteis de toda a teoria de máquinas de indução:
\begin{equation}\label{eq:v1-c7-27}
P_{j,rotor} = s\,P_{ag},
\qquad
P_{mec} = (1-s)\,P_{ag} .
\end{equation}
A Equação~\eqref{eq:v1-c7-27} é a origem da relação que o Volume II usa como
dada ao decompor as perdas do motor: a fração da potência de entreferro perdida
no rotor é numericamente igual ao escorregamento. Ela também impõe um limite
duro que costuma ser esquecido em controle de velocidade: reduzir a rotação
aumentando o escorregamento --- por redução de tensão, por exemplo --- é
termodinamicamente desastroso, porque a $s = 0{,}5$ metade de toda a potência
de entreferro vira calor no rotor. Controlar velocidade por frequência, e não
por escorregamento, não é um refinamento de conforto: é a única forma de fazê-lo
sem destruir o rendimento --- e essa é a justificativa física do inversor.

\textbf{Torque e escorregamento.} Combinando as
Equações~\eqref{eq:v1-c7-26} e~\eqref{eq:v1-c7-27} com o equivalente de
Thévenin do estator chega-se à curva torque--escorregamento, cuja forma o
Volume II usa em regime permanente. Dois fatos qualitativos dela merecem
registro por terem consequência de diagnóstico. Primeiro: o torque máximo
(torque de \emph{breakdown}) \textbf{não} depende de $R_r$; o que depende de
$R_r$ é o escorregamento em que ele ocorre. Aumentar a resistência do rotor
--- por temperatura, pela Equação~\eqref{eq:v1-c7-11}, ou por projeto de barra
profunda --- desloca o pico para escorregamentos maiores sem alterar sua
altura. Segundo: para escorregamentos pequenos, a curva é praticamente linear,
e o torque é proporcional a $s$ --- é o que torna legítimo usar corrente ou
escorregamento como estimadores de carga em regime.

\textbf{Assinaturas de falha no espectro.} A estrutura periódica do rotor
imprime frequências características na corrente do estator, e é isso que
permite diagnosticar mecanicamente uma máquina a partir de uma medida elétrica.
Uma barra rompida quebra a simetria do circuito rotórico, produzindo bandas
laterais em torno da fundamental afastadas de $2sf$; uma excentricidade
modula o entreferro e produz bandas afastadas de múltiplos da frequência de
rotação. O Volume II usa essas assinaturas como indicadores; a razão de elas
existirem está na dependência de $\lambda(\theta)$ que a
Equação~\eqref{eq:v1-c7-21} já anunciou: qualquer defeito que altere a
geometria magnética altera essa dependência, e o faz na frequência com que a
geometria se repete.

\subsection{Acionamento por inversor de frequência}\label{sub:v1-c7-inversor}

\textbf{Por que $V/f$.} O fluxo concatenado por polo em uma máquina alimentada
por tensão é aproximadamente proporcional à razão entre tensão e frequência,
Equação~\eqref{eq:v1-c7-28}, resultado que sai diretamente da lei de Faraday
integrada sobre um enrolamento de $N$ espiras ($\hat{\Phi}$ é o fluxo por
polo e $V_f$ o valor eficaz da tensão de fase):
\begin{equation}\label{eq:v1-c7-28}
\hat{\Phi} \approx \frac{V_{f}}{4{,}44\,f\,N}
\quad\Longrightarrow\quad
\frac{V}{f} = \text{constante} \;\Leftrightarrow\; \hat{B} = \text{constante} .
\end{equation}
Reduzir a frequência sem reduzir a tensão aumentaria $\hat{B}$ e levaria o aço
à saturação da Seção~\ref{sub:v1-c7-constitutivas}, com corrente de
magnetização explodindo e perda no ferro crescendo pela
Equação~\eqref{eq:v1-c7-12}. Manter $V/f$ mantém o fluxo, e com ele a
capacidade de torque, em toda a faixa abaixo da rotação nominal --- por isso
essa região é de \emph{torque constante}. Acima da rotação nominal a tensão não
pode mais crescer (o barramento CC é o teto), $V/f$ cai, o fluxo enfraquece e a
máquina passa à região de \emph{potência constante}, com torque decrescente.
A relação $V/f$ que o Volume II enuncia como regra de operação do inversor é,
portanto, uma consequência direta da não linearidade do meio magnético.

\textbf{O que a comutação produz além da fundamental.} O inversor sintetiza a
tensão desejada por modulação de largura de pulso: a tensão de saída é uma
sequência de degraus entre os níveis do barramento CC, cujo valor médio em cada
período de comutação segue a senoide de referência. A indutância do motor
filtra a \emph{corrente}, que resulta razoavelmente senoidal; ela não filtra a
\emph{tensão}, que chega aos terminais como degraus com o
$\mathrm{d}v/\mathrm{d}t$ íngreme já discutido na
Seção~\ref{sub:v1-c7-concentrados}. Daí decorrem duas solicitações que não
existem em alimentação senoidal: a solicitação dielétrica do isolamento, agravada
pela sobretensão de reflexão, e a excitação das capacitâncias parasitas.

\textbf{A tensão de modo comum, e de onde ela vem.} Este é o ponto em que a
cadeia se fecha. Define-se a tensão de modo comum como a média instantânea das
três tensões de fase em relação a uma referência,
Equação~\eqref{eq:v1-c7-29}:
\begin{equation}\label{eq:v1-c7-29}
v_{cm}(t) = \frac{v_a(t) + v_b(t) + v_c(t)}{3} .
\end{equation}
Em uma fonte senoidal equilibrada, a Equação~\eqref{eq:v1-c7-29} é
identicamente nula --- é o que significa ser equilibrado. Em um inversor
trifásico de dois níveis, cada braço só pode estar conectado ao polo positivo
ou ao negativo do barramento, e não existe combinação de estados dos três
braços que anule a soma: $v_{cm}$ salta em degraus de $V_{cc}/3$ a cada
comutação. A tensão de modo comum não é um defeito de projeto do inversor; é
uma consequência topológica de sintetizar três tensões com seis chaves.

Sendo uma grandeza de sequência zero, ela não encontra caminho pelos
enrolamentos de um motor sem neutro, como a
Seção~\ref{sub:v1-c7-trifasico} observou. Procura, então, os caminhos
capacitivos parasitas do Modo~3 da Seção~\ref{sub:v1-c7-concentrados}, e o
divisor formado por essas capacitâncias impõe uma fração dela ao eixo do rotor.
Quando essa tensão de eixo excede a rigidez dielétrica do filme de óleo do
rolamento, o filme rompe e a energia armazenada se descarrega em um pulso que
erode a pista. É o mecanismo de correntes de rolamento induzidas por inversor,
documentado quantitativamente por Erdman \textit{et al.} \cite{erdman} e
modelado no Volume II como um divisor capacitivo com critério de disparo. Este
capítulo fornece as três peças de que aquele modelo depende e que ele não
deriva: a existência das capacitâncias (lei de Gauss,
Seção~\ref{sub:v1-c7-bases}), a razão de elas conduzirem em alta frequência
(falha da hipótese concentrada, Seção~\ref{sub:v1-c7-concentrados}) e a origem
da tensão que as excita (Equação~\eqref{eq:v1-c7-29}).

\textbf{O lado da rede.} O retificador de entrada também não é transparente. Um
retificador de seis pulsos drena corrente com harmônicos de ordem
$h = 6k\pm1$ ($5^{\underline{a}}$, $7^{\underline{a}}$, $11^{\underline{a}}$,
$13^{\underline{a}}$, \dots), com amplitudes que decaem aproximadamente como
$1/h$. São esses harmônicos que alimentam a $\mathrm{THD}_i$ da
Equação~\eqref{eq:v1-c7-17} e o fator de distorção da
Equação~\eqref{eq:v1-c7-18}. Em uma rede isolada offshore, de potência de
curto-circuito limitada, eles distorcem também a \emph{tensão} do barramento,
afetando todas as demais cargas --- razão pela qual a distorção é um requisito
de projeto de instalação, e não apenas uma característica do acionamento.

\subsection{Perdas, rendimento e o que é observável em
campo}\label{sub:v1-c7-perdas}

A contabilidade de perdas de uma máquina elétrica é normatizada, e o método de
perdas segregadas da IEC 60034-2-1 \cite{iec60034} separa cinco categorias, que
a Tabela~\ref{tab:v1-c7-3} reúne com sua origem física neste capítulo e seu
comportamento com a carga. A tabela é o insumo direto do modelo de rendimento
do Volume II, e sua coluna mais importante é a última: é ela que diz o que se
pode esperar que mude quando a máquina degrada.

\begin{table}[H]
\centering
\caption[Perdas segregadas de máquina de indução]{Categorias de perda da
IEC 60034-2-1, origem física neste capítulo e dependência da carga.}
\label{tab:v1-c7-3}
\small
\resizebox{\textwidth}{!}{%
\begin{tabular}{@{}llll@{}}
\toprule
\textbf{Categoria} & \textbf{Origem física} & \textbf{Equação} & \textbf{Dependência da carga} \\
\midrule
Joule no estator & $\mathbf{J}\cdot\mathbf{E}$ no cobre & \eqref{eq:v1-c7-4}, \eqref{eq:v1-c7-11} & $\propto I^2$; cresce com $T$ \\
Joule no rotor & Corrente induzida no rotor & \eqref{eq:v1-c7-27} & $= s\,P_{ag}$; cresce com a carga \\
Núcleo (ferro) & Histerese + correntes parasitas & \eqref{eq:v1-c7-12} & Quase constante ($\hat{B}$, $f$ fixos) \\
Atrito e ventilação & Mancais e ventilador próprio & Cap.~\ref{sec:mecanica} & Constante em rotação fixa \\
Suplementares & Harmônicos espaciais, dispersão & --- & Cresce com $I^2$; medida por ensaio \\
\bottomrule
\end{tabular}}
\end{table}

Somando as categorias, o rendimento fecha pela
Equação~\eqref{eq:v1-c7-30}:
\begin{equation}\label{eq:v1-c7-30}
\eta_{mot} = \frac{P_{mec} - P_{atr}}{P_{el,ent}}
= \frac{(1-s)P_{ag} - P_{atr}}{P_{ag} + P_{j,est} + P_{fe}} .
\end{equation}
onde $P_{atr}$ reúne atrito e ventilação e as perdas suplementares foram
omitidas por conveniência de escrita --- elas são de poucos décimos de ponto
percentual e entram, na prática de ensaio, agregadas às demais parcelas fixas.

Duas leituras da Equação~\eqref{eq:v1-c7-30} importam para o resto da coleção.
A primeira é a forma da curva de rendimento em função da carga: como as perdas
de núcleo, atrito e ventilação são praticamente constantes enquanto a potência
útil cai, o rendimento é aproximadamente plano entre $50\,\%$ e $100\,\%$ da
carga nominal e desaba abaixo de cerca de $40\,\%$. Um compressor que passa a
maior parte do tempo em carga parcial paga esse desabamento continuamente, e é
essa a razão termodinâmica --- não apenas de conveniência de controle --- de o
acionamento de velocidade variável melhorar o consumo específico de uma
planta de ar comprimido.

A segunda leitura é sobre \textbf{observabilidade}, e conecta este capítulo ao
problema central do Volume II. Das grandezas da
Equação~\eqref{eq:v1-c7-30}, apenas $P_{el,ent}$ e o escorregamento são
medidos rotineiramente em campo; $R_s$ é medível com a máquina parada; $R_r$ é
inferível; e as perdas de núcleo, atrito e ventilação só se obtêm por ensaio a
vazio e rotor bloqueado, em bancada. Elas entram no modelo como \emph{valores
de catálogo} --- exatamente a classe de parâmetro que o Volume II identifica
como fonte de erro quando usada cegamente, e que o Gêmeo Digital recalibra a
partir de dados de operação. Este capítulo torna explícito \emph{por que} esses
parâmetros são de catálogo: eles dependem de propriedades constitutivas
medidas ($k_h$, $k_e$, $\sigma_e(T)$, a curva $B(H)$ do aço específico) que
nenhuma lei de balanço fornece --- a mesma conclusão a que a
Seção~\ref{sub:v1-c2-constitutivas} chegou para viscosidade e condutividade
térmica, e pela mesma razão.

\subsection{Síntese: premissas e ligação com o Volume
II}\label{sub:v1-c7-sintese}

A Tabela~\ref{tab:v1-c7-4} lista os resultados centrais do capítulo com as
premissas de que dependem, no mesmo formato da Tabela~\ref{tab:v1-c3-2}, e
serve como lista de verificação antes de aplicar qualquer um deles fora do
contexto em que foi deduzido.

\begin{table}[H]
\centering
\caption[Premissas dos resultados do capítulo]{Resultados centrais do
Capítulo~\ref{sec:eletrica} e as premissas de que dependem.}
\label{tab:v1-c7-4}
\small
\resizebox{\textwidth}{!}{%
\begin{tabular}{@{}lll@{}}
\toprule
\textbf{Resultado} & \textbf{Equação} & \textbf{Premissas} \\
\midrule
Conservação de carga & \eqref{eq:v1-c7-1} & Nenhuma; contida em \eqref{eq:v1-c7-2} \\
Leis de Kirchhoff & \eqref{eq:v1-c7-6}, \eqref{eq:v1-c7-7} & Regime quase-estático, $\ell/\lambda_{em}\ll1$ \\
Lei de Ohm integral & \eqref{eq:v1-c7-10} & Seção uniforme; $r_{\text{cond}}\ll\delta_p$; $\sigma_e$ isotrópico \\
Impedância e fasores & \eqref{eq:v1-c7-13} & Linearidade; frequência única; regime permanente \\
Triângulo de potências & \eqref{eq:v1-c7-15} & Tensão e corrente senoidais \\
Potência trifásica constante & \eqref{eq:v1-c7-16} & Sistema equilibrado nas três fases \\
Fator de potência verdadeiro & \eqref{eq:v1-c7-18} & Tensão senoidal; corrente distorcida \\
Componentes simétricas & \eqref{eq:v1-c7-19} & Sistema linear; decomposição sempre válida \\
Torque por coenergia & \eqref{eq:v1-c7-21} & Perdas separáveis do campo de acoplamento \\
Torque por tensão de Maxwell & \eqref{eq:v1-c7-23} & $\bar{\sigma}_t$ limitada pela saturação do aço \\
Velocidade síncrona & \eqref{eq:v1-c7-24} & Enrolamento distribuído senoidalmente; $p$ polos \\
Repartição $s\,P_{ag}$ & \eqref{eq:v1-c7-27} & Circuito equivalente por fase; regime permanente \\
Regra $V/f$ & \eqref{eq:v1-c7-28} & Queda em $R_s$ desprezível; abaixo da rotação nominal \\
Tensão de modo comum & \eqref{eq:v1-c7-29} & Inversor de dois níveis; três braços \\
Rendimento & \eqref{eq:v1-c7-30} & Perdas segregadas conforme IEC 60034-2-1 \\
\bottomrule
\end{tabular}}
\end{table}

Por fim, a Tabela~\ref{tab:v1-c7-5} explicita a função deste capítulo dentro da
coleção: cada resultado do subsistema elétrico que o Volume II utiliza como
ponto de partida tem, aqui, a sua origem. É a leitura recomendada para quem
chegar a este volume vindo de lá.

\begin{table}[H]
\centering
\caption[Ligação com o Volume II]{Correspondência entre os fundamentos deste
capítulo e os modelos do subsistema elétrico do Volume II.}
\label{tab:v1-c7-5}
\small
\resizebox{\textwidth}{!}{%
\begin{tabular}{@{}lll@{}}
\toprule
\textbf{O que o Volume II usa} & \textbf{Fundamento aqui} & \textbf{Seção} \\
\midrule
Regra $V/f$ do inversor & Saturação magnética e fluxo por polo & \ref{sub:v1-c7-constitutivas}, \ref{sub:v1-c7-inversor} \\
Torque eletromagnético (Thévenin) & Coenergia e tensão de Maxwell & \ref{sub:v1-c7-conversao} \\
Escorregamento como medida de carga & Necessidade física do escorregamento & \ref{sub:v1-c7-inducao} \\
Perda de rotor igual a $s\,P_{ag}$ & Repartição da potência de entreferro & \ref{sub:v1-c7-inducao} \\
$\eta_{mot}(s,T)$ e $P_{\text{fixas}}$ & Perdas segregadas; $\sigma_e(T)$; Steinmetz & \ref{sub:v1-c7-constitutivas}, \ref{sub:v1-c7-perdas} \\
Índices de desbalanço NEMA & Impedância de sequência negativa & \ref{sub:v1-c7-trifasico} \\
Bandas laterais de barra quebrada & Dependência de $\lambda(\theta)$ & \ref{sub:v1-c7-inducao} \\
\textit{Derating} por $\mathrm{THD}_i$ & Fator de distorção; perda extra & \ref{sub:v1-c7-potencia} \\
Tensão de eixo e EDM & Modo comum e capacitâncias parasitas & \ref{sub:v1-c7-concentrados}, \ref{sub:v1-c7-inversor} \\
Degradação térmica do isolamento & Reflexão de onda e $\mathrm{d}v/\mathrm{d}t$ & \ref{sub:v1-c7-concentrados} \\
\bottomrule
\end{tabular}}
\end{table}


%=============================================================================
% CAPÍTULO 8 — FUNDAMENTOS DE MECÂNICA
%=============================================================================
\section{Fundamentos de Mecânica}\label{sec:mecanica}

\subsection{Dinâmica rotacional}

A dinâmica de corpos rígidos em rotação é análoga à dinâmica translacional. A
segunda lei de Newton para rotação é
\begin{equation}\label{eq:v1-c8-1}
\tau = J_p \, \alpha,
\end{equation}
onde $\tau$ é o torque, $J_p$ é o momento de inércia em relação ao eixo de
rotação e $\alpha$ é a aceleração angular. Os símbolos seguem os do
Volume II: o torque é $\tau$ --- e não $T$, que é a temperatura nos
Capítulos~\ref{sec:termodinamica} a~\ref{sec:calor} --- e o momento de inércia
é $J_p$ --- e não $I$, que é a corrente elétrica no
Capítulo~\ref{sec:eletrica}, nem $J$, que é o jacobiano da deformação no
Capítulo~\ref{sec:continuo}.

A energia cinética de rotação é
\begin{equation}\label{eq:v1-c8-2}
E_c = \frac{1}{2} J_p \,\omega^2.
\end{equation}

\subsection{Vibrações mecânicas}

Um sistema massa-mola-amortecedor unidimensional é governado por
\begin{equation}\label{eq:v1-c8-3}
m\,\ddot{x} + c\,\dot{x} + k\,x = F(t),
\end{equation}
onde $m$ é a massa, $c$ é o coeficiente de amortecimento, $k$ é a rigidez e
$F(t)$ é a força externa.

A frequência natural não amortecida é
\begin{equation}\label{eq:v1-c8-4}
\omega_n = \sqrt{\frac{k}{m}}.
\end{equation}

A frequência natural amortecida é
\begin{equation}\label{eq:v1-c8-5}
\omega_d = \omega_n\sqrt{1-\zeta^2},
\end{equation}
onde $\zeta = c/(2\sqrt{km})$ é a razão de amortecimento.

\subsection{Desbalanceamento e forças rotativas}

Uma massa excêntrica $m_u$ girando a uma velocidade angular $\omega$ produz
uma força de desbalanceamento
\begin{equation}\label{eq:v1-c8-6}
F_u = m_u \, e \, \omega^2,
\end{equation}
onde $e$ é a excentricidade. Essa força é transmitida aos mancais e pode
provocar vibração excessiva se a velocidade de operação estiver próxima a uma
frequência crítica do eixo.

\subsection{Fadiga}

A fadiga é a falha por aplicação repetida de tensões abaixo do limite de
escoação. Para carregamento completamente reverso, a lei de Basquin descreve
a vida à fadiga $N_f$ em função da amplitude de tensão $\sigma_a$:
\begin{equation}\label{eq:v1-c8-7}
\sigma_a^m \, N_f = C,
\end{equation}
onde $m$ e $C$ são constantes do material.


%=============================================================================
% CAPÍTULO 9 — INSTRUMENTAÇÃO E MEDIÇÃO
%=============================================================================
\section{Instrumentação e Medição}\label{sec:instrumentacao}

\subsection{O papel da instrumentação}

A instrumentação é o elo entre o modelo físico e o mundo real. Sem medições
confiáveis, não é possível calibrar modelos, validar simulações nem implementar
controle preditivo. Este capítulo apresenta os princípios de medição, os
principais tipos de sensores empregados em máquinas rotodinâmicas e as
noções básicas de incerteza.

\subsection{Princípios de medição}

Uma medição é uma estimativa de uma grandeza física com uma determinada
incerteza. O resultado de uma medição é completo apenas quando acompanhado de
sua incerteza:
\begin{equation}\label{eq:v1-c9-1}
y = y_{med} \pm U,
\end{equation}
onde $U$ é a incerteza expandida, geralmente calculada para um fator de
cobertura $k=2$ (aproximadamente 95\% de confiança).

A incerteza combinada $u_c$ é obtida a partir das incertezas das fontes
individuais $u_i$:
\begin{equation}\label{eq:v1-c9-2}
u_c = \sqrt{\sum_i \left(\frac{\partial f}{\partial x_i}\right)^2 u_i^2},
\end{equation}
para uma grandeza $y = f(x_1, x_2, \ldots, x_n)$.

\subsection{Sensores para máquinas rotodinâmicas}

A Tabela~\ref{tab:v1-c9-1} resume os sensores mais comuns no monitoramento de
compressores, bombas, turbinas e ventiladores.

\begin{table}[H]
\centering
\caption{Sensores típicos em máquinas rotodinâmicas.}
\label{tab:v1-c9-1}
\small
\begin{tabular}{@{}p{2.0cm}p{3.0cm}p{4.0cm}p{4.2cm}@{}}
\toprule
\textbf{Grandeza} & \textbf{Tipo de sensor} & \textbf{Princípio} & \textbf{Observações} \\
\midrule
Pressão & Piezoelétrico, piezorresistivo & Deformação de diafragma & Dinâmica ou estática \\
Temperatura & Termopar, PT100 & Efeito termoelétrico/resistivo & Termopares rápidos \\
Vazão & Orifício, turbina, ultrasônico & Diferença de pressão, rotação ou tempo de trânsito & Cada tipo tem faixa própria \\
Vibração & Acelerômetro, velocímetro & Piezoelétrico ou eletromagnético & Análise espectral \\
Velocidade & Encoder, tacogerador & Pulso eletromagnético ou óptico & Precisão de posição \\
Corrente & Transformador de corrente & Indução eletromagnética & Não invasivo \\
\bottomrule
\end{tabular}
\end{table}

\subsection{Aquisição de dados}

A aquisição de dados (DAQ) converte sinais analógicos em digitais para
processamento. A taxa de amostragem $f_s$ deve respeitar o teorema de Nyquist:
\begin{equation}\label{eq:v1-c9-3}
f_s \geq 2\,f_{max},
\end{equation}
onde $f_{max}$ é a maior frequência de interesse no sinal.

\subsection{Diagnóstico e monitoramento de condição}

O monitoramento de condição baseia-se na comparação entre valores medidos e
valores esperados do modelo. Um resíduo é definido como
\begin{equation}\label{eq:v1-c9-4}
r = y_{medido} - y_{modelo}.
\end{equation}
Quando o resíduo ultrapassa um limiar estatístico, indica-se uma possível
anomalia. Essa ideia é a base dos algoritmos de Gêmeo Digital apresentados no
Volume II.


%=============================================================================
% CAPÍTULO 10 — SÍNTESE E LIGAÇÃO COM MÁQUINAS ROTODINÂMICAS
%=============================================================================
\section{Síntese e Ligação com Máquinas Rotodinâmicas}\label{sec:sintese}

\subsection{Unindo os fundamentos}

Uma máquina rotodinâmica real exige a combinação de todas as áreas apresentadas
neste volume. Tomemos como exemplo um compressor de parafuso acionado por motor
elétrico:

\begin{itemize}
\item A \textbf{mecânica do contínuo} fornece as equações de conservação de
  massa, quantidade de movimento e energia no interior dos rotores.
\item A \textbf{termodinâmica} define o trabalho de compressão, a elevação de
  temperatura e a eficiência politrópica.
\item A \textbf{mecânica dos fluidos} quantifica as perdas por atrito nos
  ductos de admissão e descarga.
\item A \textbf{transferência de calor} dimensiona os trocadores que
  resfriam o ar entre estágios.
\item A \textbf{eletricidade} dimensiona o motor e o inversor de frequência.
\item A \textbf{mecânica} garante que o eixo opere abaixo das frequências
  críticas e suporte as forças de desbalanceamento.
\item A \textbf{instrumentação} fornece os dados para calibrar e validar o
  Gêmeo Digital.
\end{itemize}

\subsection{Do modelo físico ao Gêmeo Digital}

O Gêmeo Digital não substitui o modelo físico; ele o atualiza continuamente
com dados de campo. O fluxo de informação é:
\begin{enumerate}
\item O modelo físico prediz o comportamento esperado da máquina.
\item Os sensores medem o comportamento real.
\item Os resíduos entre predição e medição são analisados.
\item Algoritmos de estimação (EKF), detecção de anomalia (VAE) e otimização
  (MPC/\-PSO) ajustam o modelo e prescrevem ações.
\end{enumerate}

Esse fluxo só é possível porque cada fenômeno físico foi previamente modelado
a partir dos fundamentos corretos — o que justifica a existência deste volume.

\subsection{Próximos passos}

Com os fundamentos estabelecidos, o leitor está apto a acompanhar o Volume II,
onde essas mesmas teorias são aplicadas a um sistema de compressão de ar
offshore completo. As equações, convenções e notação deste volume serão
reutilizadas lá, de modo que a transição seja direta.


%=============================================================================
% BIBLIOGRAFIA
%=============================================================================
\newpage
\begin{thebibliography}{99}

\bibitem{malvern1969} Malvern, L. E. \textit{Introduction to the Mechanics of a Continuous Medium}. Prentice-Hall, 1969.

\bibitem{gurtin1981} Gurtin, M. E. \textit{An Introduction to Continuum Mechanics}. Academic Press, 1981.

\bibitem{sonntag2003} Sonntag, R. E., Borgnakke, C., Van Wylen, G. J. \textit{Fundamentos da Termodinâmica}. 7ª ed. Edgard Blücher, 2003.

\bibitem{moran2018} Moran, M. J., Shapiro, H. N. \textit{Fundamentos de Engenharia Termodinâmica}. 8ª ed. LTC, 2018.

\bibitem{white2011} White, F. M. \textit{Mecânica dos Fluidos}. 6ª ed. McGraw-Hill, 2011.

\bibitem{fox2020} Fox, R. W., McDonald, A. T., Pritchard, P. J. \textit{Introduction to Fluid Mechanics}. 10ª ed. Wiley, 2020.

\bibitem{incropera2008} Incropera, F. P., DeWitt, D. P. \textit{Fundamentos de Transferência de Calor e de Massa}. 6ª ed. LTC, 2008.

\bibitem{bejan1993} Bejan, A. \textit{Heat Transfer}. Wiley, 1993.

\bibitem{chapman2012} Chapman, S. J. \textit{Fundamentos de Máquinas Elétricas}. Cengage Learning, 2012.

\bibitem{fitzgerald2003} Fitzgerald, A. E., Kingsley, C., Umans, S. D. \textit{Máquinas Elétricas}. 6ª ed. Bookman, 2003.

\bibitem{erdman} Erdman, J.M.; Kerkman, R.J.; Schlegel, D.W.; Skibinski, G.L. \textit{Effect of PWM Inverters on AC Motor Bearing Currents and Shaft Voltages}. IEEE Transactions on Industry Applications, 32(2), 250--259, 1996. DOI: 10.1109/28.491472

\bibitem{iec60034} International Electrotechnical Commission. \textit{IEC 60034-2-1:2024 — Rotating electrical machines — Part 2-1: Standard methods for determining losses and efficiency from tests (excluding machines for traction vehicles)}, Ed. 3.0. Geneva, 2024.

\bibitem{nemamg1} National Electrical Manufacturers Association. \textit{ANSI/NEMA MG 1-2021 — Motors and Generators}. Rosslyn, VA, 2021. A Parte 14 (Application Data) define o desbalanço de tensão como o maior desvio absoluto em relação à média das três fases, dividido pela média, e recomenda o limite de 1\% nos terminais do motor em carga nominal, com curva de \textit{derating} acima desse valor.

\bibitem{rao2011} Rao, S. S. \textit{Mechanical Vibrations}. 5ª ed. Prentice-Hall, 2011.

\bibitem{shigley2011} Shigley, J. E., Mischke, C. R., Budynas, R. G. \textit{Projeto de Elementos de Máquinas}. 8ª ed. Bookman, 2011.

\bibitem{figliola2011} Figliola, R. S., Beasley, D. E. \textit{Theory and Design for Mechanical Measurements}. 5ª ed. Wiley, 2011.

\bibitem{international2008} Joint Committee for Guides in Metrology (JCGM). \textit{Evaluation of Measurement Data — Guide to the Expression of Uncertainty in Measurement (GUM 1995 with minor corrections)}. 2008.

\end{thebibliography}

\end{document}
